\subsection{Which hypotheses suffice}
\label{ssec:closure}

The implications proved above, together with those of \Cref{ssec:bullseye},
determine exactly which combinations of statements yield the conjecture. The
logical skeleton of the paper and of the literature it quotes is finite: a
set of statements, each of them proved here, quoted from the literature, or
taken as a hypothesis, and a set of inference rules, each of them one
theorem. Consequences can then be taken.

The answer is only as strong as the rule set. A set of statements the rules
do not carry to the conjecture is \emph{insufficient} in the sense of the
rules, and that means no more than that no rule derives the conjecture from
it. An implication that is true and left out makes a set look necessary when
it is not. One such implication
is elementary, and it decides the shape of the answer, so we prove it first.

\begin{proposition}[The conjecture for the varieties that are not abelian
is the conjecture]\label{prop:f3ishc}
Let $A$ be a complex abelian variety of dimension $g\ge1$ and let
$X=A\times\PP^{1}$. Then $X$ is a smooth complex projective variety that is
not an abelian variety, and if every Hodge class of degree $2p$ on $X$ is
algebraic, then so is every Hodge class of degree $2p$ on $A$. Consequently
the statement
\begin{quote}
the Hodge conjecture holds for every smooth complex projective variety that
is not an abelian variety
\end{quote}
implies the Hodge conjecture for every abelian variety, and it is equivalent
to $\mathbf{HC}$.
\end{proposition}

\begin{proof}
$X$ is smooth and projective of dimension $g+1$. By the K\"unneth formula
$h^{1,0}(X)=h^{1,0}(A)+h^{1,0}(\PP^{1})=g$, while an abelian variety of
dimension $g+1$ has $h^{1,0}=g+1$; so $X$ is not an abelian variety. (Equally,
$X$ contains the rational curves $\{a\}\times\PP^{1}$, and every morphism
from $\PP^{1}$ to a complex torus is constant, since it lifts to the
universal cover $\CC^{g+1}$ \cite[Chapter~1]{BL04}.)

Let $\mathrm{pr}\colon X\to A$ be the projection and $D=A\times\{0\}$, a
smooth divisor that $\mathrm{pr}$ maps isomorphically onto $A$, so that
$\mathrm{pr}_{*}[D]=1$ in $H^{0}(A,\QQ)$. Let $\alpha$ be a Hodge class of
degree $2p$ on $A$. Pull-back along a morphism is a morphism of Hodge
structures, so $\mathrm{pr}^{*}\alpha$ is a Hodge class of degree $2p$ on
$X$, and by hypothesis $\mathrm{pr}^{*}\alpha=[Z]$ for an algebraic cycle
$Z$ on $X$ with rational coefficients. The projection formula gives
\[
  \mathrm{pr}_{*}\bigl([Z]\cup[D]\bigr)
  =\mathrm{pr}_{*}\bigl(\mathrm{pr}^{*}\alpha\cup[D]\bigr)
  =\alpha\cup\mathrm{pr}_{*}[D]=\alpha .
\]
The cycle class map is compatible with intersection products and with proper
push-forward \cite[Chapter~19]{Ful98}, so the left side
is the class of the cycle $\mathrm{pr}_{*}(Z\cdot D)$ and $\alpha$ is
algebraic. A point, the abelian variety of dimension zero, satisfies the
conjecture trivially.

So the displayed statement gives the conjecture for every abelian variety,
and with it the conjecture for every smooth complex projective variety. The
converse is immediate: the displayed statement is a special case of
$\mathbf{HC}$.
\end{proof}

Since the conjecture for the varieties that are not abelian is the whole
conjecture, the statement that belongs on a frontier is weaker: it asks for
algebraicity only modulo what abelian varieties supply.

\begin{definition}[The conjecture modulo abelian varieties]
\label{def:f3prime}
For a smooth complex projective variety $X$ and $p\ge0$ let
$\mathrm{Ab}^{p}(X)\subseteq H^{2p}(X,\QQ)$ be the span of the classes
$\Gamma_{*}\gamma=\mathrm{pr}_{X*}\bigl([\Gamma]\cup\mathrm{pr}_{B}^{*}\gamma\bigr)$
of degree $2p$, where $B$ runs over the complex abelian varieties, the point
included, $\gamma$ over the Hodge classes on $B$, and $\Gamma$ over the
algebraic cycles on $B\times X$ with rational coefficients. The statement
\textup{(F3$'$)} is: for every $X$ and every $p$,
\[
  \Hdg^{p}(X)=\mathrm{Ab}^{p}(X).
\]
\end{definition}

Each $\Gamma_{*}\gamma$ is a Hodge class, since cup product with an
algebraic class and push-forward preserve Hodge classes, so
$\mathrm{Ab}^{p}(X)\subseteq\Hdg^{p}(X)$; and taking $B$ a point shows that
$\mathrm{Ab}^{p}(X)$ contains every algebraic class. So \textup{(F3$'$)} says
that every Hodge class is algebraic modulo the images of Hodge classes of
abelian varieties under algebraic correspondences, which is the form of (F3)
that \Cref{rem:f3collapse} asks for.

\begin{proposition}[What is known of \textup{(F3$'$)}]\label{prop:f3prime}
\leavevmode
\begin{enumerate}[label=\textup{(\roman*)},leftmargin=2.6em]
\item $\mathbf{HC}$ implies \textup{(F3$'$)}, and \textup{(F3$'$)} together
  with the conjecture for abelian varieties implies $\mathbf{HC}$.
\item \textup{(F3$'$)} holds on every abelian variety and on every product
  $A\times\PP^{1}$ with $A$ abelian; the argument of \Cref{prop:f3ishc} gives
  nothing for it.
\item Let $\mathcal{A}$ be the class of smooth complex projective varieties
  $X$ for which $H^{*}(X,\QQ)$ is spanned by classes $\Gamma_{*}y$ with $B$
  abelian, $\Gamma$ an algebraic cycle on $B\times X$ and
  $y\in H^{*}(B,\QQ)$. Every $X\in\mathcal{A}$ satisfies \textup{(F3$'$)}.
  The class $\mathcal{A}$ contains every abelian variety and every smooth
  projective curve, and it is closed under products and under images of
  surjective morphisms. So it contains every product of curves, every
  symmetric product of a curve, and every smooth projective variety onto
  which a product of curves or an abelian variety maps surjectively.
\item Every $X$ of dimension $n$ satisfies \textup{(F3$'$)} in the degrees
  $2p$ with $p\in\{0,1,n-1,n\}$, so every $X$ of dimension at most three
  satisfies it in every degree.
\end{enumerate}
\end{proposition}

\begin{proof}
(i) If $\mathbf{HC}$ holds, then $\Hdg^{p}(X)$ consists of algebraic classes,
which lie in $\mathrm{Ab}^{p}(X)$. Conversely, if the conjecture holds for
abelian varieties, every $\gamma$ in \Cref{def:f3prime} is algebraic, so
every $\Gamma_{*}\gamma$ is, since composition with a correspondence carries
algebraic classes to algebraic classes \cite[Chapter~16]{Ful98}; then
\textup{(F3$'$)} says that every Hodge class is algebraic.

(ii) On an abelian variety $A$ take $B=A$ and $\Gamma$ the diagonal. On
$X=A\times\PP^{1}$ the K\"unneth formula gives
$H^{*}(X,\QQ)=\mathrm{pr}^{*}H^{*}(A,\QQ)\oplus
\mathrm{pr}^{*}H^{*}(A,\QQ)\cup\eta$, with $\eta$ the class of a point of
$\PP^{1}$ pulled back to $X$, which is algebraic. The first summand is
$\Gamma_{*}H^{*}(A,\QQ)$ for $\Gamma$ the transpose of the graph of
$\mathrm{pr}$, and the second is $\Gamma'_{*}H^{*}(A,\QQ)$ for
$\Gamma'=\Gamma\cdot\mathrm{pr}_{X}^{*}\eta$. So $X\in\mathcal{A}$, and
(iii) applies.

(iii) Let $X\in\mathcal{A}$ and fix $p$. Since $H^{2p}(X,\QQ)$ is finite
dimensional, finitely many classes $\Gamma_{j*}y_{j}$ span it, and we may
take each $\Gamma_{j}$ of pure codimension $c_{j}$ on $B_{j}\times X$ and each
$y_{j}$ of pure degree $k_{j}$, with $k_{j}+2c_{j}-2\dim B_{j}=2p$. Then
$\Gamma_{j*}\colon H^{k_{j}}(B_{j},\QQ)\to H^{2p}(X,\QQ)$ is a morphism of
Hodge structures up to a Tate twist, and the sum of these maps is a
surjective morphism of polarisable Hodge structures onto $H^{2p}(X,\QQ)$.
Polarisable Hodge structures are semisimple, so its kernel has a complement
that is a sub-Hodge structure and maps isomorphically onto $H^{2p}(X,\QQ)$;
hence every Hodge class of degree $2p$ is the image of a Hodge class, that is
$\sum_{j}\Gamma_{j*}\gamma_{j}$ with each $\gamma_{j}$ a Hodge class on
$B_{j}$. So $\Hdg^{p}(X)=\mathrm{Ab}^{p}(X)$.

An abelian variety lies in $\mathcal{A}$ through its diagonal. A curve $C$
of genus $g\ge1$ lies in $\mathcal{A}$ through an Abel--Jacobi embedding
$\iota\colon C\to J(C)$: the transpose of its graph induces $\iota^{*}$,
which is an isomorphism on $H^{0}$ and $H^{1}$ and sends the class of the
theta divisor to $g$ times the class of a point, so it is surjective; the
projective line is the image of an elliptic curve under a double cover, and
is covered by the last closure property. For products, if $H^{*}(X_{i},\QQ)$
is spanned by the $\Gamma_{i*}y$, then by the K\"unneth formula
$H^{*}(X_{1}\times X_{2},\QQ)$ is spanned by the
$(\Gamma_{1}\times\Gamma_{2})_{*}(y_{1}\otimes y_{2})$, with
$\Gamma_{1}\times\Gamma_{2}$ an algebraic cycle on
$(B_{1}\times B_{2})\times(X_{1}\times X_{2})$. For images, let
$f\colon Y\to X$ be a surjective morphism with $Y\in\mathcal{A}$, let
$r=\dim Y-\dim X$ and let $h$ be an ample class on $Y$. The correspondence
$\Psi=[\Gamma_{f}]\cdot\mathrm{pr}_{Y}^{*}h^{r}$ on $Y\times X$ induces
$\Psi_{*}y=f_{*}(y\cup h^{r})$, and
$\Psi_{*}f^{*}x=x\cup f_{*}h^{r}=c\,x$ with $c>0$ the degree of $h^{r}$ on a
general fibre. So $\Psi_{*}$ is surjective, and $H^{*}(X,\QQ)$ is spanned by
the classes $(\Psi\circ\Gamma)_{*}y$, with $\Psi\circ\Gamma$ algebraic
\cite[Chapter~16]{Ful98}. A symmetric product of a curve is the image of a
power of it.

(iv) Every class of degree $0$ or $2n$ is algebraic, and every Hodge class of
degree $2$ is algebraic by the Lefschetz theorem on $(1,1)$ classes. For
$n\ge2$, the hard Lefschetz theorem makes cup product with $h^{n-2}$, for $h$
an ample class, an isomorphism $H^{2}(X,\QQ)\to H^{2n-2}(X,\QQ)$ of Hodge
structures up to a Tate twist, so every Hodge class of degree $2n-2$ is
$h^{n-2}$ times a Hodge class of degree $2$, and is algebraic. Algebraic
classes lie in $\mathrm{Ab}^{p}(X)$. When $n\le3$ every $p$ between $0$ and
$n$ is one of $0,1,n-1,n$.
\end{proof}

\begin{remark}[The reach of algebraic correspondences]\label{rem:f3primeopen}
\textup{(F3$'$)} is the statement that the Hodge conjecture reduces to
abelian varieties along algebraic correspondences, and
\Cref{prop:f3prime} records the cases of such a reduction that are theorems:
the varieties of $\mathcal{A}$, whose whole cohomology is reached from abelian
varieties by algebraic correspondences, and the degrees in which the
conjecture itself is a theorem. The next degree is degree four on fourfolds
outside $\mathcal{A}$ and outside the scope of \Cref{prop:f3primesmall} and
\Cref{prop:f3primefibration}. The square $S\times S$ of a K3 surface is the standard
example when $\End_{\mathrm{Hdg}}(T(S))$ is a totally real field of degree
at least two, that is, when $S$ has real multiplication. When that field is
$\QQ$ or a CM field, the Hodge conjecture holds on $S\times S$: by
\Cref{prop:orthpowers}(i) in the first case, and in the second by
\cite{Bus19} and the fact that the elements of norm one span a CM field, as in
\cite{RM08} (\Cref{prop:k3powers}). Its Hodge classes in
$H^{2}(S)\otimes H^{2}(S)$ include the classes
that induce the Hodge endomorphisms of the transcendental lattice $T(S)$, and
$S\times S$ lies in $\mathcal{A}$ as soon as the Kuga--Satake embedding of
$T(S)$ is induced by an algebraic cycle, the Kuga--Satake Hodge conjecture for
$S$: the transpose of such a cycle maps the
cohomology of the Kuga--Satake variety onto $T(S)$, and the rest of
$H^{*}(S,\QQ)$ is algebraic. The rational Hodge isometries are algebraic
\cite{Bus19}, but a real multiplication is not an isometry, and
\Cref{rem:mumfordks} meets exactly this obstacle for the surfaces
$S_{\lambda}$. With the Hodge conjecture for products of K3 surfaces of
\cite{OAI26f} as hypothesis the obstacle disappears for K3 surfaces
(\Cref{cor:k3all}), and the step itself remains.

For varieties whose cohomology is not of abelian type the position is
starker. By \Cref{prop:noabeliantype} no correspondence moves the $H^{3}$ of a
very general quintic threefold into the cohomology of an abelian variety, so
such a variety does not lie in $\mathcal{A}$. That is no obstruction to
\textup{(F3$'$)} itself, since a Hodge class spans a Hodge structure of Tate
type, which is of abelian type, and nothing in Hodge theory forbids it from
lying in $\mathrm{Ab}^{p}(X)$. Nor is it a route: to put a Hodge class into
$\mathrm{Ab}^{p}(X)$ one must construct an algebraic cycle on $B\times X$
whose class carries it, which is the kind of construction the Hodge
conjecture asks for, now on $B\times X$ instead of on $X$. The cases of
\textup{(F3$'$)} proved below outside $\mathcal{A}$ come from the Hodge
conjecture itself, from hypotheses on zero-cycles, and, in degree four, from
correspondences with members of $\mathcal{A}$
(\Cref{prop:k3powers} to \Cref{prop:f3primedominant}); none of them reaches a
variety in general position, and \textup{(F3$'$)} enters \Cref{thm:closure} as
a hypothesis.
\end{remark}

Two further facts fix the position of \textup{(F3$'$)}. The first says how
strong it is.

\begin{proposition}[The strength of \textup{(F3$'$)}]\label{prop:f3primestrength}
\textup{(F3$'$)} implies \textup{(M)}: every Hodge class on every smooth
complex projective variety is motivated in the sense of Andr\'e. In
particular it implies that every such class is absolutely Hodge. Adding the
implication $\textup{(F3$'$)}\Rightarrow\textup{(M)}$ to the rule set of
\Cref{setup:rules} leaves the five minimal sufficient sets of
\Cref{thm:closure}(ii) unchanged.
\end{proposition}

\begin{proof}
Let $\gamma\in\Hdg^{p}(X)$. By \textup{(F3$'$)},
$\gamma=\sum_{j}\Gamma_{j*}\beta_{j}$ with $\beta_{j}$ a Hodge class on an
abelian variety $B_{j}$ and $\Gamma_{j}$ an algebraic cycle on $B_{j}\times X$.
Every Hodge class on an abelian variety is motivated
\cite[Th\'eor\`eme~0.6.2]{And96b}, and the motivated classes form a graded
$\QQ$-algebra that contains the algebraic classes and is stable under
pull-back and push-forward along morphisms \cite[\S2]{And96b}. So
$\mathrm{pr}_{B_{j}}^{*}\beta_{j}\cup[\Gamma_{j}]$ is motivated, and so is its
push-forward $\Gamma_{j*}\beta_{j}$; hence $\gamma$ is motivated. Motivated
classes are absolutely Hodge \cite[\S2]{And96b}. For the last assertion,
the new rule has the premise (F3$'$) and the conclusion (M). The sets $C_{1}$
and $C_{2}$ of the proof of \Cref{thm:closure} omit (F3$'$), and $C_{3}$ and
$C_{4}$ contain both statements, so the four sets stay closed and that proof
applies without change; the five sets remain sufficient, since adding a rule
only enlarges closures.
\end{proof}

So on the route $\{\mathrm{L},\mathrm{F3}'\}$ the second member already
contains \textup{(M)}, and that route is never easier than
$\{\mathrm{L},\mathrm{M}\}$. The second fact locates, on the powers of a variety of orthogonal type,
the first class that \textup{(F3$'$)} has to reach.

\begin{proposition}[Powers of a variety of orthogonal type]
\label{prop:orthpowers}
Let $Y$ be a smooth complex projective variety of even dimension $m$. Suppose
that $H^{k}(Y,\QQ)$ is spanned by algebraic classes for $k\ne m$, that
$H^{m}(Y,\QQ)=N\oplus V$ with $N$ spanned by algebraic classes, the
intersection form $\psi$ nondegenerate on $N$ and $V=N^{\perp}$, and that the
Hodge group of $V$ is $\SO(V,\psi)$, with $t=\dim V\ge3$. Let
$\det(V)\in H^{mt}(Y^{t},\QQ)$ be the image of a generator of
$\bigwedge^{t}V\subseteq V^{\otimes t}$ under the K\"unneth embedding
$V^{\otimes t}\subseteq H^{mt}(Y^{t},\QQ)$.
\begin{enumerate}[label=\textup{(\roman*)},leftmargin=2.6em,itemsep=2pt]
\item The Hodge conjecture holds for $Y^{n}$ for every $n<t$.
\item $\det(V)$ is a Hodge class, and $\Hdg^{*}(Y^{t})$ is spanned by
  algebraic classes and $\det(V)$.
\item $\det(V)$ does not lie in the $\QQ$-algebra generated by the pull-backs
  along the projections of the Hodge classes of the $Y^{j}$, $j<t$.
\item The Hodge conjecture holds for every power of $Y$ if and only if
  $\det(V)$ is algebraic, and \textup{(F3$'$)} holds for every power of $Y$
  if and only if $\det(V)\in\mathrm{Ab}(Y^{t})$.
\end{enumerate}
\end{proposition}

\begin{proof}
Algebraic classes have even degree, so $H^{k}(Y,\QQ)=0$ for odd $k$. Let
$[\Delta]=\sum_{k}\pi_{k}$ be the K\"unneth decomposition of the class of the
diagonal, $\pi_{k}\in H^{2m-k}(Y,\QQ)\otimes H^{k}(Y,\QQ)$. For $k\ne m$,
$\pi_{k}=\sum_{i}e_{i}^{\vee}\otimes e_{i}$ for a basis $(e_{i})$ of
$H^{k}(Y,\QQ)$ and the Poincar\'e dual basis of $H^{2m-k}(Y,\QQ)$, both
spanned by algebraic classes, so $\pi_{k}$ is algebraic; likewise the part
$\pi_{N}$ of $\pi_{m}$ in $N\otimes N$, since $\psi$ is nondegenerate on $N$.
So $\pi_{V}=[\Delta]-\sum_{k\ne m}\pi_{k}-\pi_{N}$, the class of $\psi^{-1}$ in
$V\otimes V\subseteq H^{2m}(Y\times Y,\QQ)$, is algebraic.

The Hodge group of $Y^{n}$ is $\SO(V,\psi)$, acting on
$H^{*}(Y^{n},\QQ)=\bigl(H^{*}(Y,\QQ)\bigr)^{\otimes n}$ through the factors $V$
and trivially on the rest, which is spanned by algebraic classes; so the Hodge
classes are the $\SO(V)$-invariants. By the first fundamental theorem for the
special orthogonal group \cite[Chapter~5]{GW09}, the invariants of $\SO(V)$ in
$V^{\otimes j}$ are spanned by the complete contractions with $\psi^{-1}$ when
$j<t$, and, when $j=t$, by these and the determinant tensor; in general they
are spanned by products of contractions and determinant tensors on disjoint
sets of $t$ factors. A contraction of the factors $a$ and $b$ of $Y^{n}$ is the
pull-back of $\pi_{V}$ along the projection to those two factors, a
determinant on a set $I$ of $t$ factors is the pull-back of $\det(V)$ along
the projection to $Y^{I}$, and products of these with pull-backs of algebraic
classes span the Hodge classes. This gives (i), (ii) and the first half of
(iv). For (iii), the orthogonal group $\mathrm{O}(V,\psi)$ acts on
$H^{*}(Y^{j},\QQ)$ in the same way, and for $j<t$ every $\SO(V)$-invariant is
$\mathrm{O}(V)$-invariant; pull-backs along projections and cup products
commute with the diagonal action of $\mathrm{O}(V)$, so every element of the
algebra in (iii) is $\mathrm{O}(V)$-invariant, while a reflection acts on
$\det(V)$ by $-1$. For the second half of (iv), $\mathrm{Ab}^{*}(Y^{n})$
contains the algebraic classes and is stable under pull-back along projections
and under cup product: $\Gamma_{*}\gamma\cup\Gamma'_{*}\gamma'$ is the image of
$\gamma\otimes\gamma'$ on $B\times B'$ under the correspondence
$(\Gamma\times\Gamma')$ followed by the diagonal of $Y^{n}$.
\end{proof}

For a K3 surface $S$ with $\End_{\mathrm{Hdg}}(T(S))=\QQ$ the Hodge group of
$T(S)$ is $\SO(T(S))$ \cite{Zar83}; for a very general one, $t=21$. So the
Hodge conjecture holds for $S^{n}$ with $n\le20$, and on $S^{21}$ it is
equivalent to the algebraicity of the single class $\det(T(S))$. The argument
is the classical one from invariant theory, and we make no claim of priority
for it. What it adds here is the location of the first instance of
\textup{(F3$'$)} on the powers of such a variety that invariant theory does
not reach: one class, in one degree, on one power.

The same invariant theory, with the motivic decompositions of the
literature and the theorems of Buskin and Markman on Hodge isometries, gives
the Hodge conjecture itself, and so \textup{(F3$'$)}, on the other varieties
built from a K3 surface and on varieties of K3$^{[n]}$ type, up to a real
multiplication and one determinant class.

\begin{proposition}[Powers, Hilbert schemes and moduli spaces of a K3
surface]\label{prop:k3powers}
Let $S$ be a projective K3 surface, $T=T(S)_{\QQ}$, $t=\dim T$ and
$E=\End_{\mathrm{Hdg}}(T)$, a totally real field or a CM field \cite{Zar83}.
\begin{enumerate}[label=\textup{(\roman*)},leftmargin=2.6em,itemsep=2pt]
\item If $E$ is a CM field, the Hodge conjecture holds for every power
  $S^{k}$ \cite[Theorem~5.4]{RM08}, \cite{Bus19}, hence for every Hilbert
  scheme $S^{[n]}$ \cite{dCM02}, for every smooth projective moduli space of
  Gieseker stable (twisted) sheaves on $S$, or of $\sigma$-stable objects in
  $D^{b}(S,\alpha)$ for a generic stability condition $\sigma$ and a Brauer
  class $\alpha$ \cite[Theorem~0.1, Remark~2.1]{Bue20}, and for the crepant
  resolutions of the ten-dimensional moduli spaces $M_{\sigma}(2v_{0})$ of
  $\sigma$-semistable objects in $D^{b}(S,\alpha)$, with $v_{0}$ primitive,
  $v_{0}^{2}=2$ and $\sigma$ $v_{0}$-generic \cite[Theorem~1.3]{FFZ21}.
\item If $E=\QQ$, the Hodge conjecture holds for $S^{[n]}$ for every
  $n<t(t+1)/2$, and for every moduli space $M$ as in \textup{(i)} with
  $\dim M<t$. For $N=t(t+1)/2$ the following are equivalent: the Hodge
  conjecture for $S^{[N]}$; the algebraicity of the class $\det(T)$ on $S^{t}$
  of \Cref{prop:orthpowers}; the Hodge conjecture for every power of $S$; and
  the Hodge conjecture for every $S^{[n]}$. For a very general $S$, the Hodge
  conjecture holds for $S^{[n]}$ with $n\le230$, and for $S^{[231]}$ if and
  only if $\det(T(S))$ on $S^{21}$ is algebraic.
\item If $E$ is totally real of degree at least two, the Hodge conjecture for
  $S\times S$, the Hodge conjecture for $S^{[2]}$, and the statement that
  every element of $E$ acts on $T$ through an algebraic cycle on $S\times S$
  are equivalent.
\end{enumerate}
Wherever the Hodge conjecture holds, so does \textup{(F3$'$)}. This is a
consequence of \cite{RM08,Bus19,dCM02,Bue20,FFZ21} and
\Cref{prop:orthpowers}; we make no claim of priority.
\end{proposition}

\begin{proof}
We use three facts. First, $H^{*}(S,\QQ)=A\oplus T$ with
$A=H^{0}\oplus\NS(S)_{\QQ}\oplus H^{4}$ spanned by algebraic classes; the
K\"unneth components of the diagonal of $S$ are algebraic \cite{Mur90}, and
so is its part $\pi_{T}$ in $T\otimes T$, as in the proof of
\Cref{prop:orthpowers}. The Hodge classes of $S^{k}$ are the invariants in
$(A\oplus T)^{\otimes k}$ of the Hodge group of $T$, which acts trivially on
$A$. Second, if $H^{*}(X,\QQ)$ is spanned by the images of the
$H^{*}(Y_{i},\QQ)$ under algebraic correspondences and the Hodge conjecture
holds for every $Y_{i}$, it holds for $X$: by the semisimplicity argument of
the proof of \Cref{prop:f3prime}(iii) every Hodge class of $X$ is a sum of
images of Hodge classes of the $Y_{i}$. Third, the correspondences
$\Gamma_{\nu}$ of the proof of \Cref{prop:aclosure}(ii) give an isomorphism
of Chow motives between $h(S^{[n]})$ and
$\bigoplus_{\nu}h(S^{(\nu)})(l(\nu)-n)$ \cite{dCM02}, so each
$h(S^{(\nu)})(l(\nu)-n)$ is a direct summand of $h(S^{[n]})$, and the lift
$\widetilde\Gamma_{\nu*}$ takes the same value on a class of $S^{l(\nu)}$ and
on its average over $G_{\nu}$. Hence the Hodge
conjecture holds for $S^{[n]}$ if and only if, for every partition $\nu$ of
$n$, every $G_{\nu}$-invariant Hodge class of $S^{l(\nu)}$ is algebraic. By
\cite[Theorem~0.1, Remark~2.1]{Bue20} the Chow motive of a moduli space $M$
as in (i) is a direct summand of a sum of Tate twists of motives
$h(S^{k_{i}})$ with $1\le k_{i}\le\dim M$, and by
\cite[Corollary~4.5]{FFZ21} the Chow motive of a crepant resolution as in (i)
lies in the tensor subcategory generated by $h(S)$, whose objects are direct
summands of sums of
Tate twists of the $h(S^{k})$, the dual of $h(S)$ being $h(S)(2)$.

(i) Over $\CC$ the Hodge group of $T$ is a product of groups
$\GL(W_{\sigma})$ acting on
$T_{\CC}=\bigoplus_{\sigma}(W_{\sigma}\oplus W_{\sigma}^{\vee})$
\cite{Zar83}, and by the first fundamental theorem for the general linear
group \cite[Chapter~5]{GW09} its invariants in $T_{\CC}^{\otimes b}$ are
spanned by products of contractions of one factor with another. Under the
isomorphism $\End(T)\cong T\otimes T$, $\varphi\mapsto\tau_{\varphi}=
\sum_{i}t_{i}^{\vee}\otimes\varphi(t_{i})$, with $(t_{i}^{\vee})$ the dual
basis of a basis $(t_{i})$ for the intersection form, the Hodge classes of
$T\otimes T$ are the $\tau_{\varphi}$ with $\varphi\in E$, and these
contractions lie in their complex span. The rational span of the products of
the $\tau_{\varphi}$ on pairs of factors has the same complexification as the
space of invariants, so the Hodge classes of $S^{k}$ are spanned by such
products and pull-backs of algebraic classes. For $u\in E$ with $u\bar u=1$
the map $\id_{\NS}\oplus u$ is a rational Hodge isometry of
$H^{2}(S,\QQ)$, since the adjoint of $u$ is $\bar u$ \cite{Zar83}; it is
induced by an algebraic cycle \cite{Bus19}, and composing that cycle with
$\pi_{T}$ on both sides gives $\tau_{u}$, which is therefore algebraic. The
elements of norm one span $E$. Indeed, write $E=E_{0}(y_{0})$ with
$y_{0}^{2}=-\delta$, $\delta\in E_{0}$, and let $y\in E$ with $\bar y=-y$; for
$s\in\QQ$ the element $u_{s}=(1+sy)(1-sy)^{-1}$ has norm one. A linear form
on $E$ vanishing on all $u_{s}$ gives a rational function of $s$ vanishing on
$\QQ$, so it vanishes on the coefficients $1$, $2y$, $2y^{2}$ of the expansion
of $u_{s}$ in powers of $s$. As $y$ runs over $E_{0}y_{0}$, the $y$ span
$E_{0}y_{0}$ and the $y^{2}=-a^{2}\delta$, $a\in E_{0}$, span $E_{0}$, since
$2a=(a+1)^{2}-a^{2}-1$. So every $\tau_{u}$, $u\in E$, is algebraic, and the
Hodge conjecture holds for every $S^{k}$; this is the argument of
\cite[Lemma~5.3, Theorem~5.4]{RM08}, which takes the algebraicity of the
isometries from Mukai's announcement; \cite{Bus19} proves it in general.
We use \cite{RM08} only when $E$ is a CM field: for a totally real $E$ the
invariants of the Hodge group are not generated by those of degree two, since
they include determinants, as \Cref{prop:orthpowers} shows for $E=\QQ$, and
the powers do not reduce to the square. The rest of (i) follows
from the second and third facts.

(ii) Here the Hodge group of $T$ is $\SO(T)$ \cite{Zar83}, and $t\ge3$,
since a Hodge structure of K3 type of dimension two has complex
multiplication. The Hodge conjecture holds for $S^{k}$, $k<t$, by
\Cref{prop:orthpowers}(i), which gives the assertion on $M$ by the second
fact. Let $\nu$ be a partition of $n$ with $a_{j}$ parts equal to $j$, so
that $G_{\nu}=\prod_{j}\mathfrak{S}_{a_{j}}$ permutes the factors of
$S^{l(\nu)}$ within blocks of sizes $a_{j}$, one block for each distinct part
size. By the proof of \Cref{prop:orthpowers}, the Hodge classes of
$S^{l(\nu)}$ are spanned by algebraic classes and by products
$\beta=\det_{I}\cup\gamma$, where $\det_{I}$ is the pull-back of $\det(T)$
along the projection to a set $I$ of $t$ factors and $\gamma$ is pulled back
from the other factors. If two factors of $I$ lie in one block, the
transposition in $G_{\nu}$ exchanging them fixes $\gamma$ and changes the
sign of the alternating class $\det_{I}$, all classes involved being of even
degree; so it takes $\beta$ to $-\beta$, and the average of $\beta$ over
$G_{\nu}$ is zero. If $\nu$ has fewer than $t$ distinct part sizes, every set
of $t$ factors meets some block twice, so a $G_{\nu}$-invariant Hodge class,
being its own average, is the average of an algebraic class, and is
algebraic. A partition with $t$ distinct part sizes has
$n\ge1+2+\dots+t=t(t+1)/2$, so by the third fact the Hodge conjecture holds
for $S^{[n]}$ when $n<t(t+1)/2$ (\Cref{fig:k3threshold}). For $n=N$ the
partition $(1,2,\dots,t)$ has
$l(\nu)=t$ and $G_{\nu}=1$, so the Hodge class $\det(T)$ of $S^{t}$ is
algebraic if the Hodge conjecture holds for $S^{[N]}$. If $\det(T)$ is
algebraic, the Hodge conjecture holds for every power of $S$ by
\Cref{prop:orthpowers}(iv), and for every $S^{[n]}$ by the second fact; and
the conjecture for $S^{t}$ contains the algebraicity of $\det(T)$. A very
general $S$ has $E=\QQ$ and $t=21$, and $N=231$.

(iii) The Hodge structure $T$ has no nonzero Hodge class, so the Hodge
classes of $A\otimes T$ and $T\otimes A$ vanish, and those of $S\times S$
are the algebraic classes and the $\tau_{\varphi}$, $\varphi\in E$. A cycle
inducing $\varphi$ on $T$, composed with $\pi_{T}$ on both sides, gives
$\tau_{\varphi}$, and $\tau_{\varphi}$ induces $\varphi$ on $T$; this proves
the equivalence for $S\times S$. By \cite{dCM02},
$h(S^{[2]})=h(S^{(2)})\oplus h(S)(-1)$, and the Hodge conjecture holds for
$S$. The exchange of the two factors takes $\tau_{\varphi}$ to
$\tau_{\varphi^{*}}$, and the adjoint involution of a totally real $E$ is
trivial \cite{Zar83}, so every $\tau_{\varphi}$ is invariant under
$\mathfrak{S}_{2}$ and the Hodge conjecture for $S^{[2]}$ asks for the same
classes as that for $S\times S$.
\end{proof}

\begin{figure}[tbp]
\centering
  \includegraphics[width=\textwidth]{fig_k3threshold}
\caption{The threshold of \Cref{prop:k3powers}(ii). For $t=1,\dots,21$ the
bar rises to $t(t+1)/2$, the least $n$ that has a partition with $t$ distinct
part sizes, and is cut into the parts $1,2,\dots,t$ of the partition that
attains it. For a K3 surface $S$ with $E=\End_{\mathrm{Hdg}}(T(S))=\QQ$ and
$t=\dim T(S)$, which is at least $3$, the Hodge conjecture holds for
$S^{[n]}$ with $n$ below the bar; at its top it is equivalent to the
algebraicity of the class $\det T(S)$ on $S^{t}$ of \Cref{prop:orthpowers},
and to the Hodge conjecture for every power of $S$; above it, it follows
from that algebraicity. A very general $S$ has $t=21$: the
conjecture holds for $S^{[n]}$ with $n\le230$, and for $S^{[231]}$ it is
equivalent to the algebraicity of $\det T(S)$ on $S^{21}$, one of the
frontier cases of \textup{(F3$'$)} (\Cref{rem:f3primefrontier}).}
\label{fig:k3threshold}
\end{figure}

\begin{proposition}[Varieties of K3$^{[n]}$ type]\label{prop:k3ntype}
Let $X$ be a smooth projective hyperk\"ahler variety deformation equivalent
to the Hilbert scheme of $n\ge2$ points on a K3 surface, $q$ the
Beauville--Bogomolov form on $H^{2}=H^{2}(X,\QQ)$, $\NS=\NS(X)_{\QQ}$,
$T=\NS^{\perp}$ and $E=\End_{\mathrm{Hdg}}(T)$, a totally real field or a CM
field whose adjoint involution for $q$ is complex conjugation \cite{Zar83}.
Let $\mathrm{SH}(X)\subseteq H^{*}(X,\QQ)$ be the subalgebra generated by
$H^{2}$, and for $\varphi\in E$ with $\varphi^{*}=\varphi$ let
$c_{\varphi}=\sum_{i}t_{i}^{\vee}\cup\varphi(t_{i})\in H^{4}(X,\QQ)$, with
$(t_{i})$ a basis of $T$ and $(t_{i}^{\vee})$ its dual basis for $q$.
\begin{enumerate}[label=\textup{(\roman*)},leftmargin=2.6em,itemsep=2pt]
\item The Hodge classes of $\mathrm{SH}(X)$ are spanned by products of
  divisor classes and classes $c_{\varphi}$.
\item If $E$ is $\QQ$ or a CM field, every Hodge class of $\mathrm{SH}(X)$ is
  algebraic.
\item Let $n=2$. Then $\mathrm{SH}(X)=H^{*}(X,\QQ)$, with
  $b_{4}(X)=276=\dim\mathrm{Sym}^{2}H^{2}$, and $c_{2}(X)=\frac65q^{\vee}$,
  where $q^{\vee}\in\mathrm{Sym}^{2}H^{2}$ is dual to $q$. So the Hodge
  conjecture holds for $X$ when $E$ is $\QQ$ or a CM field. When $E$ is
  totally real of degree at least two, it holds for $X$ if and only if every
  element of $E$, acting on $T$ and by zero on $\NS$, is induced by an
  algebraic cycle on $X\times X$.
\end{enumerate}
This is a consequence of \cite{Mar24,CM13}; we make no claim of priority.
\end{proposition}

\begin{proof}
(i) Multiplication $\mathrm{Sym}^{k}H^{2}\to\mathrm{SH}^{2k}(X)$ is a
surjective morphism of polarisable Hodge structures, so by semisimplicity
every Hodge class of $\mathrm{SH}^{2k}(X)$ is the image of one of
$\mathrm{Sym}^{k}H^{2}=\bigoplus_{a+b=k}\mathrm{Sym}^{a}\NS\otimes
\mathrm{Sym}^{b}T$, whose Hodge classes are
$\mathrm{Sym}^{a}\NS\otimes(\mathrm{Sym}^{b}T)^{\Hdg}$. The isomorphism
$\End(T)\cong T\otimes T$ of the proof of \Cref{prop:k3powers}, now for $q$,
carries the adjoint to the exchange of the factors, so
$(\mathrm{Sym}^{2}T)^{\Hdg}$ consists of the tensors of the self-adjoint
$\varphi\in E$, whose images in $H^{4}$ are the $c_{\varphi}$. Over $\CC$ the
Hodge group of $T$ is a product of groups $\SO(T_{\sigma})$ on
$T_{\CC}=\bigoplus_{\sigma}T_{\sigma}$, with $\dim T_{\sigma}=\dim_{E}T\ge3$,
when $E$ is totally real, and of groups $\GL(W_{\sigma})$ on
$T_{\CC}=\bigoplus_{\sigma}(W_{\sigma}\oplus W_{\sigma}^{\vee})$ when $E$ is a
CM field \cite{Zar83}. By the first fundamental theorems
\cite[Chapter~5]{GW09}, the invariant polynomials of one vector of
$T_{\sigma}$, or of one vector of $W_{\sigma}$ and one of $W_{\sigma}^{\vee}$,
are generated by the quadratic form, or by the pairing: a determinant needs
$\dim T_{\sigma}$ distinct vectors. So $(\mathrm{Sym}^{b}T)^{\Hdg}\otimes\CC$
is spanned by products of elements of $(\mathrm{Sym}^{2}T)^{\Hdg}\otimes\CC$,
and the rational span of the products of classes $c_{\varphi}$, which has the
same complexification, is $(\mathrm{Sym}^{b}T)^{\Hdg}$.

(ii) Divisor classes are algebraic, so by (i) it suffices that every
$c_{\varphi}$ is. By \cite[Theorem~1.1]{CM13} the Lefschetz standard
conjecture holds for $X$, so the K\"unneth projectors and the inverse of the isomorphism
$L^{2n-2}\colon H^{2}\to H^{4n-2}$, $L$ the cup product with an ample class
$h$, are induced by algebraic cycles \cite{Kle68,Kle94}. Let $u\in E$ with
$u\bar u=1$, and $u=1$ if $E=\QQ$. The map $f_{u}=\id_{\NS}\oplus u$ is a
rational Hodge isometry of $H^{2}$, so by \cite[Theorem~1.1]{Mar24} there is
an algebraic class $Z_{u}$ on $X\times X$ (the diagonal if $u=1$) inducing on
$H^{2}$ a nonzero rational multiple of $f_{u}$: the class of the
correspondence $\tilde f$ of that theorem if it preserves degrees, and the
class of $\tilde f\circ(c_{2}(X)^{n-1}\cup\,\cdot\,)$ if it reverses them.
The class $Z_{u}$ may act in other degrees too; its K\"unneth component in
$H^{4n-2}(X)\otimes H^{2}(X)$ is algebraic and induces a nonzero multiple of
$f_{u}$ on $H^{2}$ and zero elsewhere. Applying $(L^{2n-2})^{-1}\otimes\id$
gives, up to a nonzero rational factor, the algebraic class
$\sum_{i}b_{i}\otimes f_{u}(a_{i})$, for a basis $(a_{i})$ of $H^{2}$ and
the dual basis $(b_{i})$ for $B(x,y)=\int_{X}xyh^{2n-2}$. Polarising the
relation $\int_{X}\alpha^{2n}=c_{X}q(\alpha)^{n}$ of Fujiki and Beauville
\cite{Bea83,Fuj87} gives
\[
  B(x,y)=\frac{c_{X}}{2n-1}\,q(h)^{n-2}\bigl(q(h)q(x,y)+2(n-1)q(x,h)q(y,h)\bigr),
\]
so $\NS$ and $T$ are orthogonal for $B$, and $B=\kappa q$ on $T$ with
$\kappa=c_{X}q(h)^{n-1}/(2n-1)\ne0$. With $(a_{i})$ adapted to
$\NS\oplus T$, the pull-back along the diagonal is
$\sum_{i}b_{i}\cup f_{u}(a_{i})=\nu+\kappa^{-1}c_{(u+\bar u)/2}$ with
$\nu\in\mathrm{Sym}^{2}\NS$ algebraic, because cup product on $H^{2}$ is
commutative and the exchange of factors turns $u$ into $\bar u$. So
$c_{(u+\bar u)/2}$ is algebraic. When $E$ is a CM field its elements of norm
one span it (proof of \Cref{prop:k3powers}(i)), so the $(u+\bar u)/2$ span
$E_{0}$, the self-adjoint part of $E$; when $E=\QQ$, $u=1$ suffices. So
every $c_{\varphi}$ is algebraic.

(iii) The variety $X$ has the Betti numbers of $S^{[2]}$, which by
\cite{dCM02} are those of $\mathrm{Sym}^{2}H^{*}(S)\oplus H^{*}(S)(-1)$:
$H^{\mathrm{odd}}(X)=0$ and $b_{4}(X)=253+1+22=276=\dim\mathrm{Sym}^{2}H^{2}$.
Here $c_{X}=3$, and the relation of Fujiki and Beauville gives
$\int_{X}abcd=q(a,b)q(c,d)+q(a,c)q(b,d)+q(a,d)q(b,c)$ for $a,b,c,d\in H^{2}$.
This form on $\mathrm{Sym}^{2}H^{2}$ is nondegenerate: in a $q$-orthogonal
basis $(e_{i})$ with $q(e_{i},e_{i})=d_{i}$, the $e_{i}e_{j}$ with $i<j$ are
orthogonal to each other and to the squares, with $\int_{X}(e_{i}e_{j})^{2}=
d_{i}d_{j}$, and on the squares the form has the matrix $D(J+2I)D$, with
$D=\mathrm{diag}(d_{i})$ and $J$ the matrix of ones, of determinant
$2^{22}\cdot25\prod_{i}d_{i}^{2}$. So
$\mathrm{Sym}^{2}H^{2}\to H^{4}$ is injective, hence bijective, and with
$H^{6}=h^{2}\cup H^{2}$ by the hard Lefschetz theorem,
$\mathrm{SH}(X)=H^{*}(X,\QQ)$. By \cite{Fuj87},
$\int_{X}c_{2}(X)xy=C\,q(x,y)$ for a constant $C$, and the Riemann--Roch
polynomial $\chi(L)=\binom{q(L)/2+3}{2}$ \cite{EGL01} gives $C=30$; the
Fujiki relation gives $\int_{X}q^{\vee}xy=25\,q(x,y)$, so
$c_{2}(X)=\frac65q^{\vee}$ by nondegeneracy. For $E=\QQ$ this makes
$c_{1}=q^{\vee}-q^{\vee}_{\NS}$ algebraic without \cite{CM13}, with
$q^{\vee}_{\NS}\in\mathrm{Sym}^{2}\NS$ the class dual to $q|_{\NS}$. With (i) and (ii) this proves the Hodge
conjecture for $X$ when $E$ is $\QQ$ or a CM field. Let $E$ be totally real
and $\varphi\in E$. If $\varphi\oplus0$ is induced by an algebraic cycle, the
computation of (ii) with $\varphi\oplus0$ in place of $f_{u}$ makes
$c_{\varphi}$ algebraic. Conversely, if $c_{\varphi}$ is algebraic, so is the
correspondence $x\mapsto(L^{2})^{-1}(c_{\varphi}\cup x)$. The Fujiki relation
gives $\int_{X}c_{\varphi}xy=\Tr(\varphi)q(x,y)+2q(\varphi x,y)$ for $x\in T$
and $y\in H^{2}$, and $B(z,y)=q(h)q(z,y)$ for $z\in T$, so that
correspondence acts on $T$ as $(\Tr(\varphi)+2\varphi)/q(h)$. The projector
onto $T$, the K\"unneth projector $\pi_{2}$ minus
$\sum_{j}n_{j}^{\vee}\otimes n_{j}$ over a basis $(n_{j})$ of $\NS$ and the
Hodge classes $n_{j}^{\vee}\in H^{6}$ dual to it, is algebraic by
\Cref{prop:f3prime}(iv); composing with it, $\varphi\oplus0$ is induced by an
algebraic cycle. By (i) the Hodge conjecture for $X$ is the algebraicity of
the $c_{\varphi}$, which proves the equivalence.
\end{proof}

For varieties of generalised Kummer type every Hodge class in the subalgebra
generated by $H^{2}$ is algebraic, whatever the field, and the Hodge
conjecture holds in dimension four \cite{FV25}. The preprint \cite{OG26}
announces, with \cite{Mar24}, that rational Hodge isometries between
varieties of the types K3$^{[a_{1}^{2}+1]}$ and K3$^{[a_{2}^{2}+1]}$ are
induced by algebraic cycles; it is not used here.

\begin{corollary}[Cubic fourfolds and their Fano varieties of lines]
\label{cor:fanolines}
Let $Y\subset\PP^{5}$ be a smooth cubic fourfold, $T(Y)\subset
H^{4}(Y,\QQ)$ the orthogonal complement of the Hodge classes,
$E(Y)=\End_{\mathrm{Hdg}}(T(Y))$, and $F=F(Y)$ the Fano variety of lines on
$Y$. The Hodge conjecture for $F$, the Hodge conjecture for $Y\times Y$, and
the statement that every element of $E(Y)$ acts on $T(Y)$ through an
algebraic cycle on $Y\times Y$ are equivalent. They hold when $E(Y)$ is $\QQ$
or a CM field, for instance when $Y$ is very general.
\end{corollary}

\begin{proof}
The Hodge conjecture holds for $Y$ \cite{Zuc77}, and also by
\Cref{prop:f3primesmall}, since $Y$ is covered by lines. So $H^{*}(Y,\QQ)$
is $T(Y)$ plus a space spanned by algebraic classes, on which the intersection
form is nondegenerate by the Hodge--Riemann relations, and as in the proof of
\Cref{prop:orthpowers} the part $\pi_{T}$ of the diagonal in
$T(Y)\otimes T(Y)$ is algebraic. The Hodge structure $T(Y)$ is of K3 type up
to a Tate twist and has no nonzero Hodge class, so $E(Y)$ is a field
\cite{Zar83}, and as in \Cref{prop:k3powers}(iii) the Hodge conjecture for
$Y\times Y$ is the third statement. By \cite{BD85}, $F$ is a hyperk\"ahler
fourfold of K3$^{[2]}$ type, and the universal line $P\subset F\times Y$
induces an isomorphism of Hodge structures $\alpha=P_{*}\colon
H^{4}(Y,\QQ)\to H^{2}(F,\QQ)$ up to a Tate twist. Both $T(Y)$ and $T(F)$
are the sums of the isotypic components that are not of Tate type, so
$\alpha(T(Y))=T(F)$ and $E(F)=\alpha E(Y)\alpha^{-1}$. Let $g$ be the
Pl\"ucker class on $F$ and $e=P^{t}_{*}\circ(g^{2}\cup\,\cdot\,)
\circ\alpha$, an algebraic correspondence that maps $T(Y)$ to itself, so that
$e|_{T(Y)}\in E(Y)$. For $0\ne\sigma\in H^{3,1}(Y)$,
$\int_{Y}e(\sigma)\bar\sigma=\int_{F}g^{2}\alpha(\sigma)\alpha(\bar\sigma)
=q(g)\,q(\alpha(\sigma),\alpha(\bar\sigma))\ne0$ by the Fujiki relation, since
$q(g)>0$ and $\alpha(\sigma)$ spans $H^{2,0}(F)$. So $e|_{T(Y)}$ is
invertible, its inverse is a polynomial in it and is algebraic, and the
inverse of $\alpha$ on $T(F)$,
which is $e^{-1}\circ P^{t}_{*}\circ(g^{2}\cup\,\cdot\,)$, is induced by an
algebraic cycle. Hence $\varphi\in E(Y)$ acts through an algebraic cycle on
$Y\times Y$ if and only if $\alpha\varphi\alpha^{-1}$ acts through one on
$F\times F$, and \Cref{prop:k3ntype}(iii) proves the equivalences and the
case where $E(Y)=E(F)$ is $\QQ$ or a CM field. A very general $Y$ has
$E(Y)=\QQ$.
\end{proof}

The class $\mathcal{A}$ of \Cref{prop:f3prime}(iii) is larger than the
closure properties recorded there show. Four further operations
preserve it, the last of them domination by correspondences, which brings in
the varieties whose motives are known to be built from those of abelian
varieties, and two classical families lie in it.

\begin{proposition}[More varieties of $\mathcal{A}$]\label{prop:aclosure}
Let $\mathcal{A}$ be as in \Cref{prop:f3prime}(iii).
\begin{enumerate}[label=\textup{(\roman*)},leftmargin=2.6em,itemsep=2pt]
\item If $Y\in\mathcal{A}$ and $Z\subset Y$ is a smooth closed subvariety
  with $Z\in\mathcal{A}$, then the blow-up $\mathrm{Bl}_{Z}Y$ lies in
  $\mathcal{A}$. If $Y\in\mathcal{A}$ and $V$ is a vector bundle on $Y$,
  then $\PP(V)\in\mathcal{A}$.
\item If $S\in\mathcal{A}$ is a surface, then the Hilbert scheme $S^{[n]}$
  lies in $\mathcal{A}$ for every $n\ge1$.
\item Every Fermat hypersurface
  $X^{r}_{m}=\{x_{0}^{m}+\dots+x_{r+1}^{m}=0\}\subset\PP^{r+1}$, $r\ge0$,
  $m\ge1$, lies in $\mathcal{A}$.
\item A K3 surface $S$ lies in $\mathcal{A}$ as soon as its Kuga--Satake
  correspondence is induced by an algebraic cycle, that is, as soon as some
  algebraic cycle $\Gamma$ on $S\times B$, with $B$ an abelian variety,
  induces a nonzero map $T(S)_{\QQ}\to H^{2}(B,\QQ)$.
\item If $H^{*}(X,\QQ)$ is the sum of the images of the $H^{*}(Y_{i},\QQ)$,
  with $Y_{i}\in\mathcal{A}$, under algebraic correspondences, then
  $X\in\mathcal{A}$; in particular $X\in\mathcal{A}$ when its Chow motive is
  a direct summand of a sum of Tate twists of motives of members of
  $\mathcal{A}$, for instance when it is of abelian type. So $\mathcal{A}$
  contains the $2n$-dimensional generalised Kummer varieties $K_{n}(A)$ of
  abelian surfaces $A$ \cite[Theorem~2.7, Corollary~2.8]{Xu15},
  \cite[Theorem~1.5]{FTV19}; the smooth projective moduli spaces of Gieseker
  stable (twisted) sheaves, or of $\sigma$-stable objects for a generic
  stability condition $\sigma$, on abelian surfaces and on K3 surfaces in
  $\mathcal{A}$ \cite[Theorem~0.1, Remark~2.1]{Bue20}; O'Grady's
  six-dimensional varieties $\widetilde K_{v}(A,H)$ obtained from moduli
  spaces of sheaves on an abelian surface $A$, with $v=2v_{0}$, $v_{0}$
  primitive, effective and of square $2$, and $H$ a $v$-generic polarisation
  \cite[Theorem~1.1]{Flo23}; the crepant resolutions of the ten-dimensional
  moduli spaces $M_{\sigma}(2v_{0})$ of $\sigma$-semistable objects, with
  $v_{0}$ primitive, $v_{0}^{2}=2$ and $\sigma$ $v_{0}$-generic, on abelian
  surfaces and on K3 surfaces in $\mathcal{A}$
  \cite[Theorem~1.3, Corollary~4.5]{FFZ21}; and the cubic fourfolds whose
  lattice of algebraic $2$-cycles has rank greater than $19$, with their Fano
  varieties of lines, the eightfolds of Lehn, Lehn, Sorger and van Straten of
  those containing no plane, and the tenfolds of Laza, Sacc\`a and Voisin of
  infinitely many Pfaffian cubic fourfolds \cite[Theorem~3.5 and its proof,
  Proposition~3.9, Theorem~4.2]{ABLP26}.
\end{enumerate}
Consequently \textup{(F3$'$)} holds on every smooth projective variety
obtained from abelian varieties, curves, Fermat hypersurfaces and K3 surfaces
as in \textup{(iv)} by products, surjective images, blow-ups along smooth
centres in $\mathcal{A}$, projective bundles and Hilbert schemes of points on
surfaces, and on every variety dominated by these as in \textup{(v)}.
\end{proposition}

\begin{proof}
(i) Let $c$ be the codimension of $Z$, $E=\PP(N_{Z/Y})$ the exceptional
divisor, $j\colon E\to\widetilde Y=\mathrm{Bl}_{Z}Y$ its inclusion,
$\rho\colon E\to Z$ and $\pi\colon\widetilde Y\to Y$ the projections and
$\xi=c_{1}(\cO_{E}(1))$. By the blow-up formula \cite[\S7.3.3]{Voi02},
\[
  H^{*}(\widetilde Y,\QQ)=\pi^{*}H^{*}(Y,\QQ)\oplus
  \bigoplus_{i=0}^{c-2}j_{*}\bigl(\xi^{i}\cup\rho^{*}H^{*}(Z,\QQ)\bigr).
\]
The first summand is the image of the transpose of the graph of $\pi$, and
the $i$-th of the others is the image of the algebraic cycle
$(\rho\times j)_{*}\xi^{i}$ on $Z\times\widetilde Y$. Composing with the
correspondences that exhibit $Y$ and $Z$ as members of $\mathcal{A}$, which is
allowed \cite[Chapter~16]{Ful98}, shows that $H^{*}(\widetilde Y,\QQ)$ is
spanned by classes $\Gamma_{*}y$ with $B$ abelian. For $\PP(V)$ the
projective bundle formula $H^{*}(\PP(V),\QQ)=\bigoplus_{i}\xi^{i}\cup
\pi^{*}H^{*}(Y,\QQ)$ \cite[\S7.3.3]{Voi02} does the same with the
algebraic cycles $\Gamma_{\pi}^{t}\cdot\mathrm{pr}^{*}\xi^{i}$.

(ii) By de Cataldo and Migliorini \cite{dCM02}, for every partition $\nu$ of
$n$ with $l(\nu)$ parts the incidence correspondence
$\Gamma_{\nu}\subset S^{(\nu)}\times S^{[n]}$ between the product
$S^{(\nu)}$ of symmetric products attached to $\nu$ and the Hilbert scheme
induces, after a Tate twist by $n-l(\nu)$, an isomorphism
$\bigoplus_{\nu}H^{*}(S^{(\nu)},\QQ)\to H^{*}(S^{[n]},\QQ)$. Let
$q_{\nu}\colon S^{l(\nu)}\to S^{(\nu)}$ be the quotient by the finite group
$G_{\nu}$. With rational coefficients the cycles and the cohomology of
$S^{(\nu)}$ are the $G_{\nu}$-invariants of those of $S^{l(\nu)}$
\cite[Example~1.7.6]{Ful98}, so $\Gamma_{\nu}$ pulls back to a cycle
$\widetilde\Gamma_{\nu}$ with rational coefficients on the smooth variety
$S^{l(\nu)}\times S^{[n]}$, whose action on $H^{*}(S^{l(\nu)},\QQ)$ is
$\Gamma_{\nu*}\circ q_{\nu*}$ up to a nonzero factor. As
$H^{*}(S^{(\nu)},\QQ)=q_{\nu*}H^{*}(S^{l(\nu)},\QQ)$, the $\nu$-th summand
is the image of $H^{*}(S^{l(\nu)},\QQ)$ under $\widetilde\Gamma_{\nu}$. Products of members of $\mathcal{A}$ lie in
$\mathcal{A}$, so $S^{[n]}\in\mathcal{A}$.

(iii) By induction on $r$, following the inductive structure of Shioda and
Katsura \cite{SK79}. For $r=0$, $X^{0}_{m}$ is a finite set of points, and
$X^{1}_{m}$ is a curve. Let $r\ge2$ and $r=a+b$ with $a,b\ge1$; fix
$\epsilon$ with $\epsilon^{m}=-1$ and put
\[
  \varphi\bigl((x),(y)\bigr)=\bigl(x_{0}y_{b+1}:\dots:x_{a}y_{b+1}:
  \epsilon x_{a+1}y_{0}:\dots:\epsilon x_{a+1}y_{b}\bigr),
\]
a rational map $X^{a}_{m}\times X^{b}_{m}\dashrightarrow\PP^{r+1}$; since
$\sum_{i\le a}x_{i}^{m}=-x_{a+1}^{m}$ and $\sum_{k\le b}y_{k}^{m}=-y_{b+1}^{m}$,
its image lies in $X^{r}_{m}$. It is defined by sections of
$\cO(1,1)$ whose common zero scheme is $\{x_{a+1}=0\}\times\{y_{b+1}=0\}$,
the product $Z\cong X^{a-1}_{m}\times X^{b-1}_{m}$ of two smooth
hyperplane sections, transversal because the Fermat hypersurfaces
are smooth. Blowing up that reduced scheme resolves the map, so $\varphi$
extends to a morphism $\widetilde\varphi\colon\mathrm{Bl}_{Z}(X^{a}_{m}\times
X^{b}_{m})\to X^{r}_{m}$. It is dominant: at a point $z$ of $X^{r}_{m}$ with
$\sum_{i\le a}z_{i}^{m}\ne0$ the preimages with $x_{a+1}=y_{b+1}=1$ are the
points $x_{i}=\lambda z_{i}$, $y_{k}=\lambda z_{a+1+k}/\epsilon$ with
$\lambda^{m}\sum_{i\le a}z_{i}^{m}=-1$, and both equations reduce to this one
because $\sum_{\text{all}}z_{i}^{m}=0$. So $\widetilde\varphi$ is surjective,
being proper. By induction $X^{a}_{m}$, $X^{b}_{m}$, $X^{a-1}_{m}$ and
$X^{b-1}_{m}$ lie in $\mathcal{A}$, hence so do their products, the blow-up
by (i) and $X^{r}_{m}$ as a surjective image; one may take $a=r-1$, $b=1$.

(iv) The transcendental part $T=T(S)_{\QQ}$ is an irreducible Hodge
structure, and $H^{*}(S,\QQ)=T\oplus T^{\perp}$ with $T^{\perp}$ spanned by
algebraic classes. Let $h$ be an ample class on $B$, $d=\dim B$, and let
$\phi=\Gamma^{t}_{*}\circ(\,\cdot\cup h^{d-2})\circ\Gamma_{*}$, an algebraic
correspondence from $S$ to itself and a morphism of Hodge structures. It
preserves $T$, because a morphism from the irreducible structure $T$ to the
algebraic part, which is of Tate type, vanishes; so $\phi|_{T}$ lies in the
field $\End_{\mathrm{Hdg}}(T)$. For $0\ne\sigma\in H^{2,0}(S)$,
$\Gamma_{*}\sigma\ne0$ because $T$ is irreducible and $\Gamma_{*}|_{T}\ne0$;
it is a class of type $(2,0)$ on $B$, hence primitive, and by the
Hodge--Riemann relations $\int_{S}\phi(\sigma)\,\bar\sigma=\int_{B}
\Gamma_{*}\sigma\,\overline{\Gamma_{*}\sigma}\,h^{d-2}\ne0$. So
$\phi|_{T}\ne0$, it is invertible, and $T\subseteq\Gamma^{t}_{*}H^{*}(B,\QQ)$.
With $T^{\perp}$ algebraic, $S\in\mathcal{A}$.

(v) If $H^{*}(X,\QQ)=\sum_{i}\Psi_{i*}H^{*}(Y_{i},\QQ)$ and each
$H^{*}(Y_{i},\QQ)$ is spanned by classes $\Gamma_{ij*}y$ with $B_{ij}$
abelian, then $H^{*}(X,\QQ)$ is spanned by the classes
$(\Psi_{i}\circ\Gamma_{ij})_{*}y$, and composites of algebraic
correspondences are algebraic \cite[Chapter~16]{Ful98}. If $h(X)$ is a
direct summand of $\bigoplus_{i}h(Y_{i})(n_{i})$, through $\alpha$ and
$\beta$ with $\beta\circ\alpha=\id$, then $\beta_{*}$ is surjective on
cohomology, and its components are the $\Psi_{i}$; a motive of abelian type
is such a summand with abelian varieties as the $Y_{i}$. The objects of the
tensor subcategory generated by $h(Y)$ are direct summands of sums of Tate
twists of the motives $h(Y^{k})$, since the dual of $h(Y)$ is a Tate twist of
$h(Y)$, and $Y^{k}\in\mathcal{A}$ when $Y\in\mathcal{A}$. The Chow motive of
$K_{n}(A)$ lies in the tensor subcategory generated by $h(A)$
\cite[Theorem~2.7, Corollary~2.8]{Xu15}, see also \cite[Theorem~1.5]{FTV19};
that of a moduli space as in (v) is a direct summand of a sum of Tate twists
of motives of powers of the surface \cite[Theorem~0.1]{Bue20}; those of the
varieties $\widetilde K_{v}(A,H)$ and of the ten-dimensional crepant
resolutions lie in the tensor subcategory generated by the motive of the
surface \cite[Theorem~1.1]{Flo23}, \cite[Corollary~4.5]{FFZ21}; and the cubic
fourfolds and hyperk\"ahler varieties of \cite{ABLP26} listed in (v) have
Chow motives of abelian type \cite[Theorem~3.5 and its proof,
Proposition~3.9, Theorem~4.2]{ABLP26}.

The last assertion is \Cref{prop:f3prime}(iii).
\end{proof}

The known cases of (iv) are recalled in \Cref{rem:mumfordks}; for the K3
surfaces among them whose powers are covered by \cite{Var22} the Hodge
conjecture itself holds on every power, which is stronger. What (iii) adds is
a family on which \textup{(F3$'$)} holds in every degree and dimension, while
the Hodge conjecture for Fermat hypersurfaces is proved in a range of degrees
and dimensions \cite{Shi79b}. Three recent preprints, if
correct, extend that range to every Fermat variety of degree less than $65$
other than $44$, $51$ and $52$ \cite{MMMV26}, to the Fermat fourfolds of
every odd degree at most $199$, by a computer-assisted proof \cite{Jum26},
and to the Fermat fourfold of degree $33$ \cite{FK26}; (iii) is new for the
degrees and dimensions outside these ranges. By (iii) and
\Cref{prop:f3prime}(i) the conjecture holds for $X^{r}_{m}$ as soon as it
holds for the powers of the Jacobian of the Fermat curve of degree $m$. The
inductive structure is Shioda and Katsura's, and the consequence for
\textup{(F3$'$)} is immediate; we claim no priority for it.

Item (v) collects consequences of published motivic results, and we claim no
priority for them either. For a generalised Kummer variety $K_{n}(A)$ it is
needed in general. The restriction from $A^{[n+1]}$, which lies in
$\mathcal{A}$ by (ii), reaches exactly the $A[n+1]$-invariant part of
$H^{*}(K_{n}(A),\QQ)$, since $A\times K_{n}(A)\to A^{[n+1]}$ is an \'etale
Galois cover with group $A[n+1]$. When $n+1$ is prime the rest is spanned by
the classes of the punctual Hilbert schemes at the points of $A[n+1]$, which
are algebraic, and (ii) suffices; when $n+1$ is composite the rest contains
transcendental classes, and \cite{Xu15,FTV19} are needed. These results
concern the varieties $K_{n}(A)$ themselves, not every variety of
generalised Kummer type; for products of varieties $K_{n}(A)$ the Hodge
conjecture itself is known \cite[Theorem~3.3]{Xu15}, and for the varieties
$\widetilde K_{v}(A,H)$ it holds on all self-products
\cite[Corollary~1.2]{Flo23}. If the
preprint \cite{Ped26} is correct, (v) also puts in $\mathcal{A}$ the tenfolds
of Laza, Sacc\`a and Voisin of every cubic fourfold in $\mathcal{A}$, whose
motives it makes direct summands of Tate twists of the motive of the fifth
power of the cubic fourfold. In low dimension the algebraicity results of
Bloch--Srinivas and Conte--Murre give more.

\begin{proposition}[Varieties whose zero-cycles are small]
\label{prop:f3primesmall}
Let $X$ be a smooth complex projective variety of dimension $n$ whose group
$\mathrm{CH}_{0}(X)_{\QQ}$ is supported on a closed subvariety of dimension at
most three. If $n\le5$, the Hodge conjecture holds for $X$, and so does
\textup{(F3$'$)}. This applies to every fourfold whose group of zero-cycles
is supported on a proper closed subset, in particular every uniruled
fourfold, and to every fivefold whose maximal rationally connected quotient
has dimension at most three, in particular every rationally connected
fivefold.
\end{proposition}

\begin{proof}
By Bloch and Srinivas \cite{BS83} every Hodge class of degree $4$ on $X$ is
algebraic. By \Cref{prop:f3prime}(iv) the Hodge classes of degree $0$, $2$,
$2n-2$ and $2n$ are algebraic, which settles $n\le4$. For $n=5$ the remaining
degree is $6=2n-4$, and cup product with $h$, for $h$ ample, is an isomorphism
$H^{4}(X,\QQ)\to H^{6}(X,\QQ)$ of Hodge structures up to a Tate twist by hard
Lefschetz, so every Hodge class of degree $6$ is $h$ times one of degree $4$.
A uniruled fourfold has $\mathrm{CH}_{0}$ supported on any ample divisor $D$,
because a rational curve through a point meets $D$ and any two points of
$\PP^{1}$ are rationally equivalent; this case is the theorem of Conte and
Murre \cite{CM78}. A proper closed subset of a fourfold has dimension at
most three.

Let $\phi\colon X\dashrightarrow R$ be the maximal rationally connected
fibration of Campana and of Koll\'ar, Miyaoka and Mori \cite{Cam92,KMM92}:
$R$, the maximal rationally connected quotient, is smooth and projective and
unique up to birational equivalence, and $\phi$ is dominant with connected
fibres, the general one rationally connected. Resolving its indeterminacy
\cite{Hir64} gives a smooth projective variety $\widetilde X$, a birational
morphism $\pi\colon\widetilde X\to X$, an isomorphism over the domain of
definition of $\phi$, and a surjective
morphism $\widetilde\phi\colon\widetilde X\to R$ whose general fibre is
smooth, by generic smoothness, and birational to a general fibre of $\phi$,
hence rationally connected. Suppose $\dim R\le3$, and let
$W\subset\widetilde X$ be a subvariety that $\widetilde\phi$ maps generically
finitely onto $R$, of degree $e$. Any two points of a general fibre $F$ are
joined by a rational curve, so every $x\in F$ satisfies $e[x]=[W\cap F]$ in
$\mathrm{CH}_{0}(\widetilde X)_{\QQ}$. The general fibres cover a dense open
subset of $\widetilde X$, and every zero-cycle is rationally equivalent to
one supported on a given dense open subset, so
$\mathrm{CH}_{0}(\widetilde X)_{\QQ}$ is supported on $W$ and
$\mathrm{CH}_{0}(X)_{\QQ}=\pi_{*}\mathrm{CH}_{0}(\widetilde X)_{\QQ}$ on
$\pi(W)$, of dimension $\dim R\le3$. A rationally connected fivefold is the
case where $R$ is a point.
\end{proof}

In degree four the hypothesis on zero-cycles can be relaxed to one on the
fibres of a map.

\begin{proposition}[Degree four along rationally connected fibrations]
\label{prop:f3primefibration}
Let $f\colon X\to Y$ be a surjective morphism of smooth complex projective
varieties whose general fibre is rationally connected. Then
$\Hdg^{2}(X)=\mathrm{Ab}^{2}(X)$ if and only if
$\Hdg^{2}(Y)=\mathrm{Ab}^{2}(Y)$, and every Hodge class of degree four on $X$
is algebraic if and only if every one on $Y$ is. Both properties are
therefore birational invariants of smooth projective varieties, and they hold
for $X$ exactly when they hold for a smooth projective model of its maximal
rationally connected quotient. Consequently \textup{(F3$'$)} holds for every
fourfold and every fivefold whose maximal rationally connected quotient is
birational to a variety $R$ with $\Hdg^{2}(R)=\mathrm{Ab}^{2}(R)$, for
instance a member of $\mathcal{A}$. With the Hodge conjecture for abelian
fourfolds \cite{Mar25a,Mar25b}, the Hodge conjecture holds for every fourfold
birational to an abelian fourfold, and for every fivefold admitting a
dominant rational map with rationally connected general fibre onto an
abelian fourfold.
\end{proposition}

\begin{proof}
Let $n=\dim X$ and $m=\dim Y$. Let $W$ be a component, dominating $Y$, of the
intersection of $X$ with a general linear space of codimension $n-m$, so that
$W\to Y$ is generically finite, of degree $e\ge1$; let
$i_{W}\colon\widetilde W\to X$ be a resolution of $W$ \cite{Hir64} and
$g=f\circ i_{W}$. The correspondence
$\Phi=\Gamma_{f}^{t}\circ\Gamma_{g}\circ\Gamma_{i_{W}}^{t}$ on $X\times X$
acts by $\Phi_{*}=f^{*}g_{*}i_{W}^{*}$ \cite[Chapter~16]{Ful98}. Let $\eta$ be
the generic point of $X$ and $K$ its field of functions. The restriction of
$\Phi^{t}$ to $\{\eta\}\times X$ is the zero-cycle $W_{f(\eta)}$, of degree
$e$, that $W$ cuts on the fibre of $X_{K}$ over the point $f(\eta)$, and that
of $[\Delta_{X}]$ is the point $\eta$ of the same fibre. Any two points of a
general fibre of $f$ are rationally equivalent, and the argument of Bloch and
Srinivas \cite{BS83} (a countable algebraically closed field of definition,
an embedding of an algebraic closure of $K$ into $\CC$, specialisation and a
norm) gives $e[\eta]=[W_{f(\eta)}]$ in $\mathrm{CH}_{0}(X_{K})_{\QQ}$. As
$\mathrm{CH}_{0}(X_{K})$ is the direct limit of the groups
$\mathrm{CH}^{n}(U\times X)$ over the dense open subsets $U\subseteq X$, the
cycle $e[\Delta_{X}]-\Phi^{t}$ vanishes on some $U\times X$, and by the
localisation sequence \cite[Chapter~1]{Ful98} it is supported on $D\times X$
for a divisor $D\supseteq X\smallsetminus U$. The diagonal is symmetric, so
$e[\Delta_{X}]-\Phi$ is supported on $X\times D$. Let $i\colon D'\to X$ be a
resolution of $D$. Each component of a cycle supported on $X\times D$ is, up
to a nonzero factor, the image of a subvariety of $X\times D'$ mapping
generically finitely onto it, so $e[\Delta_{X}]-\Phi=(\id\times i)_{*}\Psi$
for a cycle $\Psi$ with rational coefficients on $X\times D'$, and the
projection formula gives, for every $a\in H^{k}(X,\QQ)$,
\[
  e\,a=f^{*}g_{*}i_{W}^{*}a+i_{*}\Psi_{*}a,\qquad
  \Psi_{*}a\in H^{k-2}(D',\QQ).
\]

Let $a\in\Hdg^{2}(X)$. Then $g_{*}i_{W}^{*}a\in\Hdg^{2}(Y)$, and
$\Psi_{*}a$ is a Hodge class of degree two on $D'$, algebraic by the
Lefschetz theorem on $(1,1)$ classes. If $\Hdg^{2}(Y)=\mathrm{Ab}^{2}(Y)$,
then $f^{*}g_{*}i_{W}^{*}a\in\mathrm{Ab}^{2}(X)$, because composites of
algebraic correspondences are algebraic, and so $a\in\mathrm{Ab}^{2}(X)$; if
every Hodge class of degree four on $Y$ is algebraic, $a$ is algebraic.
Conversely, as in the proof of \Cref{prop:f3prime}(iii), for $h$ ample on
$X$ the correspondence $[\Gamma_{f}]\cdot\mathrm{pr}_{X}^{*}h^{n-m}$ sends
$f^{*}b$ to $c\,b$, with $c>0$ the degree of $h^{n-m}$ on a general fibre;
so $b\in\Hdg^{2}(Y)$ lies in $\mathrm{Ab}^{2}(Y)$, or is algebraic, as soon
as $f^{*}b$ does, or is.

A birational morphism is surjective with general fibre a point, and two
birational smooth projective varieties are dominated through birational
morphisms by a resolution of the closure of the graph \cite{Hir64}; so both
properties are birational invariants. The morphism
$\widetilde\phi\colon\widetilde X\to R$ of the proof of
\Cref{prop:f3primesmall}, from a variety birational to $X$ onto its maximal
rationally connected quotient, has rationally connected general fibre,
which gives the statement on $R$. For $n\le5$ the degrees $0$, $2$, $2n-2$
and $2n$ are settled by \Cref{prop:f3prime}(iv), and for $n=5$ cup product
with $h$ is an isomorphism $H^{4}(X,\QQ)\to H^{6}(X,\QQ)$ of Hodge structures
up to a Tate twist, so that $\Hdg^{3}(X)=h\cup\Hdg^{2}(X)$; members of
$\mathcal{A}$ satisfy \textup{(F3$'$)} by \Cref{prop:f3prime}(iii). An
abelian fourfold $B$ satisfies the Hodge conjecture \cite{Mar25a,Mar25b}. A
dominant rational map onto $B$ with rationally connected general fibre
becomes, after its indeterminacy is resolved, a morphism whose general fibre
is smooth and rationally connected, as in the proof of
\Cref{prop:f3primesmall}; so the fourfolds and fivefolds of the last
assertion have every Hodge class of degree four algebraic, and the other
degrees are settled as before.
\end{proof}

\Cref{prop:f3primefibration} reaches varieties that \Cref{prop:aclosure}
does not: conic bundles with singular fibres over members of $\mathcal{A}$ of
dimension three or four, and fourfolds and fivefolds birational to members
of $\mathcal{A}$ without being known to lie in $\mathcal{A}$. A projective
bundle over a member of $\mathcal{A}$ already lies in $\mathcal{A}$ by
\Cref{prop:aclosure}(i), and so, by \Cref{prop:aclosure}(v) and the
Leray--Hirsch theorem, does the source of a smooth morphism onto a member of
$\mathcal{A}$ on which algebraic classes restrict to a basis of the
cohomology of a fibre, such as a Severi--Brauer bundle.
\Cref{prop:f3primefibration} is the relative form of the argument of Bloch
and Srinivas, and we claim no priority for it.

\begin{remark}[Zero-cycles alone go no further]\label{rem:f3primesharp}
Hypotheses on zero-cycles cannot carry \textup{(F3$'$)} further without the
fourfolds. If $\mathrm{CH}_{0}(X)_{\QQ}$ is supported on a closed subset $Z$,
with resolution $j\colon Z'\to X$, the proof of \Cref{prop:f3primefibration},
with the closure in $X\times Z$ of a zero-cycle on $Z_{K}$ rationally
equivalent to the generic point in place of $\Phi^{t}$, gives a divisor with
resolution $i\colon D'\to X$ and cycles $\Phi$ on $Z'\times X$ and $\Psi$ on
$X\times D'$ with $a=\Phi_{*}j^{*}a+i_{*}\Psi_{*}a$ for every
$a\in H^{k}(X,\QQ)$, where $\Psi_{*}a$ has degree $k-2$; this is the
transpose of the decomposition of Bloch and Srinivas \cite{BS83}. When
$Z\ne X$ and $a$ has type $(n,0)$, $n=\dim X$, the first term vanishes, as
$\dim Z'<n$, and the second has no component of type $(n,0)$; so
$H^{n,0}(X)=0$. A uniruled fivefold has its zero-cycles
supported on an ample divisor, and for a fourfold $Y$ and
$b\in\Hdg^{2}(Y)$ the class $\mathrm{pr}^{*}b$ is a Hodge class on
$Y\times\PP^{1}$ with $b=\mathrm{pr}_{*}(\mathrm{pr}^{*}b\cup\eta)$, for
$\eta$ the class of $Y\times\{0\}$. So \textup{(F3$'$)}, and likewise the
Hodge conjecture, holds for every uniruled fivefold if and only if it holds
in degree four for every fourfold.

It also holds for every rational sixfold if and only if it holds for every
fourfold. A rational sixfold has its zero-cycles supported on a point, so
the proof of \Cref{prop:f3primesmall} settles every degree but six; and it
is the image under a birational morphism of a variety obtained from
$\PP^{6}$ by blow-ups along smooth centres \cite{Hir64}, on which, by the
blow-up formula \cite[\S7.3.3]{Voi02}, the Hodge classes of degree six are
images under algebraic correspondences of Hodge classes of degree at most
four on the centres, of dimension at most four, and lie in
$\mathrm{Ab}^{3}$, or are algebraic, when the fourfolds satisfy
\textup{(F3$'$)}, or the conjecture. Conversely, a fourfold $Y$
is birational to a hypersurface in $\PP^{5}$; place it in a hyperplane of
$\PP^{6}$, and let $\widetilde Y$ be its smooth strict transform under an
embedded resolution $M\to\PP^{6}$ \cite{Hir64}. On the rational sixfold
$X=\mathrm{Bl}_{\widetilde Y}M$, with exceptional divisor $j\colon E\to X$
and projection $\rho\colon E\to\widetilde Y$, the normal bundle of $E$ is
$\cO_{E}(-1)$, so $a=-\rho_{*}j^{*}j_{*}\rho^{*}a$ for
$a\in H^{4}(\widetilde Y,\QQ)$; this recovers $\Hdg^{2}(\widetilde Y)$ from
$\Hdg^{3}(X)$ through algebraic correspondences, and
\Cref{prop:f3primefibration} passes from $\widetilde Y$ to $Y$.

Among smooth hypersurfaces, ordered by dimension and degree, the first
family that these results do not cover is that of the sextic fourfolds:
those of dimension at most three are covered by \Cref{prop:f3prime}(iv),
the fourfolds of degree at most five are Fano, hence rationally connected
\cite{Cam92,KMM92}, and covered by \Cref{prop:f3primesmall}, while a sextic
fourfold has $h^{4,0}=1$, so that its zero-cycles are not supported on a
proper closed subset and it is its own maximal rationally connected
quotient; the Fermat sextic lies in $\mathcal{A}$ by
\Cref{prop:aclosure}(iii). For a smooth sextic fourfold $Y$ in a
hyperplane of $\PP^{6}$ no resolution is needed, and the rational sixfold
$\mathrm{Bl}_{Y}\PP^{6}$, with $h^{6,0}=0$, $h^{5,1}=1$, $h^{4,2}=426$ and
$h^{3,3}=1753$ by the blow-up formula \cite[\S7.3.3]{Voi02} and the Hodge
numbers $1$, $426$, $1752$ of $Y$, satisfies \textup{(F3$'$)} if and only if
$Y$ does. For the Hodge
conjecture these equivalences are folklore, and we claim no priority for
them.
\end{remark}

Surjective morphisms carry $\mathcal{A}$ to itself (\Cref{prop:f3prime}(iii)).
In degree four the morphism can be replaced by any dominant rational map,
which reaches some of these sextic fourfolds.

\begin{proposition}[Degree four along dominant rational maps]
\label{prop:f3primedominant}
Let $\varphi\colon Y\dashrightarrow X$ be a dominant rational map of smooth
complex projective varieties.
\begin{enumerate}[label=\textup{(\roman*)},leftmargin=2.6em,itemsep=2pt]
\item If $\Hdg^{2}(Y)=\mathrm{Ab}^{2}(Y)$, then
  $\Hdg^{2}(X)=\mathrm{Ab}^{2}(X)$; if every Hodge class of degree four on
  $Y$ is algebraic, so is every one on $X$.
\item \textup{(F3$'$)} holds on every fourfold and every fivefold that
  admits a dominant rational map from a variety $Y$ with
  $\Hdg^{2}(Y)=\mathrm{Ab}^{2}(Y)$, for instance from a member of
  $\mathcal{A}$. The Hodge conjecture holds on every fourfold that admits a
  dominant rational map from an abelian fourfold, such as a smooth
  projective model of the quotient of an abelian fourfold by a finite group,
  and on every fivefold that admits one from the product of an abelian
  fourfold and a projective space.
\end{enumerate}
\end{proposition}

\begin{proof}
(i) Let $\Gamma\subset Y\times X$ be the closure of the graph of $\varphi$
and $\widetilde Y\to\Gamma$ a resolution of singularities \cite{Hir64}. The
projections give a birational morphism $\widetilde Y\to Y$ and a morphism
$g\colon\widetilde Y\to X$, which is surjective, being proper and dominant.
Both properties are birational invariants (\Cref{prop:f3primefibration}),
so they hold for $\widetilde Y$. Let $x\in\Hdg^{2}(X)$, let $h$ be ample on
$\widetilde Y$ and $r=\dim\widetilde Y-\dim X$, and let
$\Psi=[\Gamma_{g}]\cdot\mathrm{pr}_{\widetilde Y}^{*}h^{r}$, so that
$\Psi_{*}g^{*}x=c\,x$ with $c>0$, as in the proof of
\Cref{prop:f3prime}(iii). The class $g^{*}x$ is a Hodge class of degree four
on $\widetilde Y$, hence a sum of classes $\Gamma_{j*}\gamma_{j}$ as in
\Cref{def:f3prime}, or algebraic. So $x=c^{-1}\sum_{j}(\Psi\circ
\Gamma_{j})_{*}\gamma_{j}$ lies in $\mathrm{Ab}^{2}(X)$, or $x$ is
algebraic, since composites of algebraic correspondences are algebraic
\cite[Chapter~16]{Ful98}.

(ii) On a fourfold or a fivefold the degrees other than four and six are
settled by \Cref{prop:f3prime}(iv), and the same argument, the Lefschetz
theorem on $(1,1)$ classes and hard Lefschetz, makes the Hodge classes of
those degrees algebraic; on a fivefold cup product with an ample class is an
isomorphism from $H^{4}$ onto $H^{6}$ of Hodge structures up to a Tate twist,
so degree six follows from degree four. Members of $\mathcal{A}$ satisfy
\textup{(F3$'$)} (\Cref{prop:f3prime}(iii)). An abelian fourfold $B$
satisfies the Hodge conjecture \cite{Mar25a,Mar25b}, and by the K\"unneth
formula a Hodge class of degree four on $B\times\PP^{k}$ is a sum of
pull-backs of Hodge classes of degrees four, two and zero on $B$ multiplied
by powers of the hyperplane class, so it is algebraic. A quotient $B/G$ by a
finite group is dominated by $B$, and so is every smooth projective model of
it.
\end{proof}

Lefschetz theory carries the two statements from a variety to its ample
complete intersections, outside the middle degree always and in the middle
degree for a very general member.

\begin{proposition}[Ample complete intersections]
\label{prop:f3primeample}
Let $X$ be a smooth complex projective variety of dimension $N$, $L$ an
ample line bundle on $X$, and $Y=Y_{c}\subset Y_{c-1}\subset\dots\subset
Y_{0}=X$ with each $Y_{j}$ a smooth divisor of $Y_{j-1}$ in the linear
system of $L^{m_{j}}|_{Y_{j-1}}$, $m_{j}\ge1$; put $r=\dim Y=N-c\ge1$ and
let $i\colon Y\to X$ be the inclusion.
\begin{enumerate}[label=\textup{(\roman*)},leftmargin=2.6em,itemsep=2pt]
\item If $2p\ne r$, every Hodge class of degree $2p$ on $Y$ is $i^{*}\alpha$
  for a Hodge class $\alpha$ of degree $2p$ on $X$. So if \textup{(F3$'$)}, or
  the Hodge conjecture, holds for $X$ in degree $2p$, it holds for $Y$ in
  degree $2p$; if $r$ is odd and it holds for $X$ in every degree, it holds
  for $Y$ in every degree.
\item Let $r$ be even, $c=1$ and $L^{m_{1}}$ very ample, and let
  $V\subset H^{r}(Y,\QQ)$ be the vanishing cohomology, the kernel of the
  Gysin map $i_{*}$. If $V$ is not of type $(r/2,r/2)$, then for a very
  general $Y$ in the linear system every Hodge class of degree $r$ on $Y$ is
  $i^{*}\alpha$ for a Hodge class $\alpha$ on $X$, and the conclusion of (i)
  holds in degree $r$ as well.
\end{enumerate}
So the Hodge conjecture holds for every smooth ample divisor of odd dimension
in $B\times\PP^{k}$, for $B$ an abelian variety of dimension at most five, and
for the very general member of every very ample linear system on an abelian
fivefold; and \textup{(F3$'$)} holds for every smooth ample divisor of odd
dimension in a member of $\mathcal{A}$.
\end{proposition}

\begin{proof}
(i) By the Lefschetz hyperplane theorem \cite[Chapter~1]{Voi03}, applied
to each $Y_{j}\subset Y_{j-1}$, the restriction $i^{*}\colon H^{k}(X,\QQ)\to
H^{k}(Y,\QQ)$ is an isomorphism for $k<r$ and injective for $k=r$. Let
$2p<r$. Then $i^{*}$ is an isomorphism of Hodge structures in degree $2p$,
so every Hodge class on $Y$ is $i^{*}$ of a Hodge class on $X$. Let $2p>r$.
The Gysin map $i_{*}\colon H^{2p}(Y,\QQ)\to H^{2p+2c}(X,\QQ)$ is Poincar\'e
dual to $i^{*}$ in degree $2r-2p<r$, hence an isomorphism, and
$i_{*}i^{*}$ is cup product with $[Y]=m_{1}\cdots m_{c}\,c_{1}(L)^{c}$. By the
hard Lefschetz theorem \cite[Chapter~6]{Voi02}, cup product with
$c_{1}(L)^{c}$ maps $H^{2p}(X,\QQ)$ onto $H^{2p+2c}(X,\QQ)$ when $2p\ge N-c$,
so $i^{*}$ is onto in degree $2p$. A surjective morphism of polarisable
Hodge structures has a section that is a morphism of Hodge structures, so
again every Hodge class on $Y$ is $i^{*}$ of a Hodge class on $X$. The
restriction $i^{*}$ is induced by the transpose of the graph of $i$, so it
carries $\mathrm{Ab}^{p}(X)$ into $\mathrm{Ab}^{p}(Y)$ and algebraic classes
to algebraic classes, which gives the rest of (i).

(ii) The projection formula shows that $V$ is the orthogonal complement of
$i^{*}H^{r}(X,\QQ)$ for the intersection form, and
$H^{r}(Y,\QQ)=i^{*}H^{r}(X,\QQ)\oplus V$ as Hodge structures
\cite[Chapter~2]{Voi03}; the Hodge classes in the first summand are the
$i^{*}\alpha$ with $\alpha$ a Hodge class on $X$. Over the open set $U$ of
the linear system parametrising smooth members, $V$ is a local system, and
its monodromy group $\Gamma$ acts irreducibly \cite[Chapter~3]{Voi03}. For
each class $v$ of the lattice of $V$ at a point of the universal cover of
$U$, the set where the flat continuation of $v$ is of Hodge type is a closed
analytic subset, either everything or nowhere dense; so at a very general
point the Hodge classes of $V$ are those whose continuation is of Hodge type
everywhere. They form a subspace $W$ stable under $\Gamma$, hence $W=0$ or
$W=V$, and $W=V$ would make $V$ of type $(r/2,r/2)$. So $W=0$.

For the consequences, $B\times\PP^{k}$ satisfies the Hodge conjecture by the
K\"unneth formula, since $B$ does \cite{Mar25a,Mar25b,MZ99}. For a very ample
divisor $Y$ of an abelian fivefold $B$, $r=4$, and $V$ is not of type $(2,2)$:
$h^{4,0}(Y)=h^{0}(K_{Y})=h^{0}(Y,\cO_{Y}(Y))=h^{0}(B,\cO_{B}(Y))+4$, by the
sequence $0\to\cO_{B}\to\cO_{B}(Y)\to\cO_{Y}(Y)\to0$ and $h^{1}(\cO_{B}(Y))=0$,
exceeds $h^{4,0}(B)=5$ because $h^{0}(\cO_{B}(Y))\ge2$, so $V$ has a nonzero
part of type $(4,0)$. Members of $\mathcal{A}$ satisfy \textup{(F3$'$)} by
\Cref{prop:f3prime}(iii).
\end{proof}

For the Hodge conjecture these transfers are classical, and we claim no
priority. What they add to \textup{(F3$'$)} is a large class outside
$\mathcal{A}$: the vanishing cohomology of an ample divisor of odd dimension
is not reached from abelian varieties in general, yet it carries no Hodge
class in even degree, since its only degree is the odd middle one.

A hypersurface of Delsarte type is defined by as many monomials as
variables. Shioda observed that such a surface is dominated by a Fermat
surface \cite{Shi86}; the same monomial substitution works in every
dimension, and with \Cref{prop:f3primedominant} it gives the following.

\begin{proposition}[Delsarte fourfolds]\label{prop:delsarte}
Let $A=(a_{ij})_{0\le i,j\le r+1}$ be an invertible matrix of non-negative
integers whose rows all sum to $m\ge2$, let $F_{A}=\sum_{i}\prod_{j}
x_{j}^{a_{ij}}$, suppose that $X_{A}=\{F_{A}=0\}\subset\PP^{r+1}$ is smooth,
and let $d$ be a positive integer with $dA^{-1}$ integral, such as
$|\det A|$.
\begin{enumerate}[label=\textup{(\roman*)},leftmargin=2.6em,itemsep=2pt]
\item There is a dominant rational map from the Fermat hypersurface
  $X^{r}_{d}$ onto $X_{A}$, given by monomials.
\item $\Hdg^{2}(X_{A})=\mathrm{Ab}^{2}(X_{A})$. For $r=4$, $X_{A}$
  satisfies \textup{(F3$'$)}, and it satisfies the Hodge conjecture if
  $X^{4}_{d}$ does.
\item If $m\ge3$ and $F_{A}$ is a sum, in disjoint sets of variables, of
  terms $x^{m}$, chains $x_{1}^{m-1}x_{2}+\dots+x_{k-1}^{m-1}x_{k}+x_{k}^{m}$
  and loops $x_{1}^{m-1}x_{2}+\dots+x_{k-1}^{m-1}x_{k}+x_{k}^{m-1}x_{1}$
  with $k\ge2$, then $A$ is invertible and $X_{A}$ is smooth. For $r=4$ and
  $m=6$ these polynomials have $29$ shapes up to a renumbering of the
  variables, the Fermat sextic being one; each defines a smooth sextic
  fourfold, with $h^{4,0}=1$, $h^{3,1}=426$ and primitive $h^{2,2}$ equal
  to $1751$. The least $d$ is $6$ for the Fermat sextic, $24$ or $30$ for
  six others, those whose terms other than $x^{m}$ are all loops of length
  two or all chains of length two, at least $120$ for the remaining
  twenty-two, and $5^{6}-1=15624$ for the loop
  $x_{0}^{5}x_{1}+x_{1}^{5}x_{2}+\dots+x_{5}^{5}x_{0}$.
\end{enumerate}
\end{proposition}

\begin{proof}
(i) The matrix $dA^{-1}$ is integral for $d=|\det A|$, by the formula
for the inverse through the adjugate. Let $B=dA^{-1}$. From
$A\mathbf{1}=m\mathbf{1}$, with $\mathbf{1}$ the vector of ones,
$B\mathbf{1}=(d/m)\mathbf{1}$, so the rows of the integral matrix $B$ all
sum to the integer $d/m$, and $AB=d\cdot I$. Choose $c\ge0$
with every $b_{jk}+c\ge0$ and let $\psi(y)_{j}=\prod_{k}y_{k}^{b_{jk}+c}$;
the coordinates of $\psi$ are monomials of the same degree, so $\psi$ is a
rational map $\PP^{r+1}\dashrightarrow\PP^{r+1}$, and
\[
  F_{A}(\psi(y))=\sum_{i}\prod_{k}y_{k}^{\sum_{j}a_{ij}(b_{jk}+c)}
  =(y_{0}\cdots y_{r+1})^{cm}\sum_{i}y_{i}^{d}.
\]
On the torus $T$ of points with all coordinates nonzero, $\psi$ is induced
by the homomorphism $y\mapsto(\prod_{k}y_{k}^{b_{jk}})_{j}$ of
$\mathbb{G}_{m}^{r+2}$, which is surjective because $\det B\ne0$. So every
$z\in X_{A}\cap T$ is $\psi(y)$ for some $y\in T$, and the identity puts $y$
on $X^{r}_{d}$. A smooth hypersurface of degree at least two and positive
dimension is irreducible, so it is not contained in a coordinate hyperplane,
and $X_{A}\cap T$ and $X^{r}_{d}\cap T$ are dense. Hence $\psi$ restricts to
a dominant rational map $X^{r}_{d}\dashrightarrow X_{A}$.

(ii) $X^{r}_{d}$ lies in $\mathcal{A}$ by \Cref{prop:aclosure}(iii), so (i)
and \Cref{prop:f3primedominant} give both assertions, the degrees other than
four being settled as in the proof of \Cref{prop:f3primedominant}(ii).

(iii) The exponent matrix is block diagonal, with blocks $m$,
$(m-1)I+N+E$ for a chain, $N$ the nilpotent shift and $E$ the matrix whose
only nonzero entry is a $1$ in the last diagonal place, and $(m-1)I+P$ for a
loop, $P$ the cyclic shift; their determinants are $m$, $m(m-1)^{k-1}$ and
$(m-1)^{k}-(-1)^{k}$, which are nonzero. The partial derivatives of a sum in
disjoint sets of variables are those of the summands, so it suffices that
the partials of each summand vanish together only at $0$. For $x^{m}$ this
is clear. For a chain, $\partial F/\partial x_{1}=(m-1)x_{1}^{m-2}x_{2}$
vanishes only if $x_{1}=0$ or $x_{2}=0$, and if $x_{2}=0$ then
$\partial F/\partial x_{2}=x_{1}^{m-1}$; so $x_{1}=0$, the same argument
applies to $x_{2},\dots,x_{k-1}$, and then $\partial F/\partial
x_{k}=mx_{k}^{m-1}$ gives $x_{k}=0$. For a loop, $\partial F/\partial
x_{j}=(m-1)x_{j}^{m-2}x_{j+1}+x_{j-1}^{m-1}$, indices modulo $k$. If one
coordinate vanishes, so does the one before it, hence all of them; if none
vanishes, the product of the relations $x_{j-1}^{m-1}=-(m-1)x_{j}^{m-2}
x_{j+1}$ gives $\prod_{j}x_{j}^{m-1}=(1-m)^{k}\prod_{j}x_{j}^{m-1}$, which
is impossible because $|1-m|\ge2$. The number of shapes, drawn in \Cref{fig:delsarte}, is the
coefficient of $t^{6}$ in $(1-t)^{-1}\prod_{k\ge2}(1-t^{k})^{-2}$, which is $29$, and
the least $d$ is the least common multiple of the orders of the cyclic groups
attached to the terms in the proof of \Cref{cor:delsarteall}(iii): $6$ for
$x^{6}$, $6\cdot5^{k-1}$ for a chain and $5^{k}-(-1)^{k}$ for a loop of length
$k$; so $24$ or $30$ when the terms other than $x^{6}$ are all loops or all
chains of length two, and at least $\operatorname{lcm}(24,30)=120$, $5^{3}+1$
or $6\cdot5^{2}$ when both kinds of length two, or a loop or chain of length at
least three, occur. The
Hodge numbers are those of every smooth sextic fourfold, read off from the
Jacobian ring \cite[Chapter~6]{Voi03}.
\end{proof}

\begin{figure}[tbp]
\centering
  \includegraphics[width=\textwidth]{fig_delsarte}
\caption{The $29$ smooth sextic fourfolds of \Cref{prop:delsarte}(iii), up
to a renumbering of the variables. Each is drawn on its six variables: an
arrow from $x_{i}$ to $x_{j}$ is a term $x_{i}^{5}x_{j}$, and a ring is a
term $x^{6}$, so a chain is a path of arrows ending at a ring and a loop is
a cycle of arrows. Under each is the least $d$ with $dA^{-1}$ integral, the
degree of the Fermat fourfold that covers it. All of them satisfy
\textup{(F3$'$)}.}
\label{fig:delsarte}
\end{figure}

A sextic fourfold has $h^{4,0}=1$, so no hypothesis on zero-cycles or on
rationally connected fibrations reaches it (\Cref{rem:f3primesharp}), and
among the smooth sextic fourfolds only the Fermat one is covered by
\Cref{prop:aclosure}. The Delsarte sextics are finitely many
up to isomorphism, since the torus rescales the coefficients of a Delsarte
polynomial to $1$, so they are isolated points in the $426$-dimensional
moduli of smooth sextic fourfolds: \Cref{prop:delsarte} adds points, not a
family. By \Cref{prop:delsarte}(ii) the Hodge conjecture for a Delsarte
fourfold follows from that for one Fermat fourfold. For six of the sextic
shapes that Fermat fourfold has degree $24$ or $30$, inside the range claimed
in the preprint \cite{MMMV26}, so if it is correct the Hodge conjecture
itself holds on these six; for the others the degree ranges from $120$ to
$18750$, and from $120$ to $3750$ with the lattice degree of
\Cref{cor:delsarteall}, beyond every known or claimed range
\cite{Shi79b,MMMV26,Jum26}.
The monomial cover is Shioda's, and we claim no priority for it.

\Cref{prop:delsarte} uses the covering map only in degree four, where
algebraicity modulo abelian varieties descends along every dominant rational
map. In higher degrees a birational modification can add classes, so the
boundary of a compactification has to be controlled. For a hypersurface of a
torus defined by affinely independent monomials the boundary is again of the
same kind, and the whole cohomology can be traced back to Fermat varieties
of a single degree. We use toric varieties as in \cite{CLS11}: a torus $T$
with character lattice $M$ and cocharacter lattice $N$, a fan $\Sigma$ in
$N_{\RR}$ and its toric variety $Y_{\Sigma}$, and for a cone
$\tau\in\Sigma$ the orbit $O(\tau)$, whose character lattice is
$\tau^{\perp}\cap M$, and its closure $V(\tau)$, the toric variety of the
fan formed by the images of the cones containing $\tau$ in
$N_{\RR}/\RR\tau$ \cite[Chapter~3]{CLS11}; for a ray $\rho$ we write
$D_{\rho}=V(\rho)$. In what follows varieties need not be connected: a
smooth projective variety lies in $\mathcal{A}$ when each of its connected
components does, and the closure properties of \Cref{prop:f3prime}(iii) and
\Cref{prop:aclosure} hold component by component.

\begin{definition}[Simplex type]\label{def:simplextype}
Let $T$ be a torus of dimension $n$ with character lattice $M$. A Laurent
polynomial $f=\sum_{i=0}^{k}c_{i}\chi^{u_{i}}$ on $T$ is of \emph{simplex
type} if every $c_{i}$ is nonzero and $u_{0},\dots,u_{k}\in M$ are affinely
independent. Let $L'\subseteq M$ be the lattice spanned by the
$u_{i}-u_{0}$, $L$ its saturation, the set of $x\in M$ with a nonzero
multiple in $L'$, and $e(f)$ the exponent of the finite group $L/L'$; put
$V_{f}=\{f=0\}\subseteq T$. A fan $\Sigma$ in $N_{\RR}$ is \emph{adapted} to
$f$ if for every cone $\sigma\in\Sigma$ there is an index $i_{\sigma}$ with
$\langle u_{i_{\sigma}},v\rangle\le\langle u_{j},v\rangle$ for all $j$ and
all $v\in\sigma$, that is, if every cone of $\Sigma$ lies in a cone of the
normal fan of the simplex spanned by the $u_{i}$. For a cone $\tau$ of an
adapted fan let $I_{\tau}$ be the set of indices $i$ with
$\langle u_{i},v\rangle=\min_{j}\langle u_{j},v\rangle$ for every
$v\in\tau$, fix $i_{\tau}\in I_{\tau}$, and let
$f_{\tau}=\sum_{i\in I_{\tau}}c_{i}\chi^{u_{i}-u_{i_{\tau}}}$, a Laurent
polynomial of simplex type on the torus $O(\tau)$, since its exponents lie
in $\tau^{\perp}\cap M$.
\end{definition}

The compactification is Khovanskii's \cite{Kho77}; for a simplex the proof
is short, and we give it.

\begin{lemma}[The compactification]\label{lem:simplexchart}
Let $f$ be of simplex type on $T$, let $\Sigma$ be a smooth fan adapted to
$f$, and let $Z$ be the closure of $V_{f}$ in $Y_{\Sigma}$.
\begin{enumerate}[label=\textup{(\roman*)},leftmargin=2.6em,itemsep=2pt]
\item $V_{f}$ is empty if $k=0$. If $k\ge1$, it is a finite \'etale cover,
  of degree $[L:L']$, of the product of a torus of dimension $n-k$ and the
  complement of $k+1$ hyperplanes in general position in $\PP^{k-1}$; in
  particular it is smooth.
\item $Z$ is smooth, and for every $\tau\in\Sigma$, $Z\cap O(\tau)=V_{f_{\tau}}$
  and $Z$ meets $V(\tau)$ transversally: $Z\cap V(\tau)$ is empty or smooth of
  codimension $\dim\tau$ in $Z$, and it is the closure of $V_{f_{\tau}}$ in
  $V(\tau)$. So the divisors $Z\cap D_{\rho}$ of $Z$ are smooth and cross
  normally, and $Z\cap D_{\rho}$ is empty when $I_{\rho}$ has one element.
\item The fan of $V(\tau)$ is smooth and adapted to $f_{\tau}$.
\end{enumerate}
\end{lemma}

\begin{proof}
(i) If $k=0$, $f$ is a monomial and has no zero on $T$. Let $k\ge1$, choose a
complement $C$ of the saturated lattice $L$ in $M$ and a basis
$b_{k+1},\dots,b_{n}$ of $C$. The characters $u_{1}-u_{0},\dots,u_{k}-u_{0},
b_{k+1},\dots,b_{n}$ span a sublattice of index $[L:L']$ in $M$, so the map
$T\to(\CC^{\times})^{n}$ that they define is an isogeny of that degree. The
function $\chi^{-u_{0}}f$ is the pull-back of $c_{0}+c_{1}z_{1}+\dots+c_{k}
z_{k}$, so $V_{f}$ is the full preimage of the product of
$(\CC^{\times})^{n-k}$ and the set $H$ of points of $(\CC^{\times})^{k}$ with
$c_{0}+\sum_{i}c_{i}z_{i}=0$. The projective closure of this affine hyperplane is a hyperplane
$\PP^{k-1}$ of $\PP^{k}$, with homogeneous coordinates $w_{0},\dots,w_{k}$,
on which the $k+1$ linear forms $w_{i}$ satisfy the single relation
$\sum_{i}c_{i}w_{i}=0$, with every $c_{i}\ne0$; so any $k$ of them are
independent, the hyperplanes $\{w_{i}=0\}$ of $\PP^{k-1}$ are in general
position, and $H$ is their complement.

(ii) Let $\sigma\in\Sigma$, spanned by part $v_{1},\dots,v_{s}$ of a basis
$v_{1},\dots,v_{n}$ of $N$, let $w_{1},\dots,w_{n}$ be the dual basis of $M$,
and let $t_{j}=\chi^{w_{j}}$, the coordinates of
$U_{\sigma}\cong\CC^{s}\times(\CC^{\times})^{n-s}$. Put $i=i_{\sigma}$. Every
$u_{j}-u_{i}$ lies in the dual cone of $\sigma$, so
$g=\chi^{-u_{i}}f=\sum_{j}c_{j}\chi^{u_{j}-u_{i}}$ is a regular function on
$U_{\sigma}$, and $Z\cap U_{\sigma}\subseteq\{g=0\}$. A face $\tau$ of
$\sigma$ is spanned by the $v_{j}$ with $j$ in a subset $J$, and $O(\tau)$ is
the locus where exactly the $t_{j}$ with $j\in J$ vanish. A character
$\chi^{x}$ regular on $U_{\sigma}$ restricts to $O(\tau)$ as $\chi^{x}$ if
$x\in\tau^{\perp}$ and as $0$ otherwise \cite[Chapter~3]{CLS11}; since
$\langle u_{j}-u_{i},v\rangle\ge0$ on $\sigma$, $u_{j}-u_{i}\in\tau^{\perp}$
exactly when $j\in I_{\tau}$. So $h=g|_{O(\tau)}$ is
$\chi^{u_{i_{\tau}}-u_{i}}f_{\tau}$, a unit times $f_{\tau}$. Let $p\in
O(\tau)$ with $h(p)=0$ and put $a_{j}=c_{j}\chi^{u_{j}-u_{i}}(p)$ for $j\in
I_{\tau}$. The derivative of $h$ at $p$ along the one-parameter subgroup of
$O(\tau)$ given by $y$ in its cocharacter lattice is $\sum_{j}a_{j}\langle
u_{j}-u_{i},y\rangle$. If all of these vanished, then, with
$\sum_{j}a_{j}=h(p)=0$, the vector $\sum_{j}a_{j}(1,u_{j}-u_{i})$ of
$\CC\oplus(\tau^{\perp}\cap M)\otimes\CC$ would be zero, as the cocharacter
lattice of $O(\tau)$ is dual to $\tau^{\perp}\cap M$; affine independence
would give every $a_{j}=0$, which is impossible at a point of the torus
$O(\tau)$. So $dh(p)\ne0$. Hence $\{g=0\}$ is smooth at $p$ and transversal
to $O(\tau)$ and to every $\{t_{j}=0,\ j\in J'\}$ with $J'\subseteq J$, whose
tangent space at $p$ contains that of $O(\tau)$. No component of $\{g=0\}$
lies in a divisor $\{t_{j}=0\}$, since $g$ restricted to the orbit of the
ray of $v_{j}$ is a unit times a nonzero Laurent polynomial; so every
component meets $T$, where $\{g=0\}$ is $V_{f}$, and
$\{g=0\}=Z\cap U_{\sigma}$. This shows that $Z$ is smooth, that
$Z\cap O(\tau)=V_{f_{\tau}}$ and that $Z\cap V(\tau)$ is smooth of
codimension $\dim\tau$. As $Z\cap O(\tau')$ has codimension at least
$\dim\tau'$ in $Z$ for every cone $\tau'\supsetneq\tau$, no component of
$Z\cap V(\tau)$ lies in the smaller orbits, and $Z\cap V(\tau)$ is the
closure of $V_{f_{\tau}}$. If $I_{\rho}$ has one element, $f_{\rho}$ is a
nonzero constant.

(iii) The fan of $V(\tau)$ is smooth because $\Sigma$ is. Let $\sigma\supseteq
\tau$. The index $i_{\sigma}$ attains the minimum on $\sigma$, hence on
$\tau$, so $i_{\sigma}\in I_{\tau}$, and for $j\in I_{\tau}$ and $v\in\sigma$,
$\langle u_{i_{\sigma}}-u_{i_{\tau}},v\rangle\le\langle
u_{j}-u_{i_{\tau}},v\rangle$. Both sides depend only on the image of $v$ in
$N_{\RR}/\RR\tau$, because the exponents of $f_{\tau}$ lie in $\tau^{\perp}$;
so the image of $\sigma$ satisfies the condition of adaptedness for
$f_{\tau}$.
\end{proof}

\begin{lemma}[Cohomology across a toric modification]\label{lem:toricmod}
Let $f$ be of simplex type on $T$, let $\Sigma_{1}$ and $\Sigma_{0}$ be
smooth projective fans adapted to $f$, with $\Sigma_{1}$ refining
$\Sigma_{0}$, let $\varphi\colon Y_{1}=Y_{\Sigma_{1}}\to Y_{0}=Y_{\Sigma_{0}}$
be the toric morphism, and let $Z_{1}$ and $Z_{0}$ be the closures of
$V_{f}$. For each ray $\rho$ of $\Sigma_{1}$ that is not a ray of
$\Sigma_{0}$ let $E_{\rho}=Z_{1}\cap D_{\rho}$, with inclusion $i_{\rho}$.
Then the kernel of $\varphi_{*}\colon H^{k}(Z_{1},\QQ)\to H^{k}(Z_{0},\QQ)$ is
contained in $\sum_{\rho}i_{\rho*}H^{k-2}(E_{\rho},\QQ)$, and
\[
  H^{k}(Z_{1},\QQ)=\varphi^{*}H^{k}(Z_{0},\QQ)
  +\sum_{\rho}i_{\rho*}H^{k-2}(E_{\rho},\QQ).
\]
\end{lemma}

\begin{proof}
Let $W\subseteq Y_{0}$ be the union of the open sets $U_{\tau}$ with
$\tau\in\Sigma_{0}\cap\Sigma_{1}$. A cone of $\Sigma_{1}$ contained in such a
$\tau$ is a face of it, since two cones of a fan meet in a common face; so
$\varphi^{-1}(U_{\tau})=U_{\tau}$, and $\varphi$ is an isomorphism over $W$. A
cone of $\Sigma_{1}$ whose rays are all rays of $\Sigma_{0}$ lies in a cone
$\tau_{0}$ of $\Sigma_{0}$, its rays are rays of $\tau_{0}$, and $\tau_{0}$
is simplicial, so it is a face of $\tau_{0}$ and lies in $\Sigma_{0}$. Hence
$Y_{1}\smallsetminus\varphi^{-1}(W)$ is the union of the orbits of the cones
with a ray not in $\Sigma_{0}$, that is, of the divisors $D_{\rho}$ of those
rays. So $U_{1}=Z_{1}\smallsetminus\bigcup_{\rho}E_{\rho}$ is the closure of
$V_{f}$ in $\varphi^{-1}(W)$, and $\varphi$ maps it isomorphically onto
$U_{0}=Z_{0}\cap W$, with $U_{1}=Z_{1}\cap\varphi^{-1}(U_{0})$. The morphism
$Z_{1}\to Z_{0}$ is birational on each component, so $\varphi_{*}\varphi^{*}
=\id$ by the projection formula, and every class is $\varphi^{*}\varphi_{*}a$
plus a class in the kernel of $\varphi_{*}$. Let $\varphi_{*}b=0$.
Push-forward along a proper map commutes with restriction to an open set and
its preimage, so $b|_{U_{1}}=(\varphi|_{U_{1}})^{*}\bigl((\varphi_{*}b)|_{U_{0}}
\bigr)=0$. By \Cref{lem:simplexchart}(ii), $\bigcup_{\rho}E_{\rho}$ is a
divisor with normal crossings and smooth components in the smooth projective
variety $Z_{1}$, and by Deligne \cite[\S3.2]{Del71} the weight spectral
sequence of its complement $U_{1}$ degenerates at $E_{2}$; its term
$E_{2}^{0,k}$, the cokernel of the Gysin map
$\bigoplus_{\rho}H^{k-2}(E_{\rho},\QQ)\to H^{k}(Z_{1},\QQ)$, is the image of
$H^{k}(Z_{1},\QQ)$ in $H^{k}(U_{1},\QQ)$. So $b$ lies in the image of the
Gysin map.
\end{proof}

\begin{theorem}[Hypersurfaces of simplex type]\label{thm:simplextype}
Let $f$ be a Laurent polynomial of simplex type on a torus $T$ of dimension
$n$, with $e=e(f)$, let $Y$ be a projective toric variety with torus $T$,
not necessarily smooth, and let $Z$ be the closure of $V_{f}$ in $Y$.
Suppose that $Z$ is smooth.
\begin{enumerate}[label=\textup{(\roman*)},leftmargin=2.6em,itemsep=2pt]
\item $Z\in\mathcal{A}$, so \textup{(F3$'$)} holds on $Z$ in every degree.
\item If the Hodge conjecture holds for the Fermat varieties $X^{j}_{e}$
  with $j\le\dim Z$, then it holds for $Z$.
\end{enumerate}
\end{theorem}

\begin{proof}
By induction on $n$. If $k=0$, $Z$ is empty; let $k\ge1$.

\emph{The Fermat form.} With $C$ and $b_{j}$ as in the proof of
\Cref{lem:simplexchart}(i), let $\widetilde M$ be the lattice with basis
$w_{i}=(u_{i}-u_{0})/e$ for $1\le i\le k$ and $w_{j}=b_{j}$ for $k<j\le n$.
It contains $M$, since $eL\subseteq L'$ gives $L\subseteq\frac1eL'$. Let
$\widetilde T$ be the torus with character lattice $\widetilde M$,
$\widetilde N\subseteq N$ its cocharacter lattice and
$\pi\colon\widetilde T\to T$ the isogeny dual to $M\subseteq\widetilde M$.
In the coordinates $z_{i}=\chi^{w_{i}}$ of $\widetilde T$,
$\pi^{*}(\chi^{-u_{0}}f)=c_{0}+\sum_{i\le k}c_{i}z_{i}^{e}$; call it
$\widetilde f$. It is of simplex type, $V_{\widetilde f}=\pi^{-1}(V_{f})$,
and after the $z_{i}$ are rescaled it is $c_{0}(1+\sum_{i\le k}z_{i}^{e})$.

\emph{The model.} Let $\Sigma_{0}$ be the fan in $\widetilde N_{\RR}$ of
$Y_{0}=\PP^{k}\times\PP^{n-k}$, on whose torus $z_{1},\dots,z_{k}$ are the
affine coordinates of the first factor and $z_{k+1},\dots,z_{n}$ those of the
second. The closure of $V_{\widetilde f}$ in $Y_{0}$ is, after the
rescaling, $Z_{0}=X^{k-1}_{e}\times\PP^{n-k}$, since the part of the Fermat
variety in the torus is dense in it; $Z_{0}\in\mathcal{A}$ by
\Cref{prop:aclosure}(iii) and \Cref{prop:f3prime}(iii), $\PP^{n-k}$ being
the Fermat variety $X^{n-k}_{1}$. The fan $\Sigma_{0}$ is adapted to
$\widetilde f$: $\min(0,e\langle w_{1},v\rangle,\dots,e\langle
w_{k},v\rangle)$ depends only on the image of $v$ in the fan of $\PP^{k}$,
the normal fan of the standard simplex, on whose cones it is linear.

\emph{The common refinement.} Let $\Sigma_{Y}$ be the fan of $Y$, the normal
fan of a polytope $P$, as $Y$ is projective \cite[Chapter~2]{CLS11}, and
$\Sigma_{0}$ that of a polytope $Q$. The face of $P+Q$ on which a vector $v$
attains its minimum is the sum of those of $P$ and of $Q$, so the normal fan
of $P+Q$ consists of the intersections of cones of $\Sigma_{Y}$ and
$\Sigma_{0}$; it has a smooth projective refinement $\Sigma_{1}$ with
respect to $\widetilde N$ \cite[Chapter~11]{CLS11}. It is adapted to
$\widetilde f$, because $\Sigma_{0}$ is. Let
$Y_{1}=Y_{\Sigma_{1},\widetilde N}$ and $Z_{1}$ the closure of
$V_{\widetilde f}$ in it, smooth by \Cref{lem:simplexchart}(ii). The
inclusion $\widetilde N\subseteq N$ maps each cone of $\Sigma_{1}$ into a cone
of $\Sigma_{Y}$ and defines a proper toric morphism $\psi\colon Y_{1}\to Y$
extending $\pi$ \cite[Chapter~3]{CLS11}; $\psi(Z_{1})$ is closed, contains
$\pi(V_{\widetilde f})=V_{f}$ and lies in $Z$, so $\psi$ maps $Z_{1}$ onto
$Z$. The identity of $\widetilde N$ defines the toric morphism
$\varphi\colon Y_{1}\to Y_{0}$, which maps $Z_{1}$ onto $Z_{0}$ and is an
isomorphism over $V_{\widetilde f}$.

\emph{The boundary.} By \Cref{lem:simplexchart}, each $E_{\rho}=Z_{1}\cap
D_{\rho}$ is the smooth closure of $V_{\widetilde f_{\rho}}$ in the smooth
projective toric variety $D_{\rho}$ of dimension $n-1$. The exponents of
$\widetilde f_{\rho}$ are the $e(w_{i}-w_{i_{\rho}})$ with $i\in I_{\rho}$,
where $w_{0}=0$. The vectors $w_{i}-w_{i_{\rho}}$ with $i\le k$,
$i\ne i_{\rho}$, together with $w_{k+1},\dots,w_{n}$, form a basis of
$\widetilde M$, so those with $i\in I_{\rho}$ span a saturated sublattice of
$\widetilde M$, which lies in $\rho^{\perp}\cap\widetilde M$ and is
saturated there. Hence $e(\widetilde f_{\rho})=e$ if $I_{\rho}$ has at least
two elements, and $E_{\rho}$ is empty otherwise. By induction every
$E_{\rho}$ lies in $\mathcal{A}$, and it satisfies the Hodge conjecture if
the $X^{j}_{e}$ with $j\le n-2$ do.

(i) By \Cref{lem:toricmod}, $H^{*}(Z_{1},\QQ)$ is spanned by the images of
$H^{*}(Z_{0},\QQ)$ and of the $H^{*}(E_{\rho},\QQ)$ under $\varphi^{*}$ and
the $i_{\rho*}$, which are induced by algebraic cycles, the transposed
graph of $\varphi$ and the graphs of the $i_{\rho}$. So
$Z_{1}\in\mathcal{A}$ by \Cref{prop:aclosure}(v), and $Z\in\mathcal{A}$ by
\Cref{prop:f3prime}(iii), as the image of $Z_{1}$ under $\psi$.

(ii) By the K\"unneth formula every Hodge class on $X^{k-1}_{e}\times
\PP^{n-k}$ is a sum of pull-backs of Hodge classes on $X^{k-1}_{e}$
multiplied by powers of the hyperplane class, so the conjecture holds on
$Z_{0}$, as $k-1\le n-1=\dim Z$; it holds on the $E_{\rho}$ by induction.
Let $a$ be a Hodge class on $Z_{1}$. Then $\varphi_{*}a$ is a Hodge class,
hence algebraic, and so is $\varphi^{*}\varphi_{*}a$. The class
$b=a-\varphi^{*}\varphi_{*}a$ is a Hodge class in the kernel of
$\varphi_{*}$, hence in the image of the Gysin map of
\Cref{lem:toricmod}, a morphism of polarisable Hodge structures after a
Tate twist; as in the proof of \Cref{prop:f3prime}(iii) it is the image of a
Hodge class, which is algebraic, so $b$ is algebraic. So the conjecture
holds on $Z_{1}$, and it passes to $Z$ along $\psi$: for a Hodge class $x$
on $Z$, $\psi^{*}x$ is algebraic, and $\Psi_{*}\psi^{*}x=c\,x$ with $c>0$ for
the algebraic correspondence $\Psi$ of the proof of
\Cref{prop:f3prime}(iii).
\end{proof}

\begin{corollary}[Delsarte hypersurfaces of every dimension]
\label{cor:delsarteall}
Let $X_{A}\subset\PP^{r+1}$, $r\ge1$, be a smooth Delsarte hypersurface of
degree $m$ as in \Cref{prop:delsarte}, let $M=\{x\in\ZZ^{r+2}:\sum_{j}x_{j}=0\}$,
and let $e_{A}$ be the exponent of $M/M'$, where $M'$ is spanned by the
differences $a_{i}-a_{0}$ of the rows of $A$.
\begin{enumerate}[label=\textup{(\roman*)},leftmargin=2.6em,itemsep=2pt]
\item $X_{A}\in\mathcal{A}$, so \textup{(F3$'$)} holds on $X_{A}$ in every
  degree, and the Hodge conjecture holds on $X_{A}$ if it holds on the Fermat
  varieties $X^{j}_{e_{A}}$ with $j\le r$. The number $e_{A}$ divides every
  $d$ of \Cref{prop:delsarte}, and $[M:M']=|\det A|/m$.
\item The same holds, with $j\le r+1$, for the cyclic cover
  $\{y^{k}=F_{A}(x)\}$ of $\PP^{r+1}$ in the weighted projective space
  $\PP(1,\dots,1,m/k)$, for every $k$ dividing $m$, with
  $M=\{x\in\ZZ^{r+3}:\sum_{j\le r+1}x_{j}+(m/k)x_{r+2}=0\}$ and $M'$ spanned
  by the differences of the exponents of $F_{A}-y^{k}$; it is smooth
  because $X_{A}$ is.
\item For the $29$ sextic shapes of \Cref{prop:delsarte}(iii), $e_{A}$ is the
  least $d$ there, except for the chain $x_{0}^{5}x_{1}+x_{1}^{5}x_{2}+\dots
  +x_{4}^{5}x_{5}+x_{5}^{6}$ and the loop, for which it is $3125$ and $2604$
  instead of $18750$ and $15624$. For the Klein quartic
  $x^{3}y+y^{3}z+z^{3}x$ it is $7$.
\end{enumerate}
\end{corollary}

\begin{proof}
(i) On the torus of $\PP^{r+1}$, with character lattice $M$, $X_{A}$ is the
closure of $V_{f}$ for $f=\chi^{-a_{0}}F_{A}$, as in the proof of
\Cref{prop:delsarte}(i). The rows of $A$ are linearly independent and lie on
the affine hyperplane $\sum_{j}x_{j}=m$, which does not contain $0$, so they
are affinely independent, $f$ is of simplex type, $L'=M'$ has rank $r+1$ and
$L=M$; so $e(f)=e_{A}$, and \Cref{thm:simplextype} gives the first two
assertions. If $dA^{-1}$ is integral and $x\in M$, then $dx=\lambda A$ with
$\lambda=dxA^{-1}$ integral and $\sum_{i}\lambda_{i}=d\sum_{j}x_{j}/m=0$,
as $A^{-1}\mathbf{1}=\mathbf{1}/m$; so $dx=\sum_{i}\lambda_{i}(a_{i}-a_{0})
\in M'$, and $e_{A}$ divides $d$. The lattice $M'$ is the intersection of $M$
with the lattice $R$ spanned by the rows, whose image under
$x\mapsto\sum_{j}x_{j}$ is $m\ZZ$; the exact sequence $0\to M/M'\to
\ZZ^{r+2}/R\to\ZZ/m\ZZ\to0$ gives $[M:M']=|\det A|/m$.

(ii) The exponents of $F_{A}-y^{k}$ are linearly independent and lie on the
hyperplane of weighted degree $m$, so the polynomial is of simplex type on
the torus of $\PP(1,\dots,1,m/k)$, a projective toric variety whose
character lattice is the given $M$. On the affine cone, the partial
derivatives vanish together with $y^{k}-F_{A}$ only where $y=0$ and $x$ is a
singular point of the cone over $X_{A}$, that is $x=0$; so the cover is
quasi-smooth, and it misses the only singular point $(0:\dots:0:1)$ of the
weighted projective space, where $y^{k}=F_{A}(0)$ would force $y=0$. So it is
smooth. It is irreducible, $F_{A}$ being irreducible, and not contained in a
coordinate hyperplane, so it is the closure of its part in the torus, and
\Cref{thm:simplextype} applies.

(iii) The least $d$ with $dA^{-1}$ integral is the exponent of
$\ZZ^{r+2}/R$, and the map $x\mapsto\sum_{j}x_{j}$ modulo $m$ is well defined
on it because the rows sum to $m$. For a polynomial as in
\Cref{prop:delsarte}(iii) the group $\ZZ^{r+2}/R$ is the direct sum of one
group for each term. The relations of a chain are $x_{j+1}=(1-m)x_{j}$ and
$mx_{k}=0$, those of a loop are $x_{j+1}=(1-m)x_{j}$, indices modulo $k$, and
that of $x^{m}$ is $mx=0$. So each summand is cyclic, generated by the class
of its first variable, of order $N=m(m-1)^{k-1}$, $(m-1)^{k}-(-1)^{k}$ and
$m$ respectively, and the sum map sends that generator to $1$, since
$1-m\equiv1$ modulo $m$; each $N$ is divisible by $m$. The least $d$ is the
least common multiple of these orders, which gives the values of
\Cref{prop:delsarte}(iii). The group $M/M'$ is the kernel of the sum map. With
one term only, the loop or the chain in six variables, the kernel is
$m\ZZ/N\ZZ$, of exponent $N/m$, which is $15624/6=2604$ and $18750/6=3125$;
for the Klein quartic, a loop with $m=4$ and $k=3$, it is $28/4=7$. With at
least two terms, of orders $N_{1},N_{2},\dots$, the kernel contains the
differences of the generators of the first and of the $b$-th summand, of order
$\operatorname{lcm}(N_{1},N_{b})$; so its exponent is the least common multiple
of all the $N_{b}$, which is the least $d$.
\end{proof}

For the Klein quartic the construction is classical: the isogeny has kernel
of order $e_{A}^{2}/[M:M']=7$, and the Klein quartic is the quotient of the
Fermat septic curve by a group of order seven; \Cref{fig:simplextype} draws
its Newton triangle and the smooth adapted fan of the proof. For $r=4$,
\Cref{cor:delsarteall}(i) recovers \Cref{prop:delsarte}(ii), with the Fermat
degree lowered to $e_{A}$. Its content lies in dimension five and more, where
otherwise only the Fermat hypersurfaces are covered
(\Cref{prop:aclosure}(iii)), and in the cyclic covers, among them the Calabi--Yau
varieties $\{y^{2}=F_{A}(x)\}$ with $F_{A}$ of degree $2r+4$ in $r+2$
variables. Khovanskii's compactification and Shioda's cover are classical, and
we claim no priority for the two lemmas. What \Cref{thm:simplextype} adds is
that the boundary of the compactification is again of simplex type with the
same exponent $e$, so that the whole cohomology, not only its part of top
Hodge level, comes from Fermat varieties of one degree.

\begin{remark}[Quotients by diagonal automorphisms]\label{rem:delsartequotient}
Let $G\subset T$ be a finite subgroup of the torus of $\PP^{r+1}$ on which
every character $\chi^{a_{i}-a_{0}}$ is trivial, so that $G$ acts on $X_{A}$;
these $G$ are the subgroups of the group $G_{A}$ dual to $M/M'$. Then $f$
descends to a Laurent polynomial of simplex type on $T/G$, whose character
lattice $M_{G}$, the characters trivial on $G$, contains $M'$; so $L'=M'$,
$L=M_{G}$, and $e(f)$ becomes the exponent $e_{G}$ of $M_{G}/M'$, a divisor of
$e_{A}$. As $\PP^{r+1}/G$ is a projective toric variety with torus $T/G$ in
which $X_{A}/G$ is the closure of $V_{f}/G$, \Cref{thm:simplextype} applies
to every projective model $Z$ of $X_{A}/G$ that is the closure of $V_{f}/G$
in a projective toric variety with torus $T/G$ and is smooth, for instance in
one whose fan is smooth and adapted to $f$ (\Cref{lem:simplexchart}): such a
$Z$ lies in $\mathcal{A}$, and the Hodge conjecture holds on it if it holds
on the $X^{j}_{e_{G}}$ with $j\le r$. For $G=G_{A}$, $M_{G}=M'$ and $e_{G}=1$,
so the conjecture holds on $Z$. For the Fermat hypersurface of degree
$m=r+2$ and the group of diagonal automorphisms of determinant one modulo
scalars, the starting point of the mirror construction for $X_{A}$, $M'=mM$
and $M_{G}$ is the set of $x\in M$ whose coordinates are congruent modulo
$m$, so $M_{G}/M'\cong\ZZ/m$ and $e_{G}=m=e_{A}$: the quotient needs the
same Fermat varieties as $X_{A}$, and nothing is gained. For any finite
group $G$ of automorphisms of $X_{A}$, every smooth projective model $X'$ of
$X_{A}/G$ has $\Hdg^{2}(X')=\mathrm{Ab}^{2}(X')$, so \textup{(F3$'$)} holds on
it when it is a fourfold or a fivefold, by \Cref{prop:f3primedominant} and
\Cref{cor:delsarteall}(i).
\end{remark}

\begin{figure}[tbp]
\centering
  \includegraphics[width=\textwidth]{fig_simplextype}
\caption{The Klein quartic $x^{3}y+y^{3}z+z^{3}x$ in the proof of
\Cref{thm:simplextype}. Left: its Newton triangle in the lattice $M$ of the
torus of $\PP^{2}$, in the chart $z=1$, with normalised area $7=[M:M']$ and
three interior points, the genus. Right: the rays of a smooth fan adapted to
it, the three rays of $\PP^{2}$, the three inner normals of the triangle and
the six rays inserted to make every cone unimodular. The curve meets only
the three divisors of the inner normals, in one point each, the lattice
length of the edge, so its Euler number is $-7+3=-4$. With $e_{A}=7$ the
isogeny of the proof makes it the quotient of the Fermat septic by a group
of order seven.}
\label{fig:simplextype}
\end{figure}

The Fermat hypersurfaces are the complete intersections of one diagonal
hypersurface. Complete intersections of several diagonal hypersurfaces of
the same degree lie in $\mathcal{A}$ as well when their coefficients are of
Vandermonde type, because they are quotients of powers of one curve. Fix
$d\ge2$, $N\ge2$ and $1\le c\le N-1$, put $r=N-c$, and call a $c\times(N+1)$
matrix $(a_{ki})$ \emph{of Vandermonde type} if there are distinct
$\lambda_{0},\dots,\lambda_{N}\in\CC$, nonzero $\mu_{0},\dots,\mu_{N}$ and
$g\in\mathrm{GL}_{c}(\CC)$ with
$a_{ki}=\mu_{i}\sum_{l=1}^{c}g_{kl}\lambda_{i}^{l-1}$; for $c\ge2$ this says
that the columns are distinct points of a rational normal curve of degree
$c-1$ in $\PP^{c-1}$. Put $w_{i}=\prod_{l\ne i}(\lambda_{i}-\lambda_{l})^{-1}$
and
\[
  X_{d,r}(\lambda)=\Bigl\{x\in\PP^{N}:\ \sum_{i=0}^{N}w_{i}\lambda_{i}^{k}x_{i}^{d}=0,
  \ \ 0\le k\le c-1\Bigr\}.
\]
For $r=1$ this is the generalised Fermat curve of type $(d,N)$ of
Gonz\'alez-Diez, Hidalgo and Leyton \cite{GDHL09}, the Galois cover of
$\PP^{1}$ with group $H=\mu_{d}^{N+1}/\mu_{d}$, acting by
$u_{i}\mapsto h_{i}u_{i}$, branched over the $\lambda_{i}$. The characters of
$H$ are the $\chi_{a}(h)=\prod_{i}h_{i}^{a_{i}}$ with
$a\in(\ZZ/d)^{N+1}$ and $\sum_{i}a_{i}=0$; write $V_{a}\subset
H^{1}(C,\CC)$ for the eigenspace on which $h^{*}$ acts by $\chi_{a}(h)$, $[a]$
for the orbit of $a$ under $(\ZZ/d)^{\times}$, and $B_{[a]}$ for the image of
a multiple of the idempotent of $\QQ[H]$ attached to $[a]$, an abelian
subvariety of the Jacobian $J(C)$ whose first cohomology is the part
$\bigoplus_{a'\in[a]}V_{a'}$ of $H^{1}(C,\CC)$.

\begin{theorem}[Diagonal complete intersections of Vandermonde
type]\label{thm:vandermonde}
Let $X\subset\PP^{N}$ be the complete intersection of $c$ hypersurfaces
$\sum_{i}a_{ki}x_{i}^{d}=0$, $1\le k\le c$, with $(a_{ki})$ of Vandermonde
type, and let $C=X_{d,1}(\lambda)$ for the $\lambda_{i}$ of $(a_{ki})$.
\begin{enumerate}[label=\textup{(\roman*)},leftmargin=2.6em,itemsep=2pt]
\item $X$ is smooth of dimension $r$ and isomorphic to $X_{d,r}(\lambda)$,
  and $C$ is a smooth curve of genus
  $1+\tfrac12d^{N-1}\bigl((N-1)(d-1)-2\bigr)$.
\item The map
  $\Phi(u^{(1)},\dots,u^{(r)})=\bigl(\prod_{j}u^{(j)}_{i}\bigr)_{i}$ is a
  surjective morphism $C^{r}\to X_{d,r}(\lambda)$, and it induces an
  isomorphism $C^{r}/G\cong X_{d,r}(\lambda)$, where $G$ is the group of
  order $r!\,d^{N(r-1)}$ generated by the permutations of the factors and the
  $(h_{1},\dots,h_{r})\in H^{r}$ with $h_{1}\cdots h_{r}=1$.
\item $X\in\mathcal{A}$, so \textup{(F3$'$)} holds on $X$ in every degree.
\item $\dim V_{a}=\#\{i:a_{i}\ne0\}-2$ for $a\ne0$, and the Hodge conjecture
  holds on $X$ as soon as it holds on the abelian varieties $B_{[a]}$ with
  $\dim V_{a}\ge r$, which have dimension $\tfrac12|[a]|\dim V_{a}$.
\end{enumerate}
\end{theorem}

\begin{proof}
Everything rests on the Lagrange identity: interpolating a polynomial $f$
of degree at most $N$ at the $\lambda_{i}$ and comparing the coefficients of
$t^{N}$ gives $\sum_{i}w_{i}f(\lambda_{i})=0$ when $\deg f\le N-1$.%
\footnote{Lagrange interpolation writes
$f(t)=\sum_{i}f(\lambda_{i})\prod_{l\ne i}(t-\lambda_{l})/(\lambda_{i}-\lambda_{l})$
for $\deg f\le N$, and the coefficient of $t^{N}$ on the right is
$\sum_{i}w_{i}f(\lambda_{i})$. For $N=2$ and $f=1$ it says that
$\sum_{i}\prod_{l\ne i}(\lambda_{i}-\lambda_{l})^{-1}=0$ for any three
distinct numbers. The
identity says that the vector $(w_{i}\lambda_{i}^{k})_{i}$ is orthogonal to the
values of every polynomial of degree at most $N-1-k$, which is how the
equations of $X_{d,r}(\lambda)$ become the condition that the $d$-th powers of
the coordinates are the values of one polynomial of degree at most $r$.}

(i) Replacing the equations by their combinations under $g^{-1}$ and $x_{i}$
by $\nu_{i}x_{i}$ with $\nu_{i}^{d}\mu_{i}=w_{i}$ turns them into those of
$X_{d,r}(\lambda)$. Every $c\times c$ minor of $(a_{ki})$ is
$\det g\prod\mu_{i}$ times a Vandermonde determinant in distinct
$\lambda_{i}$, so any $c$ columns are independent. At a point $x$ of $X$ the
Jacobian matrix is $(d\,a_{ki}x_{i}^{d-1})$; if $S=\{i:x_{i}\ne0\}$ had at most
$c$ elements, $\sum_{i\in S}x_{i}^{d}a_{i}=0$ would be a relation between at
most $c$ columns, so $|S|>c$ and the Jacobian has rank $c$.%
\footnote{The columns of the Jacobian matrix with $i\in S$ are the columns of
$(a_{ki})$ multiplied by the nonzero numbers $d\,x_{i}^{d-1}$, and any $c$ of
the columns of $(a_{ki})$ are independent; so $c$ of them already give rank
$c$. Without the Vandermonde condition this fails: for two diagonal
quadrics in $\PP^{3}$ with proportional columns $i=0,1$, the points with
$x_{2}=x_{3}=0$ and $a_{0}x_{0}^{2}+a_{1}x_{1}^{2}=0$ are singular.} Hence $X$ is a
smooth complete intersection of dimension $r$, connected because $r\ge1$, and
the same holds for $C$. The genus follows from the Riemann--Hurwitz formula
for the cover $C\to\PP^{1}$ of degree $d^{N}$ described below, which has
$N+1$ branch points of ramification index $d$.

(ii) For $1\le k\le N-1$ let $\Lambda_{k}$ be the space of
$y\in\CC^{N+1}$ with $\sum_{i}w_{i}\lambda_{i}^{l}y_{i}=0$ for
$0\le l\le N-k-1$. The vectors $(F(\lambda_{i}))_{i}$ with $\deg F\le k$ lie
in $\Lambda_{k}$ by the identity, since $\deg t^{l}F\le N-1$; they form a space
of dimension $k+1$, since $k<N+1$; and $\Lambda_{k}$ has dimension $k+1$,
the $N-k$ equations having independent Vandermonde rows. So
$\Lambda_{k}=\{(F(\lambda_{i}))_{i}:\deg F\le k\}$, and $C$ and
$X_{d,r}(\lambda)$ are the sets of points whose $d$-th powers lie in
$\Lambda_{1}$ and in $\Lambda_{r}$. The point $u\in C$ determines
$(s:t)\in\PP^{1}$ with $u_{i}^{d}=\kappa(t-\lambda_{i}s)$ for all $i$; this is
the quotient $C\to C/H=\PP^{1}$, whose fibres are the $H$-orbits. If
$(u^{(j)}_{i})^{d}=t_{j}-\lambda_{i}s_{j}$, then
$\Phi(u)_{i}^{d}=F(\lambda_{i})$ for $F(t)=\prod_{j}(t_{j}-ts_{j})$, a nonzero
polynomial of degree at most $r<N+1$; so $\Phi(u)\in X_{d,r}(\lambda)$ and
$\Phi(u)\ne0$, and $\Phi$, given by forms of multidegree $(1,\dots,1)$ with no
common zero, is a morphism. It is $G$-invariant. Conversely, if $x\in
X_{d,r}(\lambda)$, then $(x_{i}^{d})_{i}=(F(\lambda_{i}))_{i}$ for a unique
nonzero $F$ of degree at most $r$, which factors as
$\prod_{j}(t_{j}-ts_{j})$, uniquely up to order and to rescaling the factors.
Choosing $u^{(j)}$ over $(s_{j}:t_{j})$ gives $\Phi(u)_{i}=\epsilon_{i}x_{i}$
with $\epsilon_{i}\in\mu_{d}$, and multiplying $u^{(1)}$ by
$(\epsilon_{i}^{-1})$ gives a preimage.\enlargethispage{2pt} Two preimages have the same roots, so
after a permutation $u'^{(j)}=h_{j}u^{(j)}$ with $h_{j}\in H$, and $\prod_{j}h_{j}$
acts trivially on the coordinates where $x_{i}\ne0$; where $x_{i}=0$ some
$u^{(j)}_{i}$ vanishes and $h_{j,i}$ may be changed freely, so we may take
$\prod_{j}h_{j}=1$. Hence the fibres of $\Phi$ are the $G$-orbits, $\Phi$
induces a finite bijective morphism from the normal variety $C^{r}/G$ to the
smooth variety $X_{d,r}(\lambda)$, and by Zariski's main theorem this is an
isomorphism.%
\footnote{In characteristic zero a finite morphism that is bijective on
closed points is generically \'etale of degree one, hence birational, and a
finite birational morphism onto a normal variety is an isomorphism
\cite[Chapter~III, Corollary~11.4]{Har77}. The group $G$ acts faithfully on $C^{r}$,
since $H$ is the group of deck transformations of $C\to\PP^{1}$, so
$\Phi$ has degree $|G|$.}
The order of $G$ is $r!\,|H|^{r-1}$.

(iii) $C$ is a curve, so $C^{r}\in\mathcal{A}$, and $X$ is the image of
$C^{r}$ under a surjective morphism (\Cref{prop:f3prime}(iii)).

(iv) For $a\ne0$, the forms
$t^{j}\prod_{i}(t-\lambda_{i})^{-b_{i}/d}\,dt$ with $b_{i}\in\{0,\dots,d-1\}$,
$b_{i}\equiv-a_{i}$, and $0\le j\le\sum_{i}b_{i}/d-2$, are holomorphic on $C$
and span the $\chi_{a}$-eigenspace of $H^{1,0}(C)$;%
\footnote{On the affine chart $s=1$ put $u_{i}^{d}=t-\lambda_{i}$, so that the
form is $t^{j}\prod_{i}u_{i}^{-b_{i}}\,dt$ and $h\in H$ multiplies it by
$\prod_{i}h_{i}^{-b_{i}}=\chi_{a}(h)$; it is well defined on
$H=\mu_{d}^{N+1}/\mu_{d}$ because $\sum_{i}b_{i}\equiv0$. Over $\lambda_{i}$
the cover is ramified of index $d$, with a local parameter $z$ satisfying
$z^{d}=t-\lambda_{i}$, and the form is a unit times
$z^{-b_{i}}\,d(z^{d})=d\,z^{d-1-b_{i}}\,dz$, holomorphic as $b_{i}\le d-1$.
Over $t=\infty$, with $w=1/t$ and $v^{d}=w$, it is a unit times
$v^{\,d(\sum_{i}b_{i}/d-j-2)+d-1}\,dv$, holomorphic exactly when
$j\le\sum_{i}b_{i}/d-2$, the exponent of $v$ being an integer. Forms with
different characters, or with different $j$, are independent, and their
number is
$\tfrac12\sum_{a\ne0}\bigl(\#\{i:a_{i}\ne0\}-2\bigr)
=\tfrac12\bigl((N+1)(d^{N}-d^{N-1})-2(d^{N}-1)\bigr)$, which is the genus of
(i); so they form a basis of $H^{1,0}(C)$, and each eigenspace is spanned by
the forms of its character. For $d=2$ and $N=4$ they are the five forms
$dt/\sqrt{(t-\lambda_{i_{1}})\cdots(t-\lambda_{i_{4}})}$, one for each choice
of four of the five indices, on a curve of genus five.}
its dimension
$\sum_{i}\langle-a_{i}/d\rangle-1$ and that of the $\chi_{-a}$-eigenspace add up
to $\#\{i:a_{i}\ne0\}-2$, and $H^{1}(C)^{H}=H^{1}(\PP^{1})=0$. The group $H$
acts trivially on $H^{0}(C)$ and $H^{2}(C)$, so a tensor in
$H^{r}(C^{r})=\bigoplus\bigotimes_{j}H^{k_{j}}(C)$ invariant under the
$(h_{j})$ with $\prod_{j}h_{j}=1$ has all its factors in $H^{1}$, with one
character, or none: taking $h$ in one factor and $h^{-1}$ in another shows
that two factors carry the same character, and a factor in $H^{0}$ or
$H^{2}$ carries the trivial one. Taking the invariants of the permutations,
\[
  H^{r}(C^{r},\CC)^{G}=\CC\,\eta\;\oplus\;\bigoplus_{a\ne0}\textstyle\bigwedge^{r}V_{a},
\]
with $\eta$ the symmetrised product of $r/2$ classes of points when $r$ is
even and $\eta=0$ otherwise, and $W_{[a]}=\bigoplus_{a'\in[a]}\bigwedge^{r}V_{a'}$
is a sub-Hodge structure defined over $\QQ$, zero unless $\dim V_{a}\ge r$.
Let $\psi\colon C\to B_{[a]}$ be an Abel--Jacobi map followed by the
projection, so that $\psi^{*}$ maps $H^{1}(B_{[a]},\QQ)$ isomorphically onto
the $[a]$-part of $H^{1}(C,\QQ)$, and let $\sigma\colon C^{r}\to B_{[a]}$ be
$\sum_{j}\psi\circ\mathrm{pr}_{j}$. Then
$\sigma^{*}(\alpha_{1}\wedge\dots\wedge\alpha_{r})=\prod_{k}\sum_{j}
\mathrm{pr}_{j}^{*}\psi^{*}\alpha_{k}$, whose component in
$H^{1}(C)^{\otimes r}$ is the alternating sum of the
$\psi^{*}\alpha_{1}\otimes\dots\otimes\psi^{*}\alpha_{r}$. The K\"unneth
projectors of a curve are algebraic,%
\footnote{For a curve $C$ with a point $o$ they are $\pi_{0}=[o\times C]$,
$\pi_{2}=[C\times o]$ and $\pi_{1}=[\Delta_{C}]-\pi_{0}-\pi_{2}$, and
$\pi_{1}^{\otimes r}$ is their exterior product on $C^{r}\times C^{r}$. The
average $e_{G}=|G|^{-1}\sum_{g\in G}[\Gamma_{g}]$ of the graphs of the
elements of $G$ acts on $H^{*}(C^{r})$ as the projector onto the invariants,
and $\Phi^{*}$ identifies $H^{*}(X_{d,r}(\lambda),\QQ)$ with
$H^{*}(C^{r},\QQ)^{G}$.} and so is the average $e_{G}$ over $G$,
so $e_{G}\circ\pi_{1}^{\otimes r}\circ\sigma^{*}$ is induced by an algebraic
cycle on $B_{[a]}\times C^{r}$, and it maps $H^{r}(B_{[a]},\QQ)$ onto
$W_{[a]}$: the alternating tensors in the $[a]$-part are the image of
$\pi_{1}^{\otimes r}\circ\sigma^{*}$, and $e_{G}$ keeps those of one
character. A surjective morphism of polarisable Hodge structures maps Hodge
classes onto Hodge classes,%
\footnote{The category of polarisable rational Hodge structures is semisimple,
so the kernel of a surjective morphism $f\colon V\to W$ has a complement
$V'$, a sub-Hodge structure that $f$ maps isomorphically onto $W$; a Hodge
class of $W$ is then the image of a Hodge class of $V'$.} so the Hodge classes of $W_{[a]}$ are algebraic
when those of $B_{[a]}$ are. Now let $x$ be a Hodge class of degree $r$ on
$X=X_{d,r}(\lambda)$, the other degrees being multiples of powers of the
hyperplane class by the Lefschetz hyperplane theorem. Then $\Phi^{*}x$ is a
$G$-invariant Hodge class, its components in the $W_{[a]}$ and along $\eta$
are Hodge classes, being cut out by the algebraic projectors attached to
the action of $H^{r}$, and they are algebraic; so $\Phi^{*}x$ is algebraic,
and so is $x=\Phi_{*}\Phi^{*}x/|G|$.\enlargethispage{2pt} Finally $\dim B_{[a]}$ is half of
$\sum_{a'\in[a]}\dim V_{a'}=|[a]|\dim V_{a}$.
\end{proof}

\begin{corollary}[Two diagonal hypersurfaces]\label{cor:twodiagonal}
\leavevmode
\begin{enumerate}[label=\textup{(\roman*)},leftmargin=2.6em,itemsep=2pt]
\item Every smooth complete intersection of two diagonal hypersurfaces
  $\sum_{i}a_{i}x_{i}^{d}=\sum_{i}b_{i}x_{i}^{d}=0$ of the same degree in
  $\PP^{N}$ is of Vandermonde type, so it lies in $\mathcal{A}$ and
  satisfies \textup{(F3$'$)} in every degree.
\item Let $d\in\{3,4,6\}$ and $N\le6$, or $d=2$ and $N\le12$. Then the
  Hodge conjecture holds on every complete intersection of Vandermonde type
  in $\PP^{N}$ of $N-r$ diagonal hypersurfaces of degree $d$. In particular
  it holds on every smooth complete intersection of two diagonal
  hypersurfaces of degree $3$, $4$ or $6$ in $\PP^{6}$.
\item For $d\in\{3,4,6\}$ and $N=6$, each orbit $[a]$ of order at least
  three with six nonzero coordinates and $\sum_{i}\langle a_{i}/d\rangle=3$ makes $B_{[a]}$ an
  abelian fourfold of Weil type for $K=\QQ(\zeta_{d})$, and $W_{[a]}$ is its
  space of Weil classes. On each of these fourfolds $X$ the Weil classes of
  the $70$, $490$ and $6125$ such orbits, for $d=3,4,6$, give Hodge classes of
  degree four outside $\QQ h^{2}$, so the Hodge classes of $X$ form a space of
  dimension at least $141$, $981$ and $12251$, whereas the divisor classes of
  $X$ are the multiples of $h$.
\end{enumerate}
\end{corollary}

\begin{proof}
(i) If two columns $(a_{i},b_{i})$ and $(a_{j},b_{j})$ were proportional,
or one of them zero, $X$ would contain a point supported on $\{i,j\}$, or
the coordinate point $e_{i}$, at which the Jacobian has rank at most one;
so the $N+1$ columns are distinct points of $\PP^{1}$. Sending a point of
$\PP^{1}$ outside them to infinity writes them as
$\mu_{i}(1,\lambda_{i})$ with distinct $\lambda_{i}$, which is the
Vandermonde form for $c=2$, and \Cref{thm:vandermonde}(iii) applies.
(ii) For $d\in\{2,3,4,6\}$ every $a$ has order $1$, $2$, $3$, $4$ or $6$, so
$|[a]|\le2$, with $|[a]|=1$ for the $a$ of order two, whose $V_{a}$ is real,
so that $\dim B_{[a]}\le\dim V_{a}\le N-1$. When $d=2$ every $a\ne0$ has
order two and an even number of nonzero coordinates, at most $12$ when
$N\le12$, so $\dim B_{[a]}=\tfrac12\dim V_{a}\le5$. So every $B_{[a]}$ has
dimension at most five,%
\footnote{Both bounds are sharp: for $d=3$ and $N=7$ a character with eight
nonzero coordinates has $\dim V_{a}=6$ and $|[a]|=2$, so $\dim B_{[a]}=6$;
for $d=2$ and $N=13$ a character with fourteen nonzero coordinates has
$\dim B_{[a]}=6$.} and the Hodge
conjecture holds on abelian varieties of dimension at most five
\cite{Mar25a,Mar25b,MZ99}; \Cref{thm:vandermonde}(iv) gives the conclusion,
and (i) the last statement.
(iii) For such $a$ the eigenspaces $V_{a}$ and $V_{-a}$ have dimension four
and contain two-dimensional parts of type $(1,0)$, so $K$ acts on
$H^{1,0}(B_{[a]})$ with multiplicities $(2,2)$ and $B_{[a]}$, of dimension
four, is of Weil type; $\bigwedge^{4}V_{a}$ and $\bigwedge^{4}V_{-a}$ are its
Weil lines, of type $(2,2)$ for every $\lambda$, and they lie in
$H^{4}(C^{4})^{G}=\Phi^{*}H^{4}(X)$. The zero coordinate of $a$ is one of
seven, and the other six $a_{i}\in\{1,\dots,d-1\}$ add up to $3d$, which
they do in $20$, $141$ and $1751$ ways for $d=3,4,6$, the coefficients of
$x^{3d-6}$ in $(1+x+\dots+x^{d-2})^{6}$; removing the character of order two
with every $a_{i}=d/2$ leaves $20$, $140$ and $1750$, and $(\ZZ/d)^{\times}=
\{\pm1\}$ pairs them off, since $-a$ again has coordinate sum $3d$; this
gives $7\cdot20/2=70$, $7\cdot140/2=490$ and $7\cdot1750/2=6125$ orbits.
Distinct orbits give independent classes, and the hyperplane class
lies in the part of the trivial character, which gives the bounds; the
Picard group of a complete intersection of dimension at least three is
generated by $h$.
\end{proof}

\Cref{fig:vandermonde} draws the construction. For $c=1$ the theorem is \Cref{prop:aclosure}(iii) again,
by a different route, and for $r=1$ it restates the definition of $C$.

\begin{figure}[tbp]
\centering
\includegraphics[width=\textwidth]{fig_vandermonde}
\caption{\Cref{thm:vandermonde} and \Cref{cor:twodiagonal}. Left: the
generalised Fermat curve $C$ over $\PP^{1}$, its sheets meeting over the
branch points $\lambda_{0},\dots,\lambda_{N}$ (five sheets stand for the
$d^{N}$ of the cover). Top: $C$, the complete intersection
$X\cong C^{r}/G$ with the map $\Phi$ of degree $|G|$, and the decomposition
of $H^{r}(X)$ through $\Phi^{*}$. Bottom: where the Hodge conjecture for $X$
stands at an arbitrary configuration, according to the abelian varieties
$B_{[a]}$ that carry its Hodge classes: known when they have dimension at
most five, which happens for $d\in\{3,4,6\}$ and $N\le6$ and for $d=2$ and
$N\le12$, and open when they are larger, a case of (F2), or carry Weil
classes of $\QQ(\zeta_{e})$ with $\varphi(e)\ge4$, the case
$\mathrm{P2}_{\mathrm{split}}$ of the paper. Foot: the very general
configuration, \Cref{thm:vgvandermonde}.}
\label{fig:vandermonde}
\end{figure}

At a very general configuration the abelian varieties $B_{[a]}$ carry only
the Hodge classes that the monodromy of the curves allows, and the
conjecture can then be decided in every dimension for $d=2$, up to dimension
seven for $d=3$ and $d=4$, and up to dimension five for $d=6$. Write $U\subset\CC^{N+1}$ for the space of
$(N+1)$-tuples of distinct numbers. The varieties $X_{d,r}(\lambda)$ with
$\lambda\in U$ form a smooth projective family over $U$, and so do the
curves $C$ and the maps $\Phi$, with $H$ acting on the family of curves.

\begin{lemma}[Hodge classes of a very general member]\label{lem:vgmonodromy}
Let $\mathbb{V}$ be a polarised variation of $\ZZ$-Hodge structure of weight
$2p$ on a smooth connected quasi-projective variety $S$, and let
$\Pi=\pi_{1}(S,s)$ act on $\mathbb{V}_{s}$. If $s$ lies outside a countable
union of proper Zariski closed subsets of $S$, the space $\Hdg(\mathbb{V}_{s})$
of rational classes of type $(p,p)$ is stable under $\Pi$, and a subgroup of
finite index of $\Pi$ fixes it pointwise.
\end{lemma}

\begin{proof}
The locus of Hodge classes has countably many components
(\Cref{def:hodgelocus}, \Cref{lem:hodgeanalytic}), and the image in $S$ of
each is Zariski closed (\Cref{thm:cdk}); let $s$ avoid the images that are
not all of $S$. For $\alpha\in\Hdg(\mathbb{V}_{s})$ the component through
$(s,\alpha)$ maps onto $S$, locally biholomorphically onto its image, so the
flat section $\tilde\alpha$ through $\alpha$ on the universal cover $\tilde S$
is of type $(p,p)$ on a nonempty open set. The set where it is of type
$(p,p)$ is a closed analytic subset of the connected manifold $\tilde S$, so it
is all of $\tilde S$, and every transport of $\alpha$ along a loop at $s$ is a
Hodge class at $s$. The group $\Pi$ therefore preserves
$\Hdg(\mathbb{V}_{s})$, the lattice $\mathbb{V}_{s,\ZZ}\cap\Hdg(\mathbb{V}_{s})$
and the flat form $\tilde Q$ of the proof of \Cref{lem:boundednorm}, which is
positive definite on the real classes of type $(p,p)$; so it acts on
$\Hdg(\mathbb{V}_{s})$ through a finite group, whose kernel is the subgroup
of the statement.
\end{proof}

For the characters of order three, four and six the monodromy of the
curves $D_{a}$ is determined by induction on the number of branch points. The
step lets two branch points collide, and a lemma on two copies of $\SL$ of a
hyperplane closes it. The base has four branch points, two of them with
opposite exponents, whose full twist is a transvection
(\Cref{fig:cyclicmerge}).

\begin{setupenv}[Cyclic covers]\label{setup:cyclic}
Let $m\in\{3,4,6\}$, $\zeta=e^{2\pi i/m}$ and $K=\QQ(\zeta)$, and let
$a=(a_{1},\dots,a_{k})$ with $a_{i}\in\ZZ/m$ nonzero, $\sum_{i}a_{i}=0$ and
$\gcd(m,a_{1},\dots,a_{k})=1$, so that $a$ has order $m$. For $\lambda$ in
the space $\mathrm{Conf}_{k}$ of $k$-tuples of distinct complex numbers let
$D_{a}$ be the cyclic cover $y^{m}=\prod_{i}(t-\lambda_{i})^{a_{i}}$ of
$\PP^{1}$, with representatives $0<a_{i}<m$; it is unramified over
$\infty$. Let $V_{a}\subset H^{1}(D_{a},\CC)$ be the eigenspace of the deck
group on which $H^{1,0}$ is spanned by the forms
$t^{j}\prod_{i}(t-\lambda_{i})^{-b_{i}/m}\,dt$, where $0<b_{i}<m$, $b_{i}\equiv
-a_{i}$ and $0\le j\le\sum_{i}b_{i}/m-2$ \cite[\S2]{DM86}, and put
$n=\dim V_{a}$,
$p=\dim V_{a}^{1,0}=\sum_{i}b_{i}/m-1$ and $q=n-p$. Let $\Gamma_{a}\subset
\GL(V_{a})$ be the image of the pure braid group
$P_{k}=\pi_{1}(\mathrm{Conf}_{k})$ and $G_{a}^{0}$ the identity component of
its Zariski closure.
\end{setupenv}

\begin{lemma}[A model of the monodromy]\label{lem:foxmodel}
Let $F_{k}$ be the free group on loops $x_{1},\dots,x_{k}$ around
$\lambda_{1},\dots,\lambda_{k}$, put $t_{i}=\zeta^{a_{i}}$, let
$\partial\colon K^{k}\to K$ map $e_{i}$ to $t_{i}-1$, and let
$E_{a}=\ker\partial/K\ell$, where $\ell=\sum_{i}t_{1}\cdots t_{i-1}e_{i}$. A pure
braid $\beta$ acts on $F_{k}$ by Artin's automorphism and on $K^{k}$ by
$v\mapsto vJ_{\beta}$, where $J_{\beta}$ is the matrix of Fox derivatives
$\partial\beta(x_{i})/\partial x_{j}$ evaluated at $x_{j}\mapsto t_{j}$. This
action preserves $\ker\partial$ and fixes $\ell$, and the representation of
$P_{k}$ on $E_{a}$, of dimension $k-2$, is isomorphic to $V_{a}$, to its dual
or to its complex conjugate. So $n=k-2$, and $G_{a}^{0}\supseteq\SL(V_{a})$ if
and only if the identity component of the Zariski closure of the image of
$P_{k}$ in $\GL(E_{a})$ contains $\SL(E_{a})$.
\end{lemma}

\begin{proof}
Fox's formula $\sum_{j}(\partial w/\partial x_{j})(x_{j}-1)=w-1$ \cite{Fox53},
evaluated at $x_{j}\mapsto t_{j}$, shows that the vector
$\sum_{j}(\partial w/\partial x_{j})(t)\,e_{j}$ of a word $w$ with $t(w)=1$ lies
in $\ker\partial$. It is the class of the lift of the loop $w$ in the summand
of $H_{1}(D_{a}^{\circ},K)$ on which the deck group acts through
$x_{i}\mapsto t_{i}$, where $D_{a}^{\circ}$ is the unramified part of the
cover, and that summand is $H_{1}(F_{k},K_{t})=\ker\partial$, computed from the
free resolution of length one of the augmentation ideal of a free group.%
\footnote{By Shapiro's lemma $H_{1}(D_{a}^{\circ},K)=H_{1}(F_{k},K[\mu_{m}])$,
and its summand for the character $x_{i}\mapsto t_{i}$ is $H_{1}(F_{k},K_{t})$.
The augmentation ideal of $\ZZ[F_{k}]$ is free on the $x_{i}-1$, so
$H_{1}(F_{k},K_{t})$ is the kernel of $\bigoplus_{i}Ke_{i}\to K$,
$e_{i}\mapsto t_{i}-1$, the chain $e_{i}$ standing for the loop $x_{i}$ and a word
$w$ for $\sum_{j}(\partial w/\partial x_{j})(t)e_{j}$.} Filling in the points
over $\lambda_{i}$ adds no relation: the loop around such a point lifts
$x_{i}^{e}$, $e$ the order of $t_{i}$, whose vector
$(1+t_{i}+\dots+t_{i}^{e-1})e_{i}$ is zero. Filling in the $m$ points over
$\infty$, where the cover is unramified, kills the lift of $x_{1}\cdots
x_{k}$, whose vector is $\ell\ne0$. So $E_{a}$ is the summand of
$H_{1}(D_{a},K)$ for the character $x_{i}\mapsto t_{i}$, of dimension $k-2$. A
pure braid fixes $x_{1}\cdots x_{k}$ and maps each $x_{i}$ to a conjugate of
itself, so it preserves $t$, and the chain rule for Fox derivatives makes
$\beta\mapsto J_{\beta}$ the action of $P_{k}$ on this summand.%
\footnote{Artin's automorphism of the half twist $\sigma_{i}$ of the $i$-th
and $(i+1)$-st points is $x_{i}\mapsto x_{i}x_{i+1}x_{i}^{-1}$,
$x_{i+1}\mapsto x_{i}$, the other generators being fixed; $P_{k}$ is generated
by the full twists $A_{st}=\sigma_{t-1}\cdots\sigma_{s+1}\sigma_{s}^{2}
\sigma_{s+1}^{-1}\cdots\sigma_{t-1}^{-1}$. For a pure braid the chain rule
reads $J_{\beta\beta'}=J_{\beta'}J_{\beta}$ or $J_{\beta}J_{\beta'}$, according
to the convention for composing braids, with no twist of the coefficients,
since $t\circ\beta=t$.} Poincar\'e duality and the choice of the generator
of the deck group identify it with $V_{a}$, its dual or its conjugate, and
the automorphisms $g\mapsto{}^{t}g^{-1}$ and $g\mapsto\bar g$ of $\GL$ preserve
$\SL$.
\end{proof}

\begin{lemma}[Collision of two branch points]\label{lem:cabling}
In \Cref{lem:foxmodel} suppose that $a_{1}+a_{2}\ne0$, and let
$a'=(a_{1}+a_{2},a_{3},\dots,a_{k})$. Let $I=(1-t_{2})e_{1}-(1-t_{1})e_{2}$,
the cycle of a disc around $\lambda_{1}$ and $\lambda_{2}$, and let
$Y\subset E_{a}$ be the image of the cycles outside that disc, the vectors of
$\ker\partial$ with $v_{2}=t_{1}v_{1}$. Then $E_{a}=Y\oplus KI$, and
$\iota(v_{1},v_{3},\dots,v_{k})=(v_{1},t_{1}v_{1},v_{3},\dots,v_{k})$ induces an
isomorphism $E_{a'}\to Y$. If $\beta\in P_{k}$ is obtained from
$\beta'\in P_{k-1}$ by doubling its first strand, then $\beta$ preserves $Y$
and $KI$, acts on $Y$ as $\beta'$ acts on $E_{a'}$, and multiplies $I$ by a
power of $\zeta$.
\end{lemma}

\begin{proof}
The vectors of $Y$ have $(v_{1},v_{2})$ proportional to $(1,t_{1})$, and
$(1-t_{2},t_{1}-1)$ is proportional to $(1,t_{1})$ only if $t_{1}t_{2}=1$; so
$I\notin Y$. The map $\iota$ carries the kernel of $\partial'$, where
$\partial'(e_{1})=t_{1}t_{2}-1$, onto the cycles outside the disc, because
$v_{1}(t_{1}t_{2}-1)=v_{1}(t_{1}-1)+t_{1}v_{1}(t_{2}-1)$, and it carries the
vector $\ell'$ of $a'$ to $\ell$; as $\dim E_{a'}+1=\dim E_{a}$, both
assertions about $Y$ follow. The doubled braid moves a disc containing
$\lambda_{1},\lambda_{2}$ rigidly. Its automorphism therefore maps
$x_{1}\mapsto wx_{1}w^{-1}$, $x_{2}\mapsto wx_{2}w^{-1}$ and $x_{j}\mapsto
u_{j}$ for $j\ge3$, where $y\mapsto w'yw'^{-1}$, $x_{j}\mapsto u_{j}'$ is the
automorphism of $\beta'$ on the free group on $y,x_{3},\dots,x_{k}$, and
$w,u_{j}$ are $w',u_{j}'$ with $y$ replaced by $x_{1}x_{2}$. For such a word
$u$ the chain rule gives $\partial u/\partial x_{1}=\partial u/\partial y$
and $\partial u/\partial x_{2}=(\partial u/\partial y)\,x_{1}$, which at $t$ is
the relation $v_{2}=t_{1}v_{1}$; together with
$\partial(wx_{i}w^{-1})/\partial x_{j}=(1-t_{i})\,\partial w/\partial
x_{j}+t(w)\delta_{ij}$ at $t$, this shows that the rows of $J_{\beta}$ for
$j\ge3$, and the first row plus $t_{1}$ times the second, are the images under
$\iota$ of the rows of $J_{\beta'}$; this is the equivariance of $\iota$. The
same formula gives $IJ_{\beta}=(1-t_{2})\bigl((1-t_{1})\partial w+t(w)e_{1}\bigr)
-(1-t_{1})\bigl((1-t_{2})\partial w+t(w)e_{2}\bigr)=t(w)I$, and $t(w)$ is a
power of $\zeta$.
\end{proof}

\begin{lemma}[Two hyperplanes]\label{lem:hyperplanes}
Let $E=Y\oplus L$ be a complex vector space with $\dim Y\ge2$ and
$\dim L=1$, let $S\subset\SL(E)$ be the subgroup acting on $Y$ through
$\SL(Y)$ and trivially on $L$, and let $\gamma\in\GL(E)$ satisfy
$\gamma L\ne L$ and $\gamma Y\ne Y$. Then the group $T$ generated by $S$ and
$\gamma S\gamma^{-1}$ acts irreducibly on $E$. If $\dim Y\ge3$, then
$T=\SL(E)$. If $\dim Y=2$, every connected semisimple subgroup of $\SL(E)$
that contains $T$ is $\SL(E)$.
\end{lemma}

\begin{proof}
The group $T$ is connected. As a representation of $S$, $E$ is
the sum of the irreducible nontrivial $Y$ and the trivial $L$, so its
$S$-stable subspaces are $0$, $L$, $Y$ and $E$, and those of
$\gamma S\gamma^{-1}$ are $0$, $\gamma L$, $\gamma Y$ and $E$; by the
hypotheses and the dimensions only $0$ and $E$ are in both lists, and $T$ is
irreducible.

Let $\dim Y\ge3$. As a representation of $\mathfrak{sl}(Y)$,
\[
  \mathfrak{sl}(E)=\mathfrak{sl}(Y)\oplus\mathrm{Hom}(L,Y)\oplus
  \mathrm{Hom}(Y,L)\oplus\CC z,
\]
$z$ acting by $1$ on $Y$ and by $-\dim Y$ on $L$, and the four summands, the
adjoint representation, $Y$, $Y^{*}$ and the trivial one, are pairwise
non-isomorphic since $\dim Y\ge3$. The Lie algebra $\mathfrak{t}$ of $T$
contains $\mathfrak{sl}(Y)$ and is stable under it, so it is
$\mathfrak{sl}(Y)$ plus some of the other three summands. Without
$\mathrm{Hom}(L,Y)$ it preserves $L$, and without $\mathrm{Hom}(Y,L)$ it
preserves $Y$; both are excluded by irreducibility. The bracket of a nonzero
map $L\to Y$ with a map $Y\to L$ not vanishing on its image has a nonzero
component along $z$, so $\mathfrak{t}=\mathfrak{sl}(E)$ and $T=\SL(E)$.

Let $\dim Y=2$ and let $G\supseteq T$ be connected and semisimple. Then $E$
is a faithful irreducible representation of dimension $3$ of the Lie algebra
$\mathfrak{g}$ of $G$. An irreducible representation of a product of simple
Lie algebras is a tensor product of irreducible representations of the
factors, and $3$ is prime, so $\mathfrak{g}$ is simple. It lies in
$\mathfrak{sl}(E)$, of dimension $8$ and rank $2$, and the simple Lie
algebras of rank at most $2$ are $\mathfrak{sl}_{2}$, $\mathfrak{sl}_{3}$,
$\mathfrak{so}_{5}$ and $\mathfrak{g}_{2}$, of dimensions $3$, $8$, $10$ and
$14$. If $\mathfrak{g}\cong\mathfrak{sl}_{2}$, then $\mathfrak{g}$ equals the
Lie algebra of $S$, both being of dimension $3$, so $G=S$, which preserves
$L$, against the irreducibility of $T$. Hence $\mathfrak{g}=\mathfrak{sl}(E)$
and $G=\SL(E)$.
\end{proof}

\begin{lemma}[Merging two branch points]\label{lem:merge}
Let $m$ and $a$ be as in \Cref{setup:cyclic}, with $k\ge6$ and $p,q\ge1$.
Then there are $i\ne j$ with $a_{i}+a_{j}\ne0$ such that the vector $a'$
obtained by replacing $a_{i}$ and $a_{j}$ by $a_{i}+a_{j}$ has order $m$, and
its numbers $p'$ and $q'$ are at least $1$.
\end{lemma}

\begin{proof}
Put $f_{i}=b_{i}/m\in[1/m,1-1/m]$, so that $p+1=\sum_{i}f_{i}$ and
$q+1=\sum_{i}(1-f_{i})$. Merging $a_{i}$ and $a_{j}$ keeps $p$ and lowers $q$
by one if $f_{i}+f_{j}<1$, and lowers $p$ by one and keeps $q$ if
$f_{i}+f_{j}>1$. Since $a_{i}+a_{j}=-\sum_{l\ne i,j}a_{l}$, the order of $a'$ is
that of the subgroup generated by the $a_{l}$ with $l\ne i,j$, the rest of
$a$. Replacing $a$ by $-a$ exchanges $p$ and $q$, so we may assume $q\ge p$,
hence $q\ge2$ as $p+q=k-2\ge4$; if $p\ge2$, every merge keeps $p',q'\ge1$
and only the order matters.

\emph{The order.} For $m=3$ two equal coordinates exist, their sum is not
$0$, and any nonzero rest generates $\ZZ/3$. For $m=4$ the number of odd
coordinates is even and positive; if it is at least $4$, merge two equal odd
ones, and if it is $2$, merge an odd one with one of the at least four
coordinates equal to $2$. For $m=6$ the rest generates $\ZZ/6$ exactly when
it contains an odd coordinate and one prime to $3$. The odd coordinates
are even in number and the coordinates prime to $3$ have sum divisible by
$3$, so each kind occurs at least twice. If exactly two coordinates are
odd, the at least four others lie in $\{2,4\}$; merge two equal ones among
them. If exactly two are prime to $3$, the others equal $3$; merge a $3$
with one of the two. Otherwise each kind occurs at least three times, and
any two coordinates with nonzero sum can be merged; they exist, since
pairwise sums $0$ would force every coordinate to be $3$.

\emph{The case $p=1$.} Then $\sum_{i}b_{i}=2m$. Some $b_{i}$ equals $1$,
since otherwise $\sum_{i}b_{i}\ge2k\ge12\ge2m$, with equality only if $m=6$
and every coordinate is $4$, when $a$ has order three. Keep such a coordinate,
$a_{i}=-1$, which generates $\ZZ/m$. The other $k-1\ge5$ coordinates have
$b$-sum $2m-1$, so the two smallest values of $b$ among them add up to at most
$2(2m-1)/5<m$; merge these two coordinates. Their sum is not $0$, $p'=p=1$,
$q'=q-1=k-4\ge2$, and the rest contains $a_{i}$, so $a'$ has order $m$.
\end{proof}

\begin{lemma}[Five branch points]\label{lem:mergefive}
Let $m$ and $a$ be as in \Cref{setup:cyclic}, with $k=5$ and $p,q\ge1$. Then
$a_{i}+a_{j}=0$ for some $i\ne j$, and the other three coordinates can be
labelled $x,y,z$ so that $x+y\ne0$ and $a'=(x+y,z,a_{i},a_{j})$ has order $m$
and $p'=q'=1$.
\end{lemma}

\begin{proof}
Replacing $a$ by $-a$ we may assume $p=1$, so that $\sum_{l}b_{l}=2m$ with
five terms in $[1,m-1]$; and $a_{i}+a_{j}=0$ means $b_{i}+b_{j}=m$. For $m=3$
the terms lie in $\{1,2\}$ and exactly one is $2$; take it and a $1$. For
$m=4$, take a $3$ and a $1$ if both occur. If no term is $3$, the terms lie in
$\{1,2\}$ and at least three are $2$, which gives $2+2$; if no term is $1$,
they lie in $\{2,3\}$ and add up to at least $10$, which is impossible. For
$m=6$ suppose that no two terms add up to $6$. Then at most one term is $3$,
and neither $1$ and $5$ nor $2$ and $4$ both occur. Terms in $\{1,2,3\}$ with
at most one $3$ add up to at most $11$, and terms in $\{3,4,5\}$ to at least
$19$. Terms in $\{1,3,4\}$, respectively $\{2,3,5\}$, cannot add up to $12$,
since the number $y$ of terms $4$, respectively $5$, would satisfy
$3y\in\{7,5\}$, respectively $3y\in\{2,1\}$. This proves the first assertion.

The other three coordinates add up to $0$ and have $b$-sum $m$, so any two of
them, $x$ and $y$, have $x+y=-z\ne0$ and $b$-sum less than $m$; merging them
keeps $p'=1$, so $q'=1$, and $a'$ contains $c=a_{i}$ and $a_{j}=-c$. The order
of $a'$ is that of the subgroup generated by $z$ and $c$. If $c$ generates
$\ZZ/m$, any choice of $z$ will do. Otherwise $c$ is even, or $m=6$ and $c=3$.
As $a$ has order $m$, one of $x,y,z$ is then odd, respectively prime to $3$,
and we take it as $z$; then $z$ and $c$ generate $\ZZ/m$.
\end{proof}

\begin{proposition}[Monodromy of cyclic covers]\label{prop:cyclicmonodromy}
In \Cref{setup:cyclic}, if $p=0$ or $q=0$ then $\Gamma_{a}$ is finite, and if
$n\ge3$ and $p,q\ge1$ then $G_{a}^{0}\supseteq\SL(V_{a})$.
\end{proposition}

\begin{proof}
If $pq=0$, the polarisation is definite on $V_{a}$ up to sign, and
$\Gamma_{a}$ preserves it and the $\ZZ[\zeta]$-lattice spanned by
$H^{1}(D_{a},\ZZ)$; a discrete subgroup of a compact group is finite. Assume
$p,q\ge1$ and use the model of \Cref{lem:foxmodel}. Three facts are used
throughout. If $u\in\Gamma_{a}$ is unipotent, the Zariski closure of its
powers is the group $\exp(\tau\log u)$, so $\log u$ lies in the Lie algebra
of $G_{a}^{0}$. The group $\Gamma_{a}$ normalises $G_{a}^{0}$. And $G_{a}^{0}$
is semisimple: by Andr\'e's theorem \cite[Theorem~1]{And92} the identity
component of the Zariski closure of the monodromy group of
$H^{1}(D_{a},\QQ)$ is a normal subgroup of the derived group of the
Mumford--Tate group of a very general fibre, hence semisimple, and
$G_{a}^{0}$ is its image in $\GL(V_{a})$, the monodromy commuting with the
deck group.

\emph{Four points.} Let $k=4$, so $n=2$, and suppose that two coordinates
of $a$ are opposite; relabel them $a_{1}$ and $a_{2}$, so that $t_{1}t_{2}=1$.
The full twist $A_{12}$ maps $x_{i}\mapsto wx_{i}w^{-1}$ for $i=1,2$, with
$w=x_{1}x_{2}$, and fixes $x_{3},x_{4}$. As $t(w)=1$, the formula for
$\partial(wx_{i}w^{-1})/\partial x_{j}$ in the proof of \Cref{lem:cabling}
gives $vJ_{A_{12}}=v+\varphi(v)\,\delta$, with $\delta=e_{1}+t_{1}e_{2}$ and
$\varphi(v)=(1-t_{1})v_{1}+(1-t_{2})v_{2}$, which equals
$\sum_{j\ge3}(t_{j}-1)v_{j}$ on $\ker\partial$. Here $\delta\in\ker\partial$,
$\varphi(\delta)=0$, $\varphi(\ell)=t_{1}\cdots t_{k}-t_{1}t_{2}=0$, the
vector $\delta$ is not a multiple of $\ell$, whose third coordinate is
$t_{1}t_{2}=1$, and $\varphi\bigl((t_{3}-1)e_{1}-(t_{1}-1)e_{3}\bigr)
=-(t_{1}-1)(t_{3}-1)\ne0$. So $A_{12}$ acts on $E_{a}$ as a unipotent
$u\ne1$, and $G_{a}^{0}$ is a connected semisimple subgroup of
$\SL(E_{a})\cong\SL_{2}$ containing $\exp(\tau\log u)$. The connected
subgroups of $\SL_{2}$ of dimension one or two are solvable, so
$G_{a}^{0}=\SL(E_{a})$.

\emph{The induction.} Let $n\ge3$. If $n=3$, \Cref{lem:mergefive} gives,
after a relabelling, a vector $a'=(a_{1}+a_{2},a_{3},a_{4},a_{5})$ of order $m$
with $a_{1}+a_{2}\ne0$, $p'=q'=1$ and two opposite coordinates, so the
identity component of the Zariski closure of the monodromy of $E_{a'}$ is
$\SL(E_{a'})$ by the case of four points. If $n\ge4$, \Cref{lem:merge} and
a relabelling give $a_{1}+a_{2}\ne0$ and $a'=(a_{1}+a_{2},a_{3},\dots,a_{k})$
of order $m$ with $p',q'\ge1$, and by induction on $n$ that identity
component contains $\SL(E_{a'})$. Let $H$ be the identity component of the
Zariski closure of the doubled braids of \Cref{lem:cabling}. Then
$H\subseteq G_{a}^{0}$, $H$ preserves $Y$ and $KI$ and acts trivially on
$KI$, where the doubled braids act through a finite group, and its image in
$\GL(Y)$ is the identity component of the Zariski closure of the monodromy of
$E_{a'}$, which contains $\SL(Y)$. So $G_{a}^{0}$ contains the group $S$ of
\Cref{lem:hyperplanes} for $E_{a}=Y\oplus KI$. Let $\gamma=A_{23}$, the full
twist of the second and third points; its automorphism fixes $x_{1}$ and
maps $x_{2}\mapsto x_{2}x_{3}x_{2}x_{3}^{-1}x_{2}^{-1}$ and $x_{3}\mapsto
x_{2}x_{3}x_{2}^{-1}$. Then $\gamma(I)=I+(1-t_{1})t_{2}I_{23}$ with
$I_{23}=(1-t_{3})e_{2}-(1-t_{2})e_{3}$, and $\gamma(KI)\ne KI$, since $I$,
$I_{23}$ and $\ell$ are independent, $\ell$ having a nonzero fourth
coordinate. The functional $v\mapsto v_{2}-t_{1}v_{1}$ vanishes on $\ell$ and
cuts out $Y$, and at $\gamma(v)$ for $v\in Y$ it equals
$(t_{3}-1)(t_{1}t_{2}v_{1}-v_{3})$, which is nonzero for the cycle
$(1-t_{4})e_{3}-(1-t_{3})e_{4}$ of $Y$; so $\gamma Y\ne Y$. As $\gamma$
normalises $G_{a}^{0}$, this group contains $S$ and $\gamma S\gamma^{-1}$,
and \Cref{lem:hyperplanes}, with $\dim Y=n-1$, gives
$G_{a}^{0}=\SL(E_{a})$: for $n\ge4$ directly, and for $n=3$ because
$G_{a}^{0}$ is connected and semisimple.
\end{proof}

For $m=3$ the proposition also follows from the $\ell$-adic theorem of
Achter and Pries \cite[Corollary~3.10]{AP07}, which identifies the monodromy
of trielliptic curves with a special unitary group.

\begin{lemma}[The discriminant of the new part]\label{lem:newdisc}
In \Cref{setup:cyclic} let $n$ be even with $p=q=n/2$, let
$L=H^{1}(D_{a},\ZZ)$ with the cup product $Q$, and let
$L_{a}=L\cap(V_{a}\oplus\overline{V_{a}})$. Then $L_{a}$ is a free
$\ZZ[\zeta]$-module of rank $n$, the $K$-hermitian form $H$ with
$Q=\tfrac12\Tr_{K/\QQ}(\delta^{-1}H)$ on it, $\delta=\sqrt{-1}$ for $m=4$ and
$\delta=\sqrt{-3}$ otherwise (\Cref{lem:hermform}), has signature
$(n/2,n/2)$, and the class of $(-1)^{n/2}\det H$ in
$\QQ^{\times}/\Nm(K^{\times})$ is trivial for $m=3$ and $m=4$; for $m=6$ it is
trivial if and only if the number of $i$ with $a_{i}=3$ is even.
\end{lemma}

\begin{proof}
The deck transformation $g$ is an isometry of the unimodular lattice $L$, so
its eigenspaces for eigenvalues $\mu,\nu$ are orthogonal unless $\mu\nu=1$.
Hence $L_{a}$, the part for the primitive $m$-th roots of unity, and the
part $L^{\mathrm{old}}$ for the other eigenvalues are orthogonal primitive
sublattices, each the orthogonal complement of the other, and
$|\det(L_{a})|=|\det(L^{\mathrm{old}})|$, which is $1$ for $m=3$, when
$L^{\mathrm{old}}=0$.%
\footnote{If $M\subset L$ is primitive and $L$ is unimodular, restriction
$L\to M^{\vee}$ is onto with kernel $M^{\perp}$, so $L/(M\oplus M^{\perp})\cong
M^{\vee}/M$, and likewise $\cong(M^{\perp})^{\vee}/M^{\perp}$; the order of
$M^{\vee}/M$ is $|\det M|$.} For a quotient map $\pi\colon D_{a}\to D'$ of
prime degree $e$ that is ramified, $\pi^{*}H^{1}(D',\ZZ)$ is primitive in $L$
and the cup product on it is $e$ times that of $D'$.%
\footnote{A loop in $D'$ avoiding the branch points, composed with a loop
around a branch point, whose monodromy generates the group of order $e$,
lifts to a closed loop with the same image in $H_{1}(D',\ZZ)$; so
$\pi_{*}\colon H_{1}(D_{a},\ZZ)\to H_{1}(D',\ZZ)$ is onto, and its transpose
$\pi^{*}$ is a split injection. The factor $e$ is $\pi_{*}\pi^{*}=e$.} For
$m=4$, $L^{\mathrm{old}}=\pi^{*}H^{1}(D',\ZZ)$ for the quotient $D'$ by $g^{2}$,
ramified over the $\lambda_{i}$ with $a_{i}$ odd, so
$|\det L^{\mathrm{old}}|=2^{2g'}$, $g'$ the genus of $D'$. For $m=6$, the
quotients $D_{2}$ by $g^{2}$ and $D_{3}$ by $g^{3}$ are ramified over the
$\lambda_{i}$ with $3\nmid a_{i}$ and with $a_{i}$ odd, which exist since $a$
has order six, and $L^{\mathrm{old}}$ contains $A\oplus B$ with
$A=\pi_{2}^{*}H^{1}(D_{2},\ZZ)$ and $B=\pi_{3}^{*}H^{1}(D_{3},\ZZ)$, of determinants
$3^{2g_{2}}$ and $2^{2g_{3}}$; the quotient $L^{\mathrm{old}}/(A\oplus B)$ embeds in
$A^{\vee}/A$ and in $B^{\vee}/B$, of coprime orders, so $L^{\mathrm{old}}=A\oplus B$
and $|\det L^{\mathrm{old}}|=3^{2g_{2}}2^{2g_{3}}$.

On $L_{a}$ the deck transformation satisfies the cyclotomic equation of $\zeta$,
so $L_{a}$ is a torsion-free module over the principal ideal domain
$\ZZ[\zeta]$, free of rank $n$; the signature of $H$ is $(p,q)$
(\Cref{rem:twosignatures}). The determinant of the trace form of a hermitian
lattice gives $|\det(L_{a})|=|d_{K}|^{n}\Nm(\tfrac{1}{2\delta})^{n}(\det H)^{2}$,
that is $(\det H)^{2}$ for $K=\QQ(i)$ and $(\det H)^{2}/4^{n}$ for
$K=\QQ(\sqrt{-3})$.%
\footnote{For a free $\cO_{K}$-module with hermitian Gram matrix $M$ and
$c\in K^{\times}$, the $\ZZ$-bilinear form $\Tr_{K/\QQ}(c\,v^{T}M\bar w)$ has
determinant of absolute value $|d_{K}|^{n}\Nm(c)^{n}(\det M)^{2}$, where
$d_{K}$ is the discriminant of $K$: for $n=1$ this is the discriminant of the
trace form of $\cO_{K}$ scaled by $\Nm(cM)$. Changing the basis by
$P\in\GL_{n}(K)$ multiplies both sides by $\Nm_{K/\QQ}(\det P)^{2}$, and a
basis in which $M$ is diagonal makes the form an orthogonal sum of forms of
rank one. Here $d_{K}=-4$ and $\Nm(\tfrac{1}{2i})=\tfrac14$, or $d_{K}=-3$ and
$\Nm(\tfrac{1}{2\sqrt{-3}})=\tfrac1{12}$.}
Hence $(-1)^{n/2}\det H$ equals $2^{n}$ for $m=3$, $2^{g'}$ for $m=4$ and
$2^{n+g_{3}}3^{g_{2}}$ for $m=6$, the sign being that of a form of signature
$(n/2,n/2)$. For $m=3$ this is $\Nm(2^{n/2})$, and for $m=4$ a norm as
$2=\Nm(1+i)$. For $m=6$, $3=\Nm(\sqrt{-3})$ and $4=\Nm(2)$
are norms and $2$ is not, since $x^{2}+xy+y^{2}$ is even only when $x$ and $y$
are, so that every norm has even $2$-adic valuation; so the class is trivial
exactly when $g_{3}$ is even. By the Riemann--Hurwitz formula $D_{3}$, the
triple cover $y^{3}=\prod_{i}(t-\lambda_{i})^{a_{i}}$ branched over the $i$
with $3\nmid a_{i}$, has genus $k-c_{3}-2=n-c_{3}$, where $c_{3}$ is the number
of $i$ with $a_{i}=3$; as $n$ is even, this is the assertion.
\end{proof}

\begin{figure}[tbp]
\centering
\includegraphics[width=\textwidth]{fig_cyclicmerge}
\caption{The proof of \Cref{prop:cyclicmonodromy}. Left: two branch points
$\lambda_{1},\lambda_{2}$ collide inside a disc (\Cref{lem:cabling}); the
cycles outside the disc form the hyperplane $Y\cong E_{a'}$, the cycle $I$
of the disc spans the complementary line, and the twist $\gamma$ of the
second and third points moves both. Middle: the two copies of $\SL(Y)$ of
\Cref{lem:hyperplanes}, which generate $\SL(E_{a})$ when $\dim Y\ge3$. Right: the induction on
$n=k-2$, each merge lowering $p$ or $q$ by one (\Cref{lem:merge},
\Cref{lem:mergefive}), down to $n=2$, where the full twist of two points
with opposite exponents is a transvection.}
\label{fig:cyclicmerge}
\end{figure}

\begin{theorem}[Very general diagonal complete
intersections]\label{thm:vgvandermonde}
Keep the notation of \Cref{thm:vandermonde}, let $r\ge2$ be even, and let
$\lambda\in U$ lie outside a countable union of proper Zariski closed
subsets of $U$. For a Hodge class $x$ of degree $r$ on $X=X_{d,r}(\lambda)$,
write $x_{[a]}$ for the component of $\Phi^{*}x$ in $W_{[a]}$.
\begin{enumerate}[label=\textup{(\roman*)},leftmargin=2.6em,itemsep=2pt]
\item If $a$ has order two, then $x_{[a]}$ is a rational multiple of the
  image of $\theta_{a}^{r/2}$, where $\theta_{a}$ is the class of a
  polarisation of $B_{[a]}$, under the correspondence of the proof of
  \Cref{thm:vandermonde}(iv); in particular $x_{[a]}$ is algebraic.
\item If $a$ has order $m\in\{3,4,6\}$ and $x_{[a]}\ne0$, then $\dim V_{a}=r$
  and $\dim V_{a}^{1,0}=r/2$. In that case $B_{[a]}$ is an abelian variety of
  dimension $r$ of Weil type for $\QQ(\zeta_{m})$, and $W_{[a]}$ is the image
  of its space of Weil classes. Its discriminant is trivial for $m=3,4$, and
  for $m=6$ exactly when an even number of the coordinates of $a$ equal
  $d/2$.
\item The Hodge conjecture holds on $X_{d,r}(\lambda)$ for every $N$ when
  $d=2$, when $d\in\{3,4\}$ and $r\le6$, and when $d=6$ and $r\le4$. For
  $d=6$ and $r=6$ the Hodge classes are algebraic modulo the spaces
  $W_{[a]}$ of the characters of order six of (ii) with an odd number of
  coordinates $3$, which are spaces of Weil classes of abelian sixfolds of
  non-split Weil type for $\QQ(\sqrt{-3})$ and occur for every $N$. For
  $d\in\{2,3,4,6\}$ the Hodge classes of degree $r$ form a space of dimension
  \[
    1+\binom{N+1}{r+2}T_{d}(r+2)+\epsilon_{d}\sum_{j\ge r/2+2}\binom{N+1}{2j},
  \]
  where $T_{d}(k)$ is the number of $(s_{1},\dots,s_{k})\in\{1,\dots,d-1\}^{k}$
  with $\sum_{i}s_{i}=dk/2$, and $\epsilon_{d}$ is $1$ for $d$ even and $0$ for
  $d$ odd; thus $T_{2}(k)=1$, $T_{3}(k)=\binom{k}{k/2}$, $T_{4}(6)=141$ and
  $T_{6}(6)=1751$.
\end{enumerate}
\end{theorem}

For odd $r$ there is nothing to prove, since every Hodge class of $X$ is
then a multiple of a power of the hyperplane class. So the theorem gives the
Hodge conjecture for the very general complete intersection of Vandermonde
type of diagonal quadrics in every dimension, of diagonal cubics or quartics
up to dimension seven, and of diagonal sextics up to dimension five, in
every $\PP^{N}$.

\begin{proof}
Apply \Cref{lem:vgmonodromy} to $H^{r}(X_{d,r}(\lambda),\ZZ)$ over $U$,
polarised by the hyperplane class, and let $\Pi'\subset\pi_{1}(U,\lambda)$ be
the subgroup of finite index that fixes the Hodge classes. The monodromy
commutes with $H$ and with $\Phi^{*}$, so $\Pi'$ fixes $x_{[a]}$ and each of
its parts in the spaces $\bigwedge^{r}V_{a'}$. Let $m$ be the order of $a$,
take representatives $0\le a_{i}<d$, and put $a_{i}'=ma_{i}/d$. The function
$y=\prod_{i}u_{i}^{a_{i}}$ is invariant under the kernel of $\chi_{a}$ and
satisfies $y^{m}=\prod_{i}(t-\lambda_{i})^{a_{i}'}$, so the quotient of $C$ by
that kernel is the cyclic cover $D_{a}$ of $\PP^{1}$ of degree $m$ given by
this equation, and pull-back identifies $H^{1}(D_{a},\QQ)$ with the part
$\bigoplus_{j\ne0}V_{ja}$ of $H^{1}(C,\QQ)$, compatibly with the monodromy.
The image $B_{[a]}$ is isogenous to the abelian subvariety of $J(D_{a})$
whose first cohomology is $\bigoplus_{j\in(\ZZ/m)^{\times}}V_{ja}$, all of
$J(D_{a})$ for $m\le3$, its polarisation being a rational multiple of the
restriction of the principal one, since pull-back along a map of degree $e$
multiplies the cup product by $e$.

(i) For $m=2$, $D_{a}$ is the hyperelliptic curve
$y^{2}=\prod_{i\in S}(t-\lambda_{i})$ over the support $S$ of $a$, of genus
$g=|S|/2-1$, and $H^{1}(D_{a},\QQ)=V_{a}$. Move the base point so that the
$\lambda_{i}$ with $i\in S$ are real and consecutive, $\lambda_{s_{1}}<\dots<
\lambda_{s_{2g+2}}$, the others lying away from the segment they span. The
loop in which $\lambda_{s_{k}}$ and $\lambda_{s_{k+1}}$ turn once around each
other, inside a disc that contains no other $\lambda_{i}$, makes $D_{a}$
acquire a node with vanishing cycle $c_{k}$, the lift of the segment
$[\lambda_{s_{k}},\lambda_{s_{k+1}}]$, and acts by the square of the
Picard--Lefschetz transvection,
$T_{k}^{2}(z)=z+2\langle z,c_{k}\rangle c_{k}$ up to orientation
\cite[Chapter~3]{Voi03}.%
\footnote{Near the collision the curve is $y^{2}=(t-\lambda_{s_{k}})
(t-\lambda_{s_{k+1}})\,v$ with $v$ a unit, analytically $y^{2}=t^{2}-\epsilon$
with $\epsilon$ a multiple of $(\lambda_{s_{k+1}}-\lambda_{s_{k}})^{2}$; one
turn of the two points around each other is two turns of $\epsilon$ around
$0$, and one turn
of $\epsilon$ acts by the Dehn twist along the vanishing cycle, whose action on
$H_{1}$ is the transvection $T_{k}$.} The classes $c_{1},\dots,c_{2g}$ satisfy
$\langle c_{k},c_{k+1}\rangle=\pm1$ and $\langle c_{k},c_{l}\rangle=0$ for
$|k-l|\ge2$, so their intersection matrix has determinant $1$ and they form a
basis of $H_{1}(D_{a},\QQ)$; they are the distinguished basis of type $A$ of
\cite{ACa79}. Some power $T_{k}^{2e}$ with $e\ge1$ lies in the image of
$\Pi'$; it is unipotent, so the Zariski closure of that image contains the
one-parameter group $\exp(\tau N_{k})$, $N_{k}(z)=\langle z,c_{k}\rangle c_{k}$,
and its Lie algebra contains $N_{1},\dots,N_{2g}$. These generate the Lie
algebra $\mathfrak{sp}(V_{a})$.%
\footnote{Identify $\mathfrak{sp}(V)$ with $\mathrm{Sym}^{2}V$, the product $uv$
acting by $z\mapsto\langle z,u\rangle v+\langle z,v\rangle u$, so that
$N_{k}=\tfrac12c_{k}^{2}$ and
$[uv,wz]=\langle u,w\rangle vz+\langle u,z\rangle vw+\langle v,w\rangle uz
+\langle v,z\rangle uw$ up to an overall sign. Then
$[c_{k}^{2},c_{k+1}^{2}]=\pm4c_{k}c_{k+1}$, and if all products
$c_{i}c_{j}$ with $i,j\le k$ lie in the generated algebra, so do the
$c_{i}c_{k+1}=\pm\tfrac12[c_{i}c_{k},c_{k+1}^{2}]$ for $i<k$, because
$c_{k+1}$ meets only $c_{k}$ among $c_{1},\dots,c_{k}$. By induction every
$c_{i}c_{j}$ lies in it, and these span $\mathrm{Sym}^{2}V$.} So the Zariski
closure contains $\Sp(V_{a})$, whose invariants in
$\bigwedge^{r}V_{a}$ are the multiples of $\omega_{a}^{r/2}$, $\omega_{a}$ the
polarisation form \cite[Theorem~17.5]{FH91}.%
\footnote{The exterior algebra of $V$ is a representation of $\Sp(V)$
and of the $\mathfrak{sl}_{2}$ spanned by $\omega\wedge{}$ and its adjoint,
and $\bigwedge^{r}V=\bigoplus_{j}\omega^{j}\wedge P^{r-2j}$, where $P^{k}$, the
kernel of the contraction with $\omega$ in degree $k\le g$, is irreducible
and nontrivial for $k\ge1$, so the invariants are the multiples of
$\omega^{r/2}$, for $r$ even.} The correspondence of the proof of
\Cref{thm:vandermonde}(iv) maps $H^{r}(B_{[a]},\QQ)=\bigwedge^{r}H^{1}(B_{[a]},
\QQ)$ isomorphically and equivariantly onto $\bigwedge^{r}V_{a}$; the class
$\theta_{a}$ goes to an invariant nondegenerate form, a nonzero multiple of
$\omega_{a}$, and $\theta_{a}^{r/2}$ to a nonzero multiple of
$\omega_{a}^{r/2}$, which is algebraic.

(ii) For $m\in\{3,4,6\}$, $D_{a}$ is the cyclic cover of \Cref{setup:cyclic}
for the exponents $a_{i}'\in\{1,\dots,m-1\}$ on the support $S$ of $a$, and
$V_{a}$ is its eigenspace, of dimension $n=|S|-2$ (\Cref{lem:foxmodel}). The
map from $U$ to $\mathrm{Conf}_{|S|}$ that forgets the $\lambda_{i}$ outside
$S$ is a fibration with connected fibres, so it is onto on fundamental groups,
and $\Pi'$ maps onto a subgroup of finite index of $P_{|S|}$; the Zariski
closure of its image in $\GL(V_{a})$ therefore contains $G_{a}^{0}$. If
$n\le2$, then $\bigwedge^{r}V_{a}$ is zero unless $r=n=2$. If $V_{a}^{1,0}$ or
$V_{a}^{0,1}$ is zero, then $\bigwedge^{r}V_{a}$ has type $(r,0)$ or $(0,r)$
and contains no class of type $(r/2,r/2)$. Otherwise $n\ge3$, and
\Cref{prop:cyclicmonodromy} gives $G_{a}^{0}\supseteq\SL(V_{a})$, whose
invariants in $\bigwedge^{r}V_{a}$ vanish for $0<r<n$. In every case
$x_{[a]}\ne0$ forces $r=n$, and then $\bigwedge^{r}V_{a}$ is a line of type
$(p_{a},r-p_{a})$ with $p_{a}=\dim V_{a}^{1,0}$, which contains a class of type
$(r/2,r/2)$ only if $p_{a}=r/2$. The field $\QQ(\zeta_{m})$ then acts on
$B_{[a]}$, of dimension $|[a]|\dim V_{a}/2=r$, with multiplicities
$(r/2,r/2)$ on the tangent space and compatibly with the polarisation, the
automorphism of $D_{a}$ preserving the cup product; so $B_{[a]}$ is of Weil
type (\Cref{def:weiltype}), and
$W_{[a]}=\bigwedge^{r}V_{a}\oplus\bigwedge^{r}V_{-a}$ is the image of
$\bigwedge^{r}_{\QQ(\zeta_{m})}H^{1}(B_{[a]},\QQ)$, its space of Weil classes.
Its discriminant (\Cref{def:disc}) is that of the lattice $L_{a}$ of
\Cref{lem:newdisc}: the isogeny with the abelian subvariety of $J(D_{a})$
whose first cohomology is $L_{a}\otimes\QQ$ scales the polarisation by a
rational number and $\det H$ by its $r$-th power, a norm. The coordinates of
$a$ equal to $d/2$ are those with $a_{i}'=3$ when $m=6$, and
\Cref{lem:newdisc} gives the assertion.

(iii) Let $x$ be a Hodge class of degree $r$. The component of $\Phi^{*}x$
along $\eta$ is algebraic. For $d\in\{2,3,4,6\}$ every $a\ne0$ has order $2$,
$3$, $4$ or $6$, and by (i) and (ii) the components $x_{[a]}$ are images of
powers of polarisations or of Weil classes of abelian $r$-folds of Weil type
for $\QQ(\sqrt{-3})$ or $\QQ(i)$. These are algebraic for $r=2$ by the theorem
of Lefschetz, for $r=4$ by Markman's theorem on abelian fourfolds of Weil
type, and for $r=6$ by Markman's theorem on abelian sixfolds of split Weil
type \cite{Mar25a,Mar25b} when the discriminant is trivial: always for
$d=3,4$, and for $d=6$ except along the characters of order six with an odd
number of coordinates $3$. So $\Phi^{*}x$ is
algebraic in the cases stated, and so is $x=\Phi_{*}\Phi^{*}x/|G|$. For the
dimension, the classes found are Hodge at every $\lambda$, and by (i) and (ii)
there are no others. Each character of order two with $|S|\ge r+2$ gives one
class, the image of $\theta_{a}^{r/2}$, nonzero since $r\le\dim V_{a}$. A
character of order at least three gives a class only if $|S|=r+2$ and
$\dim V_{a}^{1,0}=\sum_{i\in S}(d-a_{i})/d-1=r/2$, that is
$\sum_{i}a_{i}=d(r+2)/2$ with representatives $0\le a_{i}<d$, and then exactly
one, in $\bigwedge^{r}V_{a}$. The characters with $|S|=r+2$ whose coordinates
on $S$ average $d/2$ are counted by $T_{d}(r+2)$; for $d$ even they include
the one of order two with all coordinates $d/2$, and the characters of order
two with $|S|\ge r+4$ give the last sum.
\end{proof}

\begin{figure}[tbp]
\centering
\includegraphics[width=\textwidth]{fig_vgmonodromy}
\caption{The proof of \Cref{thm:vgvandermonde}. Left: the hyperelliptic
curve $D_{a}$ of a character of order two, drawn over the six points of its
support for $g=2$, with the chain of vanishing cycles $c_{1},\dots,c_{5}$;
a full twist of two neighbouring branch points, drawn as a braid, acts by
the square of the transvection $T_{k}$. Middle: the chain as a graph, and the
induction in $\mathrm{Sym}^{2}V_{a}$ by which the logarithms $N_{k}$ generate
$\mathfrak{sp}(V_{a})$. Right: what the monodromy leaves of the Hodge
classes, the powers of the polarisation for characters of order two, and
for characters of order three, four and six the Weil classes of an abelian
variety of Weil type for $\QQ(\sqrt{-3})$ or $\QQ(i)$.}
\label{fig:vgmonodromy}
\end{figure}

\Cref{fig:vgmonodromy} draws the steps of the proof. The dimensions of (iii)
never exceed $h^{r/2,r/2}(X)$, and they equal it for a quadric, for
the Fermat hypersurfaces (for cubics $7$, $21$ and $71$ in dimensions $2$, $4$
and $6$, for quartics $142$ and $1108$ in dimensions $4$ and $6$, for sextics
$1752$ in dimension $4$) and for two quadrics in $\PP^{N}$ with $N$ even
($N+2$), where every class of type $(r/2,r/2)$ is algebraic. For two cubics
in $\PP^{6}$ the very general member has $141$ Hodge classes of degree four
against $h^{2,2}=267$, and for three quartics in $\PP^{7}$ it has $3950$
against $h^{2,2}=29872$ (\Cref{fig:vgcounts}). \Cref{fig:vgscope} shows, for each degree $d$ and dimension $r$,
what is known for the very general member.

\begin{figure}[tbp]
\centering
\includegraphics[width=\textwidth]{fig_vgcounts}
\caption{The dimension of the space of Hodge classes of degree $r$ on the
very general member, from \Cref{thm:vgvandermonde}(iii), against the Hodge
number $h^{r/2,r/2}$, on a logarithmic scale. Where the two agree, the
space $H^{r/2,r/2}$ is spanned by Hodge classes at every configuration;
elsewhere the very general member has fewer Hodge classes, and the members
with more lie over a countable union of proper closed subsets of the space
of configurations.}
\label{fig:vgcounts}
\end{figure}

\begin{figure}[tbp]
\centering
\includegraphics[width=\textwidth]{fig_vgscope}
\caption{The Hodge conjecture for the very general complete intersection
$X_{d,r}(\lambda)$ of Vandermonde type, for every $N$. Grass: proved, by the
theorem of Lefschetz for $r=2$ and by \Cref{thm:vgvandermonde} for $d=2$,
for $d=3,4$ with $r\le6$ and for $d=6$ with $r\le4$. Teal hatching: proved
except for the Weil classes of abelian sixfolds of non-split Weil type for
$\QQ(\sqrt{-3})$. Clay: open, the Weil classes of abelian varieties of Weil
type over $\QQ(i)$ or $\QQ(\sqrt{-3})$ in dimension at least eight, or,
hatched, of a CM field of degree at least four (\Cref{rem:vandermonde}).
For odd $r$ there is nothing to prove.}
\label{fig:vgscope}
\end{figure}

\begin{remark}[What the diagonal complete intersections give]
\label{rem:vandermonde}
\Cref{thm:vandermonde} adds to the domain of \textup{(F3$'$)} a family of
complete intersections of every dimension and every degree, parametrised by
configurations of $N+1$ points on the line, among them all smooth complete intersections of two diagonal
hypersurfaces; for $c\ge3$ and $r\ge2$ the condition on the columns is a real
restriction. \Cref{cor:twodiagonal} shows on them how \textup{(F3$'$)} and
(F2) together give the conjecture: once the Hodge classes of $X$ are
carried to the abelian varieties $B_{[a]}$, the conjecture for $X$ is the
algebraicity of their Hodge classes, a question of type (F2), and for the
fourfolds in $\PP^{6}$ of degrees $3$, $4$ and $6$ it is answered by
Markman's theorem: their exceptional classes are Weil classes of abelian
fourfolds of Weil type, $490$ families of them for $d=4$. For $d=3$ these
fourfolds are Fano and the conjecture also follows from the theorem of Conte
and Murre on uniruled fourfolds \cite{CM78}; for $d=4$ and $d=6$ they are of
general type, with $h^{2,2}=2584$ and $48588$. Complete intersections of
diagonal hypersurfaces, and of quadrics, are also the subject of Terasoma's
paper \cite{Ter88}. At a very general configuration
\Cref{thm:vgvandermonde} goes further, in every $\PP^{N}$: for $d=2$ in every
dimension, for $d=3$ and $d=4$ up to dimension seven, and for $d=6$ up to
dimension five, through the monodromy of cyclic covers of degree three, four
and six (\Cref{prop:cyclicmonodromy}). For $d=3$ and $r\ge8$ the
conjecture for the very general member is the algebraicity of the Weil
classes of the abelian $r$-folds $B_{[a]}$ of split Weil type for
$\QQ(\sqrt{-3})$, which is open from dimension eight on, and the same holds
for $d=4$ with $\QQ(i)$. For $d=6$ and $r=6$ the open classes are those of
the abelian sixfolds of non-split Weil type for $\QQ(\sqrt{-3})$ of
\Cref{thm:vgvandermonde}(iii), such as the one of the character
$(1,1,2,2,3,5,5,5)$ on eight points. Beyond these cases the
same reduction meets (F2) itself: for $d=5$ or $d\ge7$ and $N$ large enough, some
$B_{[a]}$ with $\dim V_{a}=r$ carries an action of $\QQ(\zeta_{e})$, $e\mid d$,
with $\varphi(e)\ge4$, and $W_{[a]}$ is the top exterior power of its first
cohomology over that CM field, which consists of Hodge classes when the
types are balanced: Weil classes for a CM field of degree at least four, of
the kind treated in \Cref{thm:main}; and for larger $N$ the $B_{[a]}$ have
dimension above five.
\end{remark}

\begin{remark}[Where \textup{(F3$'$)} stops]\label{rem:f3primefrontier}
With \Cref{prop:f3prime}, \Cref{prop:aclosure} and \Cref{prop:f3primesmall},
\textup{(F3$'$)} holds on every variety of dimension at most three, on every
fourfold whose zero-cycles are supported in dimension three (in particular
every uniruled fourfold), on every fivefold whose zero-cycles are supported
in dimension at most three, on every fourfold and fivefold whose maximal
rationally connected quotient is birational to a variety satisfying
\textup{(F3$'$)} in degree four (\Cref{prop:f3primefibration}), on every
fourfold and fivefold dominated by a rational map from such a variety
(\Cref{prop:f3primedominant}), and on
every member of $\mathcal{A}$, a class containing the Fermat
hypersurfaces, the smooth Delsarte hypersurfaces with the cyclic covers
branched along them (\Cref{cor:delsarteall}), and the complete intersections
of diagonal hypersurfaces of Vandermonde type, among them every smooth
complete intersection of two diagonal hypersurfaces (\Cref{thm:vandermonde},
\Cref{cor:twodiagonal}), of every degree and dimension, the Hilbert schemes of points on
abelian surfaces, on products of curves and on the K3 surfaces of
\Cref{prop:aclosure}(iv), the generalised Kummer varieties $K_{n}(A)$ and the
other varieties of \Cref{prop:aclosure}(v), and every blow-up and projective bundle
built from these. Through the Hodge conjecture itself it holds on the
complete intersections of \Cref{cor:twodiagonal}(ii), among them the
fourfolds of general type cut out by two diagonal quartics or sextics in
$\PP^{6}$, on the very general complete intersections of diagonal quadrics,
of diagonal cubics and quartics up to dimension seven and of diagonal
sextics up to dimension five, of \Cref{thm:vgvandermonde}, and on the
varieties of \Cref{prop:k3powers}, \Cref{prop:k3ntype} and
\Cref{cor:fanolines}: for a K3 surface $S$ with
$E=\End_{\mathrm{Hdg}}(T(S))$ a CM field, on every power, Hilbert scheme and
moduli space of $S$; for $E=\QQ$ and $t=\dim T(S)$, on $S^{k}$ for $k<t$, on
$S^{[n]}$ for $n<t(t+1)/2$ and on the moduli spaces of dimension less than
$t$; on every fourfold of K3$^{[2]}$ type whose field $E$ is $\QQ$ or a CM
field; on $F(Y)$ and $Y\times Y$ for every cubic fourfold $Y$ with such a
field; and on every fourfold of generalised Kummer type \cite{FV25}. Its
frontier cases are these (\Cref{fig:f3primereach}), the first two in
degree four on fourfolds whose zero-cycles are not supported on a proper
closed subset and which are not birational to a member of $\mathcal{A}$ by
any construction above: $S\times S$ and $S^{[2]}$ for a K3 surface
whose real multiplication is not induced by a cycle constructed so far, such
as $S_{\lambda}$
(the known cases are recalled in \Cref{rem:mumfordks}); the fourfolds of
K3$^{[2]}$ type with real multiplication, among them $F(Y)$ for a cubic
fourfold $Y$ with real multiplication, and with it $Y\times Y$; and, for a
very general K3 surface, the single class $\det T(S)$ on $S^{21}$ of
\Cref{prop:orthpowers}, equivalently the Hodge conjecture for $S^{[231]}$.
In the first two the
rational Hodge isometries \cite{Bus19,Mar24} act on $T$ by $\pm1$, and a
composite of Hodge similarities that returns to $T$ is an element $g\in E$
with $g^{2}\in\QQ$, the adjoint involution of $E$ being trivial; such
elements span only the compositum of the quadratic subfields of $E$, which is
$\QQ$ for the cubic field of $S_{\lambda}$. Apart from proving the
conjecture itself where zero-cycles are small, none of these results touches
the step that
\Cref{rem:f3primeopen} identifies: putting a Hodge class into
$\mathrm{Ab}^{p}(X)$ for a variety that no correspondence connects to an
abelian variety. \Cref{thm:closure} is unchanged, since every statement
proved here is a special case of one of the hypotheses of its rule set, not a
new rule. With the theorems of recent preprints as hypotheses, the first and
third frontier cases are settled, and the second when the N\'eron--Severi
space is isotropic or represents $-2$, for instance when the Picard number is
at least four (\Cref{cor:k3all}, \Cref{prop:k3twosmall}).
\end{remark}

\begin{figure}[tbp]
\centering
  \includegraphics[width=\textwidth]{fig_f3primereach}
\caption{What arguments on zero-cycles and on rationally connected
fibrations give for \textup{(F3$'$)} and for the Hodge conjecture (HC in the
figure) on a variety $X$, over the plane of $n=\dim X$ and the least dimension
$d$ of a closed subset supporting $\mathrm{CH}_{0}(X)_{\QQ}$; $d$ is at most
the dimension of the maximal rationally connected quotient $R$ of $X$, and
$d=n$ whenever $H^{n,0}(X)\ne0$ (\Cref{rem:f3primesharp}). Both hold when
$n\le5$ and $d\le3$ (\Cref{prop:f3primesmall}), hence on every variety of
dimension at most three. In the hatched cells
\textup{(F3$'$)} holds when $R$ is birational to a variety with
$\Hdg^{2}=\mathrm{Ab}^{2}$, for instance a member of $\mathcal{A}$, and the
Hodge conjecture when $R$ is birational to an abelian fourfold
(\Cref{prop:f3primefibration}). The double arrows are the equivalences of
\Cref{rem:f3primesharp}: \textup{(F3$'$)}, and the Hodge conjecture, hold for
every uniruled fivefold, and for every rational sixfold, if and only if they
hold in degree four for every fourfold; the sixfold $\mathrm{Bl}_{Y}\PP^{6}$
of a smooth sextic fourfold $Y$ in a hyperplane carries the Hodge numbers
listed. The rings are the frontier cases (\Cref{rem:f3primefrontier}): the
sextic fourfolds, the first family of smooth hypersurfaces beyond these
results
(the $29$ Delsarte sextics lie in $\mathcal{A}$ by
\Cref{cor:delsarteall}); $S\times S$ and $S^{[2]}$ for a K3
surface $S$ with real multiplication, such as
$S_{\lambda}$; the fourfolds of K3$^{[2]}$ type with real multiplication,
among them $F(Y)$; and, beyond the breaks of the axes, $S^{[231]}$ for a very
general K3 surface, which is equivalent to the algebraicity of $\det T(S)$ on
$S^{21}$ (\Cref{prop:k3powers}, \Cref{fig:k3threshold}).}
\label{fig:f3primereach}
\end{figure}

\begin{setupenv}[The rule set]\label{setup:rules}
The rule set has forty-two statements and twenty-two rules. Ten statements
are proved in this paper: the base point with an explicit cycle in every Weil
family of an imaginary quadratic field (\Cref{thm:everyfamily}) and of a CM
field of degree at least four (\Cref{thm:cmbasepoint}); for each of the two
kinds of field, the reduction of the Weil classes of a family to one class at
every point of one connected domain, and the stability of the algebraic locus
under the rational points, with dense orbits (\Cref{sec:reduction},
\Cref{prop:densealg}, \Cref{thm:cmpropagation}); and four statements on
secant objects: the form of the weakened criterion for the transform
(\Cref{thm:factor}), the conditions of Markman's strategy that bear on the
Chern character, the description of the local obstruction as a failure of
Cohen--Macaulayness (\Cref{cor:cmdichotomy}), and three of the four
ingredients of the construction in every dimension. Four are quoted:
$\mathbf{W}(K,2,\delta)$ for every $K$ and $\delta$, $\mathbf{W}(K,3,\delta_{0})$
\cite{Mar25a,Mar25b}, the deformation theorem for motivated classes
\cite[Th\'eor\`eme~0.5]{And96b}, and the accessibility of Hodge classes on
abelian varieties \cite[\S5.1, Theorem~3]{Mil20}. Sixteen are hypotheses:
(F3), (F2) of \Cref{cor:fourremain} and (F3$'$) of \Cref{def:f3prime};
(L), (M), (V) of \Cref{def:lmv}; (P2), in the polarised form of
\Cref{rem:p2prime}, for four groups of families, the split and the other
families of the imaginary quadratic fields and of the CM fields of degree at
least four, the third group giving $\mathrm{P2}_{\mathrm{split}}$; secant
objects meeting the weakened criterion in every dimension, and an affirmative
answer to \cite[Question~11.4]{Mar25b}; and, for a smooth and for a singular
Cohen--Macaulay support on an abelian fourfold, its existence with the
required extension property and the vanishing of the image of the product of
translation classes. The other twelve are derived: $\mathbf{HC}$, the
conjecture for every abelian variety, written $\mathbf{AB}$, and the Weil
classes of every CM field, of the imaginary quadratic fields, of their
families of trivial and of nontrivial discriminant, of the CM fields of
degree at least four and of their split and their other families, together
with one secant object on an abelian fourfold, $\mathbf{W}(K,4,\delta_{0})$,
and the range settled in the literature. Each rule is a finite set of
premises with one conclusion, and each is a theorem of this paper:
\Cref{tab:rules} lists twenty-one of them, and the twenty-second derives the
range settled in the literature from the two quoted Weil statements. For a
set $T$ of statements write $\Cn(T)$ for the least set that contains $T$ and
is closed under the rules, which exists because each rule has a single
conclusion and the consequence operator is therefore monotone. Write
$\mathbf{T}$ for the set of statements proved or quoted. Call a set $S$ of
hypotheses \emph{sufficient} if $\mathbf{HC}\in\Cn(\mathbf{T}\cup S)$, and
\emph{minimal} if no proper subset of $S$ is sufficient.
\end{setupenv}

\begin{table}[htbp]
\centering
\small
\renewcommand{\arraystretch}{1.14}
\begin{tabular}{@{}>{\raggedright\arraybackslash}p{0.40\linewidth}>{\raggedright\arraybackslash}p{0.27\linewidth}l@{}}
\toprule
premises & conclusion & by \\
\midrule
the conjecture for abelian varieties, (F3)
  & $\mathbf{HC}$ & \Cref{ssec:closuremap} \\
(F3)
  & abelian varieties & \Cref{prop:f3ishc} \\
the conjecture for abelian varieties, (F3$'$)
  & $\mathbf{HC}$ & \Cref{prop:f3prime} \\
the Weil classes of every CM field, (F2)
  & abelian varieties & \Cref{ssec:closuremap} \\
$K$ imaginary quadratic, the fields of degree at least four
  & every CM field & \Cref{app:albert} \\
$\delta=\delta_{0}$, $\delta\ne\delta_{0}$
  & every discriminant & \Cref{prop:separate} \\
a base point, the reduction, the orbit, (P2) for the family
  & that family & \Cref{thm:semiregclosure} \\
the same three for a CM field of degree at least four, (P2) for the family
  & that family & \Cref{thm:cmpropagation} \\
its split families, its other families
  & the fields of degree at least four & \Cref{prop:cmweil} \\
$\delta=\delta_{0}$
  & $\delta\ne\delta_{0}$ & \Cref{prop:descent} \\
the split families of the fields of degree at least four
  & their other families; $\delta=\delta_{0}$ & \Cref{prop:descent} \\
a secant object in every dimension, Question 11.4
  & $\delta=\delta_{0}$ & \Cref{prop:splitgeom} \\
a Cohen--Macaulay support, smooth or singular
  & one secant object & \Cref{cor:cmdichotomy} \\
one secant object, Question 11.4
  & $\mathbf{W}(K,4,\delta_{0})$ & \Cref{cor:markmancondition} \\
\midrule
(L), (M)
  & $\mathbf{HC}$ & \Cref{prop:lefmot} \\
(L), the deformation theorem for motivated classes
  & (V) & \Cref{prop:lefvhc} \\
(V), the accessibility of Hodge classes on abelian varieties
  & abelian varieties & \Cref{thm:vhcab} \\
\bottomrule
\end{tabular}
\caption{The rules of \Cref{setup:rules} whose conclusion is derived, with the
theorem that proves each; the rows for (P2) and for the Cohen--Macaulay
support each stand for two rules, the row for descent over the fields of
degree at least four for two, descent within a field and scalar extension
down to the imaginary quadratic fields, the second row is \Cref{prop:f3ishc},
the third is
\Cref{prop:f3prime}(i) for the weaker statement (F3$'$) of
\Cref{def:f3prime}, the last three rows are the rules of
\Cref{ssec:bullseye}, one of whose conclusions, (V), is itself a
hypothesis of the rule set, and the one remaining rule records what the literature supplies on Weil classes. The
bounded criterion (P1) is absent: by
\Cref{thm:p1equivalent} it is equivalent to its own conclusion, so it is a
restatement and not a premise.}
\label{tab:rules}
\end{table}

\begin{theorem}[The closure of what is proved here]\label{thm:closure}
With the notation of \Cref{setup:rules}:
\begin{enumerate}[label=\textup{(\roman*)},leftmargin=2.6em]
\item $\mathbf{HC}\notin\Cn(\mathbf{T})$. Nothing in this paper, combined with the
  theorems it quotes, implies the Hodge conjecture.
\item There are exactly five minimal sufficient sets:
  \[
    \{\mathrm{F3}\},\quad\{\mathrm{L},\mathrm{M}\},\quad
    \{\mathrm{L},\mathrm{F3}'\},\quad\{\mathrm{V},\mathrm{F3}'\},\quad
    \{\mathrm{P2}_{\mathrm{split}},\mathrm{F2},\mathrm{F3}'\},
  \]
  with \textup{(L)}, \textup{(M)}, \textup{(V)} as in \Cref{def:lmv},
  \textup{(F2)}, \textup{(F3)} as in \Cref{cor:fourremain}, (F3$'$) as in
  \Cref{def:f3prime}, and $\mathrm{P2}_{\mathrm{split}}$ the propagation
  statement \textup{(P2)}, read in the polarised form of \Cref{rem:p2prime},
  for the split families of the CM fields of degree at least four. No
  minimal set contains \textup{(P2)} for an imaginary quadratic field or for
  a family of nontrivial discriminant: \Cref{prop:descent} derives those.
  The single hypothesis (F3) is sufficient, and it is equivalent to
  $\mathbf{HC}$ (\Cref{prop:f3ishc}); no other single hypothesis is.
  The sets $\{\mathrm{L},\mathrm{F3}\}$, $\{\mathrm{V},\mathrm{F3}\}$ and
  $\{\mathrm{P2}_{\mathrm{split}},\mathrm{F2},\mathrm{F3}\}$ are sufficient and
  not minimal. The last of the five is the only one with three elements and the
  only one whose derivation of $\mathbf{HC}$ uses a statement proved in this
  paper.
\item In each of the five, every element is necessary: removing any one
  leaves $\mathbf{HC}$ outside the closure.
\item The secant route lies in no minimal sufficient set. Granting both of its
  demands, a secant object satisfying the weakened criterion in every
  dimension and an affirmative answer to \cite[Question 11.4]{Mar25b}, puts
  $\mathbf{W}(K,n,\delta_{0})$ in the closure, and with it, by
  \Cref{prop:descent}, $\mathbf{W}(K,n,\delta)$ for every $\delta$; it leaves
  the Weil classes of the CM fields of degree at least four outside.
\item No statement proved in this paper rests, at any depth, on a
  hypothesis, and the rule set is acyclic.
\item Every element of a minimal sufficient set other than
  $\mathrm{P2}_{\mathrm{split}}$ is a consequence of $\mathbf{HC}$, and $\{\mathrm{F3}\}$ and
  $\{\mathrm{L},\mathrm{M}\}$ are each equivalent to it; no implication from
  $\mathbf{HC}$ to \textup{(P2)} is known.
\end{enumerate}
\end{theorem}

\begin{proof}
(i) and (v). No rule has a proved or quoted statement as its conclusion, so
each of them lies in $\Cn(\mathbf{T})$ by itself, on the strength of its proof
in this paper or in the literature; this is the first assertion of (v). The
twenty-second rule is the only one whose premises all lie in $\mathbf{T}$:
every other rule has among its premises a hypothesis or a derived statement
other than the range settled in the literature. Hence $\Cn(\mathbf{T})$ is
$\mathbf{T}$ together with that range, and it does not contain $\mathbf{HC}$,
which is (i). For acyclicity, list first the thirty statements that are
proved, quoted or hypotheses other than (V); then one secant object,
$\mathbf{W}(K,4,\delta_{0})$ and the range settled in the literature; the
split and then the other families of the CM fields of degree at least four;
the families of trivial and then of nontrivial discriminant of the imaginary
quadratic fields; the imaginary quadratic fields, the CM fields of degree at
least four and every CM field; and last (V), $\mathbf{AB}$ and $\mathbf{HC}$.
The premises of every rule come before its conclusion in this list.

(ii) and (iii). The five sets are sufficient:
$\mathrm{F3}\Rightarrow\mathbf{AB}$ (\Cref{prop:f3ishc}) and
$\mathbf{AB}\wedge\mathrm{F3}\Rightarrow\mathbf{HC}$; $\mathrm{L}\wedge
\mathrm{M}\Rightarrow\mathbf{HC}$ (\Cref{prop:lefmot}); (L) and the
deformation theorem give (V) (\Cref{prop:lefvhc}), (V) and accessibility give
$\mathbf{AB}$ (\Cref{thm:vhcab}), and $\mathbf{AB}\wedge\mathrm{F3}'
\Rightarrow\mathbf{HC}$ (\Cref{prop:f3prime}(i)), which settles
$\{\mathrm{L},\mathrm{F3}'\}$ and $\{\mathrm{V},\mathrm{F3}'\}$; and
$\mathrm{P2}_{\mathrm{split}}$ with the CM base point, the reduction and the
orbit gives the split families of the CM fields of degree at least four
(\Cref{thm:cmpropagation}), descent gives their other families and the
families of trivial discriminant of the imaginary quadratic fields
(\Cref{prop:descent}), hence the fields of degree at least four
(\Cref{prop:cmweil}), the nontrivial discriminant (\Cref{prop:descent}),
the imaginary quadratic fields (\Cref{prop:separate}) and every CM field
(\Cref{app:albert}), and with (F2) this is $\mathbf{AB}$ and with (F3$'$)
it is $\mathbf{HC}$.

Let $C_{1}$, respectively $C_{2}$, be the set of all statements other than
$\mathbf{HC}$, (F3), (F3$'$) and (L), respectively (M). Let $C_{3}$ be the
set of all statements other than $\mathbf{HC}$, $\mathbf{AB}$, the Weil
classes of every CM field, of the fields of degree at least four and of their
split families, (F3), (L), (V) and $\mathrm{P2}_{\mathrm{split}}$; and let
$C_{4}$ be the set of all statements other than $\mathbf{HC}$,
$\mathbf{AB}$, (F3), (L), (V) and (F2). Each $C_{j}$ contains $\mathbf{T}$
and not $\mathbf{HC}$, and each is closed under the rules: every rule whose
conclusion it omits has a premise it omits. For $C_{1}$ and $C_{2}$ the only
omitted conclusion is $\mathbf{HC}$, and the three rules with that
conclusion have the premises (F3), (F3$'$), and (L) and (M). For $C_{3}$
and $C_{4}$ the rules with conclusion $\mathbf{HC}$ have the premises
$\mathbf{AB}$ or (L); those with conclusion $\mathbf{AB}$ have the premises
(F3), the Weil classes of every CM field together with (F2), and (V); the
only rule with conclusion (V) has the premise (L); and for $C_{3}$ the rule
for every CM field needs the fields of degree at least four, the rule for
these needs their split families, and the rule for the split families needs
$\mathrm{P2}_{\mathrm{split}}$. Since $\Cn$ is monotone, a set of
hypotheses contained in some $C_{j}$ has its closure, together with
$\mathbf{T}$, inside $C_{j}$, and is not sufficient. So a sufficient set
meets each of
\begin{gather*}
  H_{1}=\{\mathrm{F3},\mathrm{F3}',\mathrm{L}\},\qquad
  H_{2}=\{\mathrm{F3},\mathrm{F3}',\mathrm{M}\},\\
  H_{3}=\{\mathrm{F3},\mathrm{L},\mathrm{V},\mathrm{P2}_{\mathrm{split}}\},\qquad
  H_{4}=\{\mathrm{F3},\mathrm{L},\mathrm{V},\mathrm{F2}\},
\end{gather*}
the hypotheses omitted from $C_{1},\dots,C_{4}$. Let $S$ meet all four. If
$\mathrm{F3}\in S$, then $S\supseteq\{\mathrm{F3}\}$. If $\mathrm{F3}\notin S$
and $\mathrm{L}\in S$, then $S$ meets $H_{2}$ in (F3$'$) or (M). If neither is
in $S$, then $S$ contains (F3$'$), from $H_{1}$, and (V) or both
$\mathrm{P2}_{\mathrm{split}}$ and (F2), from $H_{3}$ and $H_{4}$. So a set
of hypotheses is sufficient exactly when it contains one of the five sets,
and as none of the five contains another, they are the minimal sufficient
sets. The other assertions of (ii) can be read off the list, and the
derivations above show that only the last uses a statement proved here;
(iii) is the minimality.

(iv) Granting the two demands of the secant route, the rule of
\Cref{prop:splitgeom} gives the families of trivial discriminant, descent
those of nontrivial discriminant, and \Cref{prop:separate} the imaginary
quadratic fields. The set formed by $\mathbf{T}$, the range settled in the
literature, the two demands and these three statements is closed, since
every other rule has a premise outside it: a hypothesis not granted, one
secant object, or a Weil statement about the fields of degree at least four
or about every CM field. Every base point the secant route produces is a
member $X\times\widehat{X}$, whose discriminant is trivial whatever $\beta$
is chosen (\Cref{prop:splitgeom}(ii)); the deformation argument stays inside
the family of trivial discriminant, and it is \Cref{prop:descent}(i), not
the deformation, that carries the result to the other discriminants. The
route has no base point over a field of degree at least four, so it reaches
none of them, and none of its hypotheses lies in one of the five sets.

(vi) \textup{(L)} follows from $\mathbf{HC}$ by \Cref{thm:equivalence},
\textup{(M)} by \Cref{prop:lefmot}, \textup{(V)} by \Cref{prop:lefvhc},
(F3$'$) by \Cref{prop:f3prime}(i), and (F2) and (F3) are special cases of
$\mathbf{HC}$. The equivalence of $\{\mathrm{F3}\}$ with $\mathbf{HC}$ is
\Cref{prop:f3ishc}, and $\{\mathrm{L},\mathrm{M}\}$ is equivalent to it by
\Cref{prop:lefmot} and the first two implications.
\end{proof}

\begin{corollary}[The three statements, one of which is the
conjecture]\label{cor:fourremain}
The Hodge conjecture follows from the results of this paper together with the
following three statements. No rule of \Cref{setup:rules} derives it from
(F1) and (F2) together, and it follows from (F3) alone, which is equivalent
to it
(\Cref{prop:f3ishc}); so the list is not a reduction of the conjecture to
three smaller problems. With (F3) replaced by the weaker statement
\textup{(F3$'$)} of \Cref{def:f3prime}, the conjecture modulo abelian
varieties, the three statements (F1), (F2), \textup{(F3$'$)} still suffice,
and no rule of \Cref{setup:rules} derives the conjecture from a proper
subset of them. With the rule of \Cref{prop:f2isab}, which the rule set does
not contain, (F2) and \textup{(F3$'$)} already suffice (\Cref{rem:f2rule}).
\begin{enumerate}[label=\textup{(F\arabic*)},leftmargin=3.0em,itemsep=4pt]
\item \emph{One propagation statement, for every CM field.} \textup{(P2)},
  in the polarised form of \Cref{rem:p2prime}: at one base point of each
  Weil family of each CM field some perfect complex whose Chern character is
  a nonzero multiple of a rational Weil class plus a polynomial in the
  polarisation has injective semiregularity map. Its first form, with Chern
  character exactly a multiple of the Weil class, is false for every
  imaginary quadratic field and every $n\ge2$ (\Cref{thm:p2false}): such a
  complex would have to deform to non-algebraic Weil tori, whose coherent
  sheaves have no Chern character in degrees $1$ to $2n-1$. The polarised form
  closes
  $\mathbf{W}(K,n,\delta)$
  for every imaginary quadratic $K$, every $n$ and every $\delta$ at once, by
  \Cref{thm:semiregclosure}, and
  $\mathbf{W}(F,n,\delta)$ for every CM field $F$ of higher degree, by
  \Cref{thm:cmpropagation}, because those theorems put $\Sigma=\cD$, so the
  Weil class is algebraic at every member of the family and not only at the
  general one, and because \Cref{thm:everyfamily} and \Cref{thm:cmbasepoint}
  supply the base point in every one of those families. By
  \Cref{prop:descent} it is enough to have it for the split families of the
  CM fields of degree at least four, in infinitely many dimensions.
  \Cref{thm:p2forces}, \Cref{thm:p2support}, \Cref{thm:nosum},
  \Cref{thm:p2numerical} and \Cref{thm:p2endo} describe the object of the
  first form and are vacuous by \Cref{thm:p2false}; for the polarised form
  \Cref{thm:p2primenumber} computes the number that $\dim\Ext^{2}(E,E)$ has
  to reach, $7n^{2}-2n$ for a general shape. For one family its
  conclusion is equivalent to the Lefschetz standard conjecture for one
  variety (\Cref{cor:weilequiv}). The bounded criterion \textup{(P1)} is not a
  second form of the same demand: by
  \Cref{thm:p1equivalent} it is equivalent to the conclusion it was to imply,
  and \Cref{thm:p1forces} shows that the one cycle it could have been tested
  on, the transport of the base cycle, does not satisfy it.
\item \emph{The classes beyond the Weil lines.} Every Hodge class on an
  abelian variety that lies outside the subring generated by divisor classes
  and by Weil classes of quotients is algebraic. Through dimension five there
  is no such class \cite{MZ99}; in dimension eight there are
  (\Cref{prop:mumfordproduct}), so this is not a reduction of the conjecture
  for abelian varieties to Weil classes, which is false as a statement about
  Hodge rings: it is the conjecture for those classes, and by
  \Cref{prop:f2isab} it is equivalent to the conjecture for every abelian
  variety. Its smallest case is
  brought down to one Kuga--Satake class in \Cref{thm:mumfordks}; it holds at
  every CM point (\Cref{thm:mumfordcm}), and on all members of a Mumford
  family it is equivalent to the Lefschetz standard conjecture for one
  ninefold (\Cref{cor:mumfordequiv}). Its classes are rigid along the base
  curve (\Cref{thm:mumfordrigid}), and one complex at one CM point with
  $\Ext^{2}$ of dimension $119$ would close it (\Cref{thm:mumfordnumerical}),
  as would a bound on the cycles that represent them at the closed points of
  the reductions of the curve modulo primes (\Cref{thm:modpdensity}).
\item \emph{The varieties that are not abelian.} The Hodge conjecture holds
  for every smooth complex projective variety that is not an abelian variety.
  This is the Hodge conjecture itself. The varieties it covers include
  $A\times\PP^{1}$ for every abelian variety $A$, and the conjecture for
  $A\times\PP^{1}$ gives it for $A$ (\Cref{prop:f3ishc}); so (F3) implies
  the conjecture for abelian varieties, hence (F1) and (F2), and hence
  $\mathbf{HC}$. Nor is it a reduction to abelian varieties in the other
  direction: no realisation of the cohomology of an arbitrary variety inside
  that of abelian varieties exists (\Cref{prop:noabeliantype}), and
  \Cref{thm:hyperplane} and \Cref{thm:returntrip} close the two routes
  through hyperplane sections that would be the natural candidates for one.
\end{enumerate}
\end{corollary}

\begin{proof}
The three together are sufficient by the derivation of \Cref{thm:closure},
and (F3) alone by \Cref{prop:f3ishc}. That (F1) and (F2) together are not
sufficient in the rule set is \Cref{thm:closure}(ii), since every minimal
sufficient set contains (F3), \textup{(F3$'$)} or \textup{(M)}; the
statement about (F1), (F2), \textup{(F3$'$)} is \Cref{thm:closure}(ii) and
(iii) for the last of the five minimal sets. The Weil classes of the CM
fields of degree at least four are not a fourth statement: they are derived
from (F1) by \Cref{thm:cmpropagation}. Nor is the bounded criterion a
premise, since by \Cref{thm:p1equivalent} it is a restatement of its own
conclusion.
\end{proof}

\begin{remark}[Necessity relative to a rule set, and the wording of (F2)]
\label{rem:f3collapse}
A rule set can show that a set of statements suffices; it can show that a set
is needed only relative to the implications written into it. A rule set that
used (F3) only through $\text{(abelian varieties)}\wedge
\text{(F3)}\Rightarrow\mathbf{HC}$, without the elementary rule
$\text{(F3)}\Rightarrow\text{(abelian varieties)}$ of \Cref{prop:f3ishc},
would make the three statements of \Cref{cor:fourremain} look jointly
necessary, although (F3) alone suffices.

The wording of (F2) has the same feature, one degree at a time. Let
$R(A)\subseteq H^{*}(A,\QQ)$ be the subring generated by divisor classes and
Weil classes of quotients. Suppose some Hodge class $\beta$ of degree $2p$ on
$A$ lies outside $R(A)$, and let $\alpha$ be any Hodge class of degree $2p$.
If $\alpha\notin R(A)$, (F2) makes it algebraic. If $\alpha\in R(A)$, then
$\alpha+\beta\notin R(A)$, because $R(A)$ is closed under subtraction, so
(F2) makes $\alpha+\beta$ and $\beta$ algebraic, and with them $\alpha$.
So in every degree in which (F2) says anything it is the whole conjecture for
that degree: on the square of a Mumford fourfold it gives every Hodge class of
$H^{4}$, the six products of divisor classes included. Multiplying by an
exceptional class of that square moves every Hodge class of every abelian
variety outside the subring, and \Cref{prop:f2isab} deduces from this that
(F2) is the whole conjecture for abelian varieties. The rule set does not
contain that implication; \Cref{rem:f2rule} says how \Cref{thm:closure}(ii)
changes when it is adjoined.

The form of (F3) that is not the conjecture itself asks only for
algebraicity modulo classes coming from abelian varieties. It is
\textup{(F3$'$)} of \Cref{def:f3prime}; \Cref{prop:f3prime} says what is
known of it, and \Cref{rem:f3primeopen} why it is not closed.
\end{remark}

\begin{figure}[tbp]
\centering
  \includegraphics[width=0.92\textwidth]{fig_frontier}
\caption{The rule set of \Cref{setup:rules} with the rules of
\Cref{ssec:bullseye}, drawn as a hypergraph. A rule with several premises is
a junction marked $\wedge$, joined to its premises and with one arrow to its
conclusion. The letters name the results that prove the rules: (a) is
\Cref{prop:f3ishc}; (b) and (g) are the closure map of
\Cref{ssec:closuremap}; (c) is \Cref{prop:f3prime}(i); (d) is
\Cref{prop:lefmot}; (e) is \Cref{prop:lefvhc}; (f) is \Cref{thm:vhcab};
(h) is \Cref{thm:cmpropagation}; and (i) is the chain of
\Cref{prop:descent}, \Cref{prop:separate}, \Cref{prop:cmweil} and
\Cref{app:albert}. Clay marks a hypothesis, a double frame a hypothesis
that the conjecture implies, grass what this paper proves, slate
what it quotes, and indigo a derived statement. The table lists the five
minimal sufficient sets of \Cref{thm:closure}(ii), the rules each of them
uses, and which of the hypotheses the conjecture implies; only
$\mathrm{P2}_{\mathrm{split}}$, which asks for a particular complex, is
drawn without that frame. The secant route, which reaches the imaginary
quadratic families only, is drawn in \Cref{fig:closure}.}
\label{fig:frontier}
\end{figure}

\begin{remark}[The target, and what it would take]\label{rem:bullseye}
With the rules of \Cref{ssec:bullseye}, \Cref{prop:f3ishc} and
\Cref{prop:f3prime} in place, \Cref{thm:closure} changes the reading of
\Cref{cor:fourremain} in three ways. First, the route through this paper is
a route only with (F3) replaced by \textup{(F3$'$)}. As stated, its last
statement (F3) is the conjecture (\Cref{prop:f3ishc}), so (F1) and (F2), and
everything this paper proves about Weil classes, drop out of every minimal
derivation that uses (F3). With \textup{(F3$'$)} in its place the route
$\{\mathrm{P2}_{\mathrm{split}},\mathrm{F2},\mathrm{F3}'\}$ is minimal, and it is not the
shortest: (F1) and (F2), which between them give the conjecture for abelian
varieties along this route, are replaced on the routes of the literature by
the single statement \textup{(V)} (\Cref{thm:vhcab}), or by \textup{(L)},
which implies it (\Cref{prop:lefvhc}). Second, every element of every minimal sufficient set
except $\mathrm{P2}_{\mathrm{split}}$ is a consequence of the conjecture
(\Cref{thm:closure}(vi)), and
two of the sets, $\{\mathrm{F3}\}$ and $\{\mathrm{L},\mathrm{M}\}$, are
each equivalent to it. (P2) is a statement about semiregularity maps, and a
proof of the conjecture need not prove it; in its first form, with a Chern
character that is exactly a Weil class, it is false (\Cref{thm:p2false}), and
the route runs through the polarised form of \Cref{rem:p2prime}. Third, every minimal sufficient set
contains (F3), \textup{(F3$'$)} or \textup{(M)}: the varieties that are not
abelian are reached only through the conjecture for them, which is the whole
conjecture, through the conjecture modulo abelian varieties
(\Cref{rem:f3primeopen}), or through motivated classes.

So the target at the centre is \textup{(L)}. It lies in two of the five
minimal sets, it gives \textup{(V)} and with it the whole abelian half, the
conjecture implies it, and with \textup{(M)} it is the conjecture. For each
$X$ it asks for one explicit Hodge class on $X\times X$ to be algebraic, the
class inducing $\Lambda$. This paper proves \textup{(L)}, \textup{(M)},
\textup{(V)} and \textup{(F3$'$)} on the classes of varieties named here:
\Cref{cor:mumfordlefschetz} adds the total spaces of Mumford families to the
known cases of \textup{(L)}, \Cref{prop:f3prime} records the known cases of
\textup{(F3$'$)}, and the routes of the literature use no other theorem of
this paper.
\end{remark}

\subsubsection*{Two of the three are not reductions}

One could hope that (F2) and (F3) are reductions in the ordinary sense,
statements to the effect that every Hodge class is carried, by some
correspondence, to a class on a Weil type variety or on an abelian variety,
so that the whole conjecture would follow from the Weil classes alone. Both
hopes are refutable, and we refute them, because a frontier that contains a
false statement is not a frontier.

\begin{proposition}[The Hodge ring is not generated by divisor and Weil
classes]\label{prop:mumfordproduct}
Let $X$ be a complex abelian fourfold whose Hodge group is a $\QQ$-form of an
almost direct product of three copies of $\mathrm{SL}_{2}$, acting on
$V=H^{1}(X,\QQ)$ through the tensor product $V_{1}\otimes V_{2}\otimes V_{3}$
of the three standard representations, with $\End^{0}(X)=\QQ$; such
fourfolds exist and form families of dimension one, by Mumford
\cite{Mum69}, and by \cite[2.5(i)]{MZ99} they are exactly the fourfolds
with $\End^{0}(X)=\QQ$ and Hodge group smaller than $\Sp_{8}$. Then:
\begin{enumerate}[label=\textup{(\roman*)},nosep,leftmargin=2.6em]
\item the Hodge ring of $X$ is generated by the polarisation:
  $\dim_{\QQ}\Hdg^{1}(X)=\dim_{\QQ}\Hdg^{2}(X)=1$;
\item $\dim_{\QQ}\Hdg^{1}(X\times X)=3$ and $\dim_{\QQ}\Hdg^{2}(X\times X)=8$,
  while the products of divisor classes span a subspace of dimension $6$
  in $H^{4}(X\times X,\QQ)$;
\item the two remaining classes lie in $H^{4}(X\times X,\QQ)$ and are Weil
  classes of no CM field acting on $X\times X$, on a quotient of it, or on an
  abelian subvariety of it.
\end{enumerate}
In particular the statement that every Hodge class on an abelian variety
lies in the subring generated by divisor classes and Weil classes of
quotients is false in dimension eight, though true through dimension five
\cite{MZ99}.
\end{proposition}

\begin{proof}
Hodge classes are the invariants of the Hodge group, so
$\Hdg^{k}(Y)\otimes\CC$ is the space of $\Hg(Y)_{\CC}$-invariants in
$\bigwedge^{2k}H^{1}(Y,\CC)$, and $\Hg(X\times X)=\Hg(X)$ acting
diagonally on $V\oplus V$. Over $\CC$ the group is
$\mathrm{SL}_{2}^{3}$ and the weights of $V$ under its maximal torus are the
eight sign vectors $(\pm1,\pm1,\pm1)$. Write $V=V_{1}\otimes W$ with
$W=V_{2}\otimes V_{3}$, on which $\mathrm{SL}_{2}^{2}=\mathrm{Spin}(4)$ acts
orthogonally. The Cauchy formula
$\bigwedge^{k}(V_{1}\otimes W)=\bigoplus_{|\lambda|=k}S_{\lambda}V_{1}\otimes
S_{\lambda'}W$ \cite[\S6.1]{FH91}, with $S^{2}V_{1}=S_{(3,1)}V_{1}=V(2)$,
$\bigwedge^{2}V_{1}=S_{(2,2)}V_{1}=\mathbf{1}$ and $S^{4}V_{1}=V(4)$, gives
$\bigwedge^{2}V=V(2)\otimes\bigwedge^{2}W\oplus S^{2}W$ and
$\bigwedge^{4}V=V(4)\otimes\bigwedge^{4}W\oplus V(2)\otimes S_{(2,1,1)}W
\oplus S_{(2,2)}W$. Under $\mathrm{Spin}(4)$, $\bigwedge^{2}W=V(2,0)\oplus
V(0,2)$, $S^{2}W=V(2,2)\oplus\mathbf{1}$, $\bigwedge^{4}W=\mathbf{1}$,
$S_{(2,1,1)}W\cong\mathfrak{sl}(W)=\bigwedge^{2}W\oplus S^{2}_{0}W$, and
$S_{(2,2)}W=V(4,0)\oplus V(0,4)\oplus V(2,2)\oplus\mathbf{1}$, the
decomposition of the algebraic curvature tensors in dimension four into the
self-dual and anti-self-dual Weyl parts, the traceless Ricci part and the
scalar part. Hence
\begin{gather*}
  \textstyle\bigwedge^{2}V=V(2,2,0)\oplus V(2,0,2)\oplus V(0,2,2)\oplus\mathbf{1},\\
  \textstyle\bigwedge^{4}V=V(2,2,2)\oplus\bigoplus_{\text{cyc}}V(4,0,0)
  \oplus\bigoplus_{\text{cyc}}V(2,2,0)\oplus\mathbf{1},
\end{gather*}
each with one trivial summand, which is (i). For (ii),
$\bigwedge^{4}(V\oplus V)=\bigoplus_{a}\bigwedge^{a}V\otimes\bigwedge^{4-a}V$
and the invariants in $\bigwedge^{a}V\otimes\bigwedge^{4-a}V$ are
$\Hom_{\mathrm{SL}_{2}^{3}}(\bigwedge^{a}V,\bigwedge^{4-a}V)$, since every
summand is self-dual; the counts are $1,1,4,1,1$, total $8$, the $4$ coming
from the four distinct irreducible summands of $\bigwedge^{2}V$. Likewise the
invariants in $\bigwedge^{2}(V\oplus V)$ number $1+1+1=3$. The divisor
classes of $X\times X$ are the three symplectic forms $\psi_{11}$, $\psi_{12}$,
$\psi_{22}$ on $V\oplus V$ built from the polarisation $\psi$ of $V$. Their
six pairwise products are linearly independent: they lie in the summands
$\bigwedge^{a}V\otimes\bigwedge^{4-a}V$ with $a=4,3,2,2,1,0$, each is
nonzero, and in the summand $a=2$ the classes $\psi_{11}\psi_{22}$ and
$\psi_{12}^{2}$ correspond, as maps $\bigwedge^{2}V\to\bigwedge^{2}V$, to a
multiple of the projection onto $\CC\psi$ and to a multiple of the
identity. The $\Sp_{8}$-invariants of $\bigwedge^{4}(V\oplus V)$ number
$1+1+2+1+1=6$ by the same count, since $\bigwedge^{1}V=V$,
$\bigwedge^{2}V=\bigwedge^{2}_{0}V\oplus\mathbf{1}$ and
$\bigwedge^{3}V=\bigwedge^{3}_{0}V\oplus V$ under $\Sp_{8}$; so the six
products are all the classes a Hodge group equal to $\Sp_{8}$ would allow. For (iii), the Weil classes of an abelian
variety $Y$ with respect to a CM field $F\subseteq\End^{0}(Y)$ are, by
definition, the classes in $\bigwedge^{2k}_{F}H^{1}(Y,\QQ)\subseteq
H^{2k}(Y,\QQ)$ with $2k=\dim_{F}H^{1}(Y,\QQ)$ (\Cref{setup:cmfield}). A CM
field acting on $X\times X$ lies in $\End^{0}(X\times X)=M_{2}(\QQ)$, so it is
imaginary quadratic, $2k=8$, and its Weil classes lie in $H^{8}$, while the two
classes of (ii) have degree four. Every proper
abelian subvariety and every proper quotient of $X\times X$ is isogenous to
$X$, because $X$ is simple, and $\End^{0}(X)=\QQ$ contains no CM field, so $X$
carries no Weil classes at all. This is the statement of \cite[0.1(iii)]{MZ99}
that the exceptional classes of $X\times X$ are not of Weil type. The last
sentence is (ii) and (iii) together with the theorem of \cite{MZ99} for
dimension at most five.
\end{proof}

\begin{proposition}[No Hodge-theoretic reduction to abelian varieties]
\label{prop:noabeliantype}
Let $Y$ be a very general smooth quintic threefold in $\PP^{4}$, and let
$V=H^{3}(Y,\QQ)$ with its polarised Hodge structure of weight three. Then $V$
is not isomorphic to a subquotient of any Hodge structure of the form
$\bigoplus_{i}T_{i}$, where each $T_{i}$ is obtained from the $H^{1}$ of
finitely many abelian varieties by tensor products, duals and Tate twists. In
particular the Hodge structure $H^{*}(Y,\QQ)$ is not of abelian type, and no
correspondence between $Y$ and an abelian variety, or a product of abelian
varieties, induces an injection of $V$ into the cohomology of the latter.
\end{proposition}

\begin{proof}
Write $\mathrm{MT}(W)$ for the Mumford--Tate group of a rational Hodge
structure $W$ and $\mu_{W}\colon\mathbb{G}_{m}\to\mathrm{MT}(W)_{\CC}$ for
its Hodge cocharacter, acting on $W^{p,q}$ by $z^{p}$. The adjoint action of
$\mu_{W}$ on $\operatorname{Lie}\mathrm{MT}(W)_{\CC}\subseteq\End(W_{\CC})$ has weights
among the differences $p-p'$ of two Hodge indices of $W$.

If $W$ has weight one, these differences lie in $\{-1,0,1\}$. If $W'$ is a
subquotient of a tensor construction $T$ on $W$, then $\mathrm{MT}(W')$ is the
image of $\mathrm{MT}(W)$ acting on $W'$, and $\mu_{W'}$ is the image of
$\mu_{W}$; so $\operatorname{Lie}\mathrm{MT}(W')$ is a quotient of $\operatorname{Lie}\mathrm{MT}(W)$
as a representation of $\mathbb{G}_{m}$ through $\mu$, and its weights are a
subset of $\{-1,0,1\}$. The same holds with $W=\bigoplus H^{1}(A_{i})$.
Hence for every Hodge structure of abelian type the adjoint weights of the
Hodge cocharacter lie in $\{-1,0,1\}$.

For $V=H^{3}(Y,\QQ)$ we show that the adjoint weight $3$ occurs. The
canonical bundle of a quintic threefold is trivial by adjunction, so
$h^{3,0}(Y)=1$; pick $0\ne e\in H^{3,0}$. The intersection form $\psi$ on $V$
is symplectic and pairs $H^{3,0}$ with $H^{0,3}$ perfectly, so the rank one
endomorphism $x\mapsto\psi(x,e)\,e$ of $V_{\CC}$ is nonzero, sends $H^{0,3}$
to $H^{3,0}$ and kills the rest, and lies in $\mathfrak{sp}(V_{\CC},\psi)$
because $\psi(\psi(x,e)e,y)+\psi(x,\psi(y,e)e)
=\psi(x,e)\bigl(\psi(e,y)+\psi(y,e)\bigr)=0$; its adjoint weight is
$3-0=3$. It remains to see that $\mathfrak{sp}(V_{\CC},\psi)\subseteq
\operatorname{Lie}\mathrm{MT}(V)_{\CC}$. Let $S$ be the parameter space of smooth quintic
threefolds, connected and nonsingular, and $\mathbb{V}$ the polarisable
variation of Hodge structure $R^{3}$ of the universal family. The connected
algebraic monodromy group $G^{0}$ of $\mathbb{V}$ contains the connected
monodromy group of any Lefschetz pencil of quintics, that is of hyperplane
sections of the fifth Veronese image of $\PP^{4}$; for a Lefschetz pencil of
odd fibre dimension the image of the monodromy representation on the
$\ell$-adic vanishing cohomology $E\otimes\QQ_{\ell}$ is open in
$\Sp(E\otimes\QQ_{\ell},\psi)$ \cite[Theorem 5.10]{Del74}, so the Zariski
closure of the monodromy group over $\QQ$ is $\Sp(E,\psi)$; here
$E=H^{3}(Y_{t},\QQ)$, the vanishing cohomology being the orthogonal
complement of the image of $H^{3}(\PP^{4},\QQ)=0$. Hence $G^{0}\supseteq
\Sp(V,\psi)$, and $G^{0}\subseteq\Sp(V,\psi)$ since the monodromy preserves
$\psi$. At a very general
point $Y$ the group $G^{0}$ is a normal subgroup of the derived group of
$\mathrm{MT}(V)$ \cite[Theorem 1]{And92}, so $\Sp(V,\psi)\subseteq
\mathrm{MT}(V)^{\mathrm{der}}\subseteq\Sp(V,\psi)$, and the adjoint weight
$3$ occurs in $\operatorname{Lie}\mathrm{MT}(V)$. So $V$ is not of abelian type, and
neither is any Hodge structure containing it as a subquotient, in particular
$H^{*}(Y,\QQ)$. A correspondence inducing an injection of $V$ into the
cohomology of an abelian variety would exhibit $V$ as a sub-Hodge structure
of a Tate twist of that cohomology, which is of abelian type.
\end{proof}

\begin{figure}[tbp]
\centering
  \includegraphics[width=\textwidth]{fig_closure}
\caption{The route of this paper through the rule set of \Cref{thm:closure},
drawn as a stack of plates, one per level of the derivation; the other routes
are drawn in \Cref{fig:frontier}. Heavy arrows are the rules used by the
minimal set $\{\mathrm{P2}_{\mathrm{split}},\mathrm{F2},\mathrm{F3}'\}$,
light arrows the rules it does not use. The bottom plate carries the
propagation statement, read in the polarised form of \Cref{rem:p2prime},
once for each of the four kinds of Weil family, with the secant route beside
it. The small grass blocks on the next plate are the base points of
\Cref{thm:everyfamily} and \Cref{thm:cmbasepoint}, one in every family. The
three arrows inside that plate are \Cref{prop:descent}: descent from the
trivial discriminant to the others, and scalar extension from the fields of
degree at least four to the quadratic fields. With them (P2) is needed only
for the split families of the fields of higher degree, and the other three
(P2) blocks are not used. (F2) is the conjecture for the Hodge classes that
no Weil line reaches, and (F3) is the whole conjecture
(\Cref{prop:f3ishc}), so the route is minimal only with (F3) replaced by
\textup{(F3$'$)} of \Cref{def:f3prime}. The bracket under the two quadratic
(P2) blocks records the first form of the criterion, in which the Chern
character is exactly a Weil class, as \Cref{thm:p2numerical},
\Cref{thm:p2support} and \Cref{cor:hhfactor} describe it; by
\Cref{thm:p2false} that form is met by no complex. The secant route reaches
the trivial discriminant, since the member it is built on is split, and
through descent the rest of the quadratic case; it has no edge to the fields
of higher degree. The bounded criterion (P1) stands apart, closed on itself:
it is equivalent to its own conclusion.}
\label{fig:closure}
\end{figure}


\subsection{Routes from outside algebraic geometry}\label{ssec:outside}

By \Cref{prop:f2isab} and \Cref{prop:f3prime}(i), the two statements
\textup{(F2)} and \textup{(F3$'$)} are the two halves of the conjecture: the
conjecture for abelian varieties, and the conjecture modulo abelian varieties.
A route from another domain must therefore either produce algebraic cycles or
arrive at one of the hypotheses of \Cref{setup:rules}. This subsection
records what the routes most often proposed do. One of them, the analytic
route through positive currents, can be made exact: the conjecture is the
vanishing of an integrality gap in a mass minimisation problem whose real
version is solved, and positivity, which is what analysis sees first, cannot
distinguish an algebraic class from any other.

Throughout, $X$ is a smooth complex projective variety of dimension $n$, $h\in
H^{2}(X,\ZZ)$ is the class of an ample line bundle, $\omega_{h}$ is a K\"ahler
form with class $h$, and for $c\in H^{2p}(X,\QQ)$ we put
\[
  L_{h}(c)=\int_{X}c\cup\frac{h^{\,n-p}}{(n-p)!}\in\QQ .
\]
A real $(p,p)$-form is \emph{strongly positive} if at every point it is a
nonnegative combination of forms $\mathrm{i}\beta_{1}\wedge\bbar{\beta}_{1}
\wedge\cdots\wedge\mathrm{i}\beta_{p}\wedge\bbar{\beta}_{p}$ with $\beta_{j}$
of type $(1,0)$; the closed positive currents of bidegree $(p,p)$ are the
currents nonnegative on such forms of the complementary degree, and every
smooth closed strongly positive form is one \cite[Chapter~III]{Dem12}.

\begin{proposition}[Positivity is no test]\label{prop:positivitynotest}
For every class $\alpha\in H^{p,p}(X,\RR)$, $0\le p\le n$, there is an integer
$N$ such that $\alpha+Nh^{p}$ is the class of a smooth closed strongly positive
$(p,p)$-form. Consequently every real $(p,p)$-class becomes the class of a
closed positive current after the addition of a multiple of $h^{p}$, and, since
$h^{p}$ is algebraic, a Hodge class is algebraic if and only if the positive
class $\alpha+Nh^{p}$ is: no property of a class implied by positivity, or by
its being a limit of classes of effective cycles, separates the algebraic Hodge
classes from the others.
\end{proposition}

\begin{proof}
Let $V=T_{x}X$ with the hermitian form $\omega_{h}(x)$, and let
$\mathrm{SP}_{p}\subseteq\Lambda^{p,p}_{\RR}V^{\vee}$ be the closed convex cone
generated by the forms $\mathrm{i}^{p}\beta_{1}\wedge\bbar{\beta}_{1}\wedge\cdots
\wedge\beta_{p}\wedge\bbar{\beta}_{p}$. These forms span
$\Lambda^{p,p}_{\RR}V^{\vee}$, so $\mathrm{SP}_{p}$ has nonempty interior, and
$\omega_{h}(x)^{p}$ lies in that interior: in a unitary basis $dz_{1},\dots,
dz_{n}$ of $V^{\vee}$,
\[
  \omega_{h}^{p}=p!\sum_{|J|=p}\Bigl(\frac{\mathrm{i}}{2}\Bigr)^{p}
  dz_{J}\wedge d\bbar{z}_{J},
\]
and a generator $D=\mathrm{i}^{p}\beta_{1}\wedge\bbar{\beta}_{1}\wedge\cdots\wedge
\beta_{p}\wedge\bbar{\beta}_{p}$ equals $|\det B|^{2}\,2^{p}(\mathrm{i}/2)^{p}
dz_{1}\wedge d\bbar{z}_{1}\wedge\cdots\wedge dz_{p}\wedge d\bbar{z}_{p}$ after a
unitary change of coordinates taking the span of the $\beta_{j}$ to that of
$dz_{1},\dots,dz_{p}$, with $B$ the matrix of the $\beta_{j}$ in that basis;
so $\omega_{h}^{p}-\delta D\in\mathrm{SP}_{p}$ for $0<\delta\le p!/(2^{p}
|\det B|^{2})$. As every real $(p,p)$-form is a finite real combination of
generators, $\omega_{h}^{p}+\varepsilon\gamma\in\mathrm{SP}_{p}$ for every
$\gamma$ and all small $\varepsilon$, which is interiority. The cone, the form
$\omega_{h}(x)^{p}$ and the norm defined by $\omega_{h}(x)$ are carried into
each other by unitary frames, so the radius $r>0$ of a ball around
$\omega_{h}(x)^{p}$ contained in $\mathrm{SP}_{p}$ does not depend on $x$.
Represent $\alpha$ by a smooth closed real $(p,p)$-form $a$, for instance its
harmonic representative, and let $N\ge\max_{X}|a|/r$. Then
$a+N\omega_{h}^{p}=N(\omega_{h}^{p}+a/N)$ is strongly positive at every point,
closed, and represents $\alpha+Nh^{p}$. The remaining assertions follow, the
last because a class of the form $\alpha+Nh^{p}$ is positive for every $\alpha$.
\end{proof}

The proposition says what positivity cannot do. Demailly studied the cone of
closed positive currents against the conjecture and constructed extremal
closed strongly positive currents that are not integration currents on
analytic sets \cite{Dem82}, so extremality does not single out the algebraic
cycles either. What positivity can do is turn the conjecture into a
variational statement, and the statement is exact. It is the K\"ahler case of
a formulation that Harvey and Lawson give for every calibrated manifold
\cite[Section~6, Proposition~6.13 and Remarks~6.9 and~6.11]{HL09}; we state
it for the classes of this paper, with the effective representative that
\Cref{lem:cibig} supplies, and then apply it to the Mumford square. For a
closed integral current $T$ of dimension $2(n-p)$ on $X$, Wirtinger's
inequality \cite[1.8.2]{Fed69}, \cite{HK74}, \cite[II.3]{HL82} gives
\begin{equation}\label{eq:wirtinger}
  \mathbf{M}(T)\;\ge\;T\Bigl(\frac{\omega_{h}^{\,n-p}}{(n-p)!}\Bigr)
  \;=\;L_{h}(c),
\end{equation}
where $\mathbf{M}$ is the mass and $c\in H^{2p}(X,\ZZ)$ is the Poincar\'e dual
of the homology class of $T$, with equality if and only if the approximate
tangent plane of $T$ is a complex subspace, positively oriented, at
$\|T\|$-almost every point. The right side depends on $c$ alone.

\begin{lemma}[Complete intersection classes are big]\label{lem:cibig}
For every algebraic class $\beta\in H^{2p}(X,\QQ)$ there is an integer $M$ such
that $\beta+Mh^{p}$ is a positive rational combination of classes of
irreducible subvarieties of codimension $p$; so some positive integer multiple
of $\beta+Mh^{p}$ is the class of an effective algebraic cycle.
\end{lemma}

\begin{proof}
Replacing $h$ by a multiple, which changes $M$, we may assume $h$ very ample.
Let $Z\subseteq X$ be irreducible of codimension $p$, and let $d$ be such that
the ideal sheaf $I_{Z}(d)$ is globally generated. The base locus of the linear
system of divisors of degree $d$ containing $Z$ is exactly $Z$, so by Bertini
$p$ general members $D_{1},\dots,D_{p}$ of it meet properly away from $Z$.
Every component of $D_{1}\cap\cdots\cap D_{p}$ has codimension at most $p$;
a component not contained in $Z$ has codimension exactly $p$, because the
$D_{i}$ meet properly off $Z$, and a component contained in $Z$ has
codimension at least $p$ and so equals $Z$. The intersection is therefore the
union of $Z$ and a closed set $W'$ of pure codimension $p$. Hence $d^{p}h^{p}=[D_{1}]\cdots[D_{p}]=m[Z]+[W']$ with $m\ge1$ and
$W'$ effective \cite[Chapter~7]{Ful98}, and $h^{p}-\varepsilon[Z]$ is an
effective rational class for $0<\varepsilon\le m/d^{p}$. Write
$\beta=\sum_{i}n_{i}[Z_{i}]$ with $n_{i}\in\QQ$ and choose $\varepsilon_{i}$ for
each $Z_{i}$ with $n_{i}<0$; then for $M\ge\sum_{n_{i}<0}|n_{i}|/\varepsilon_{i}$
the class $\beta+Mh^{p}$ is $\sum_{n_{i}>0}n_{i}[Z_{i}]+\sum_{n_{i}<0}
(|n_{i}|/\varepsilon_{i})(h^{p}-\varepsilon_{i}[Z_{i}])+(M-\sum_{n_{i}<0}
|n_{i}|/\varepsilon_{i})h^{p}$, a positive rational combination of effective
classes, $h^{p}$ itself being the class of a complete intersection.
\end{proof}

\begin{theorem}[The conjecture as a mass gap]\label{thm:massgap}
Let $\alpha\in\Hdg^{p}(X)$ be a rational Hodge class.
\begin{enumerate}[label=\textup{(\roman*)},leftmargin=2.6em,itemsep=2pt]
\item Let $N\ge1$ be an integer with $N\alpha$ integral, let $M$ be an
  integer for which $\alpha+Mh^{p}$ is the class of a smooth closed strongly
  positive form, as every large $M$ is by \Cref{prop:positivitynotest}, and
  put $c=N(\alpha+Mh^{p})$. The minimum of the mass over closed real normal
  currents of dimension $2(n-p)$ in the class dual to $c$ is $L_{h}(c)$, and
  it is attained by a smooth strongly positive form. The minimum over closed
  integral currents in the same class exists and is at least $L_{h}(c)$.
\item The class $\alpha$ is algebraic if and only if there are integers
  $N\ge1$ and $M\ge0$, with $N\alpha$ integral, for which the minimum of the
  mass over closed integral currents in the class dual to $N(\alpha+Mh^{p})$
  equals $L_{h}(N(\alpha+Mh^{p}))$; and then every mass minimising integral
  current in that class is the integration current of an effective algebraic
  cycle.
\item Write $\|c\|_{\ZZ}$ for the minimum in \textup{(i)} over integral
  currents. By a theorem of Federer, recorded in this form by Harvey and
  Lawson \cite[Section~6]{HL09}, \cite{Fed69,Fed74}, $\lim_{k\to\infty}
  \|kc\|_{\ZZ}/k=L_{h}(c)$: the integrality gap closes asymptotically for
  every class. The class $\alpha$ is algebraic if and only if, for every
  sufficiently large $M$, the gap is closed at some finite $k$, that is,
  $\|kc\|_{\ZZ}=kL_{h}(c)$ for some $k\ge1$.
\end{enumerate}
In particular the conjecture for $X$ is the statement that for every rational
Hodge class the integrality gap of this problem vanishes for some $N$ and $M$;
\textup{(F2)} is that statement for abelian varieties. Teh and Yang give a
different sufficient condition in the same setting, through Lipschitz chains
that are closed for $d$ and $d^{c}$ with positive $(k,k)$-part
\cite[Theorem~4.2]{TY21}.
\end{theorem}

\begin{proof}
(i) The form $\omega_{h}^{n-p}/(n-p)!$ has comass one and calibrates exactly
the complex $2(n-p)$-planes, which is Wirtinger's inequality
\cite[1.8.2]{Fed69}; so every closed real normal current $S$ in the class has
$\mathbf{M}(S)\ge S(\omega_{h}^{n-p}/(n-p)!)=L_{h}(c)$. By
\Cref{prop:positivitynotest} the class is represented by a smooth closed
strongly positive form $a$, and for a strongly positive $(p,p)$-form the mass of
the current it defines is $\int_{X}a\wedge\omega_{h}^{n-p}/(n-p)!$, because a
strongly positive $2(n-p)$-vector is a nonnegative combination of complex unit
planes, each of mass one and calibrated. So the minimum over real currents is
$L_{h}(c)$ and $a$ attains it. Closed normal currents represent every real
homology class and closed integral currents every integral one
\cite[5.11]{FF60}, and by the compactness theorem of Federer and Fleming the
mass attains its minimum on the closed integral currents in a given class
\cite[9.6]{FF60}; the lower bound is \eqref{eq:wirtinger}.

(ii) If $\alpha$ is algebraic, \Cref{lem:cibig} gives $M$ and $N$ with
$N(\alpha+Mh^{p})=\sum_{j}n_{j}[V_{j}]$, $n_{j}\in\ZZ_{>0}$, $V_{j}$
irreducible of codimension $p$. The integration current $T=\sum_{j}n_{j}[V_{j}]$
is a closed locally integral current \cite[Theorem~3.1.1]{Kin71}, its tangent
planes are complex and positively oriented, so equality holds in
\eqref{eq:wirtinger} and $\mathbf{M}(T)=L_{h}(c)$ with $c=N(\alpha+Mh^{p})$.
Conversely let $T$ be a closed integral current in the class dual to $c$ with
$\mathbf{M}(T)=L_{h}(c)$, for instance a mass minimiser when the minimum is
$L_{h}(c)$. Equality in \eqref{eq:wirtinger} makes $T$ a closed locally
rectifiable current of bidimension $(n-p,n-p)$ which is positive, and by
King's theorem \cite[Theorem~5.2.1]{Kin71} such a current is a holomorphic
cycle $\sum_{j}n_{j}[V_{j}]$ with $n_{j}$ nonzero integers and $V_{j}$
irreducible analytic subvarieties of dimension $n-p$; positivity makes every
$n_{j}$ positive. (Harvey and Shiffman characterised the holomorphic chains
among closed rectifiable currents of bidimension $(n-p,n-p)$ without
positivity, under a hypothesis on the support that Alexander removed
\cite{HS74,Ale97}.) By Chow's theorem \cite[p.~167]{GH78} the $V_{j}$
are algebraic, so $c$ is the class of an effective algebraic cycle and $\alpha$
is algebraic. This is the argument of \cite[Proposition~6.13]{HL09}.

(iii) The limit is Federer's theorem on the real and integral masses of a
class, as stated in the final note of \cite[Section~6]{HL09}; with (i) the
limit is $L_{h}(c)$, since the real minimum is attained by a smooth form. If
the gap is closed at $k$, then (ii), applied with $kN$ in place of $N$, makes
$\alpha$ algebraic. Conversely, if $\alpha$ is algebraic, \Cref{lem:cibig}
gives $M_{0}$ such that $\alpha+M_{0}h^{p}$ is a positive rational
combination of classes of subvarieties; for $M\ge M_{0}$ so is
$\alpha+Mh^{p}$, since $h^{p}$ is the class of a complete intersection, and
some multiple $kc$ is the class of an effective cycle, whose integration
current attains the bound. For a fixed small $M$ the equivalence is not asserted, since it would
need an effective multiple of every algebraic class represented by a
strongly positive form, in codimension at least two. The last sentence of the
theorem is the definition of the conjecture for $X$, and \Cref{prop:f2isab}
for \textup{(F2)}.
\end{proof}

\begin{example}[The Mumford square]\label{ex:mumfordmass}
Let $Y=X\times X$ be the square of a Mumford fourfold at a very general point
of its curve, with the product principal polarisation $\theta_{Y}$, and let
$\omega$ be one of the two exceptional classes of \Cref{prop:mumfordproduct},
normalised in $H^{2}(X)\otimes H^{2}(X)$ with $a_{0}=0$ in the notation of
\Cref{setup:omegaa}, so that it corresponds to an endomorphism of $H^{2}(X)$
killing $\theta$ with image in the primitive part. Then
$L_{\theta_{Y}}(\omega)=0$: the pairing of $\omega$ with
$\theta_{Y}^{6}=\sum_{k}\binom6k\theta^{k}\otimes\theta^{6-k}$ sees only the
term $20\,\theta^{3}\otimes\theta^{3}$, and $\int_{X}x\cup\theta^{3}=0$ for
every primitive $x\in H^{2}(X)$. Also $L_{\theta_{Y}}(\theta_{Y}^{2})=8!/6!=56$,
and $4\theta_{Y}^{2}$ is the class of the complete intersection $D\cap D'$ of two
general members of $|2\theta_{Y}|$, a base point free system, whose integration
current has mass $224$. By \Cref{thm:massgap} and \Cref{cor:mumfordone}, the
Hodge conjecture for every member of the Mumford family, and with it
$B(W\times_{C}W)$, is equivalent to the following: for some integers
$N\ge1$ and $M\ge1$ with $N\omega$ integral, the class dual to
$N(M\theta_{Y}^{2}+\omega)$ contains an integral current of mass exactly
$56NM$; and any such current is an effective algebraic cycle representing the
class. For $M$ divisible by four the class $NM\theta_{Y}^{2}$ contains one,
$NM/4$ copies of $D\cap D'$, of mass $56NM$; so the question is whether the
perturbation by $N\omega$, which changes the Wirtinger bound by nothing, can
be absorbed without mass. By \Cref{thm:mumfordcm} the answer is yes at
every CM point of the curve, with integers $N$ and $M$ depending on the
point; and if the same $N$ and $M$ served at infinitely many CM points the
answer would be yes everywhere, by \Cref{cor:mumfordbounded}, because the
components of an effective cycle of degree $56NM$ have degree at most $56NM$
and the Hilbert polynomials of subvarieties of bounded degree form a finite
set.
\end{example}

\Cref{fig:massgap} draws the theorem and the example.

\begin{figure}[tbp]
\centering
  \includegraphics[width=\textwidth]{fig_massgap}
\caption{\Cref{thm:massgap} and \Cref{ex:mumfordmass}. Left: the mass of a
current in the class dual to $kc$ against $k$. The line $k\,L_{h}(c)$ is the
minimum over closed real normal currents, attained by a smooth strongly
positive form; the minima over closed integral currents lie on or above it,
with $\|kc\|_{\ZZ}/k\to L_{h}(c)$ by Federer's theorem, and for $M$ large the
class is algebraic exactly when one of them lies on the line, where the minimiser is
an effective algebraic cycle. The heights of the dots are schematic. Right:
the numbers of the Mumford square, in the split model. Of the seven terms $\binom6k\theta^{k}\otimes\theta^{6-k}$ of
$\theta_{Y}^{6}$ only $k=3$ pairs with $\pi_{0}=\frac14\theta\otimes\theta$,
giving $20\cdot144=2880$, and none pairs with the primitive summands; so
$L_{\theta_{Y}}(\omega_{a})=4a_{0}$, the Wirtinger bound of an exceptional
class with $a_{0}=0$ is zero, and the perturbation of $NM\theta_{Y}^{2}$ by
$N\omega$ changes the bound $56NM$ by nothing.}
\label{fig:massgap}
\end{figure}

\subsection{The routes at a glance}

\begingroup
\renewcommand{\arraystretch}{1.25}
\begin{longtable}{@{}>{\raggedright\arraybackslash}p{0.30\linewidth}>{\raggedright\arraybackslash}p{0.30\linewidth}>{\raggedright\arraybackslash}p{0.30\linewidth}@{}}
\caption{The routes considered, and the boundary of each result.}
\label{tab:scope}\\
\toprule
\textbf{Route} & \textbf{Result} & \textbf{Beyond the result} \\
\midrule
\endfirsthead
\multicolumn{3}{@{}l}{\footnotesize\itshape Table \thetable, continued.}\\[2pt]
\toprule
\textbf{Route} & \textbf{Result} & \textbf{Beyond the result} \\
\midrule
\endhead
\midrule
\multicolumn{3}{r@{}}{\footnotesize\itshape continued on the next page}\\
\endfoot
\bottomrule
\endlastfoot
Ambient object, restricted to $\Av$, transported by $\Phi_{\cP}$
  & closed by \Cref{thm:divisor}
  & ambient varieties with exceptional classes \\
Any $K$-equivariant construction with the norm character
  & closed by \Cref{thm:character}
  & constructions that are not eigenclasses \\
Secant dimension count
  & no obstruction, \Cref{cor:nonumerical}
  & none \\
Semiregularity dimension count
  & no obstruction, \Cref{cor:countnotbind}
  & none \\
Semiregularity of an object whose Chern character sits on the Weil line
  & closed by \Cref{thm:rankzero} and \Cref{cor:kernelfamily} for direct
    sums, and for every object by \Cref{thm:p2false}, for every CM field by
    \Cref{thm:cmpurefalse}
  & objects of positive rank; the weaker hypothesis of
    \cite[Question 11.4]{Mar25b} \\
The polarised criterion at a quartic CM field, $n=2$
  & computed for every shape, \Cref{thm:quarticrank}: a complex meeting it
    has $\dim\Ext^{2}(E,E)\ge88$, \Cref{cor:quarticleast}
  & existence of such a complex \\
Induction through hyperplane sections
  & closed by \Cref{thm:hyperplane}, \Cref{thm:returntrip}
  & correspondences, degenerations, coniveau \\
Kuga--Satake, from a weight two structure of K3 type
  & not obstructed, \Cref{prop:ksnotrestriction}
  & available only at $n=1$, \Cref{cor:ksboundary} \\
Divisors, on a member with quaternionic multiplication
  & succeeds, \Cref{thm:quatweil}
  & fails at full Hodge group, \Cref{prop:divfails} \\
Propagation off the base locus
  & reduced to one bound or one rank, \Cref{thm:degreebound},
    \Cref{thm:semiregclosure}
  & the bound or the rank, \Cref{rem:tworemain} \\
Objects with effective Chern character at a base point
  & closed by \Cref{thm:positivity}
  & objects indecomposable in $D^{b}$, \Cref{rem:whatitmustbe} \\
The weakened criterion for a secant object of positive rank
  & reduced to two conditions on the factor, \Cref{thm:factor}, the first
    automatic for Markman's construction, \Cref{cor:markmancondition}; the
    identity fails for Markman's candidate, \Cref{thm:candidatefails}, and at every
    transversal double point, \Cref{cor:localobstruction}
  & objects whose support has no such point, \Cref{rem:candidatemeaning} \\
A secant object $\cI_{S}(b\Theta)$ with smooth support
  & invariants forced and $d\le b^{2}$, so $d\le9$ at $b=3$,
    \Cref{thm:smoothinvariants}, \Cref{cor:fourdiscriminants}
  & existence for $d\in\{1,3,5,7\}$ at $b=3$, one vanishing in
    $H^{1}(S,N_{S/X})$ and first-order extension, \Cref{rem:smoothremains} \\
The local part of the weakened criterion
  & it is a condition on depth: no obstruction at a Cohen--Macaulay point,
    \Cref{thm:cmvanishing}, \Cref{cor:cmdichotomy}
  & Cohen--Macaulay supports, where the identity becomes one global
    vanishing beside first-order extension \\
A support that is a complete intersection
  & closed for every $b$ and every $d$, \Cref{thm:noci}
  & supports of higher Hilbert--Burch rank, \Cref{prop:lowrankburch} \\
A support with a Hilbert--Burch resolution by line bundles or by simple
semi-homogeneous bundles
  & closed for every rank, twist and $d$, \Cref{thm:nosplit},
    \Cref{thm:noblocks}
  & resolutions over bundles that are not projectively flat \\
The secant route, granted every demand it makes
  & reaches $\delta=\delta_{0}$, and every $\delta$ by \Cref{prop:descent},
    \Cref{thm:closure}(iv)
  & the CM fields of degree at least four, which give the quadratic ones by
    \Cref{prop:descent}(ii) \\
Propagation from a base point
  & closes every imaginary quadratic family at once, \Cref{thm:degreebound},
    \Cref{thm:semiregclosure}
  & (P2) for the family, \Cref{cor:fourremain} \\
The bounded form (P1), read on the transported cycle
  & closed: the transport spreads each component by $c^{2n}/m_{i}$ and the
    data is unbounded whenever a component has finite stabiliser,
    \Cref{lem:multiplicity}, \Cref{thm:p1forces}
  & nothing; the other reading is closed too \\
The bounded form (P1), read on the minimal representative
  & closed: equivalent to the conjecture for the family,
    \Cref{thm:p1equivalent}
  & nothing; it is a restatement, so it is not a route \\
The pure form of (P2), for a given representative
  & forces the evaluation map to vanish on the $4n^{2}$-dimensional
    annihilator, and excludes every representative that is a sheaf on a
    smooth germ of codimension two of its support, \Cref{thm:p2forces},
    \Cref{cor:p2shape}
  & nothing: no representative exists, \Cref{thm:p2false} \\
The same form, tested by a number
  & $\dim\Ext^{2}(E,E)=2n(2n-1)$ at one base point would suffice,
    \Cref{thm:p2numerical}, and no representative attains it,
    \Cref{thm:p2false}
  & nothing \\
The same form, for a representative in codimension $n$
  & closed at every member of every family: the support must have a
    component of codimension less than $n$, and at a general member every
    component of codimension at least two has multiplicity zero,
    \Cref{lem:finiteproduct}, \Cref{thm:p2support}
  & nothing: no representative exists, \Cref{thm:p2false} \\
The Hochschild action on such a representative
  & the annihilator is the mixed part of a splitting of $HH^{1}$ and the
    action factors through a ring of dimension $2^{2n+1}-1$ in one of
    dimension $2^{4n}$, \Cref{thm:hhsplitting}, \Cref{cor:hhfactor}
  & nothing: no such complex exists, \Cref{thm:p2false} \\
The shape of such a representative
  & an indecomposable summand carries the class; at $n=2$, and at $n=3$
    granting the compatibility of \Cref{cor:hhfactor}(iii), its
    endomorphism algebra has dimension at least three, so no simple object
    serves, \Cref{thm:p2endo}
  & nothing: no object of any shape meets the pure form, \Cref{thm:p2false} \\
The pure form of (P2), Chern character exactly a Weil class
  & false for every imaginary quadratic field and every $n\ge2$: a
    semiregular such complex would deform to very general non-algebraic Weil
    tori, whose coherent sheaves have no Chern character below the top
    degree, \Cref{prop:weiltorussheaves}, \Cref{thm:p2false}
  & the polarised form (P2$'$), whose Chern character may carry powers of
    the polarisation, \Cref{rem:p2prime} \\
The reduction to Weil classes, and to abelian varieties
  & false as a statement about Hodge rings in dimension eight,
    \Cref{prop:mumfordproduct}; impossible Hodge-theoretically,
    \Cref{prop:noabeliantype}
  & the algebraicity of the classes beyond the Weil lines, which is (F2), and
    (F3), which is the whole conjecture, \Cref{prop:f3ishc} \\
The two classes on the self-product of a Mumford fourfold
  & a real multiplication by a cubic field on the primitive part of $H^{2}$,
    algebraic once the Kuga--Satake class of one K3 surface of Picard number
    thirteen is, \Cref{prop:mumfordrm}, \Cref{thm:mumfordks}
  & that Kuga--Satake class, \Cref{rem:mumfordks} \\
Other routes to the same two classes
  & degeneration and larger families closed off by rigidity; semiregularity
    reduced to $\dim\Ext^{2}=119$ at one CM point; specialisation reduced to
    a degree bound, \Cref{thm:mumfordrigid}, \Cref{thm:mumfordnumerical},
    \Cref{cor:mumfordbounded}
  & an object with $\dim\Ext^{2}=119$, or the degree bound,
    \Cref{rem:bypassaudit} \\
The cycles at a CM point, for the same two classes
  & the fourfold is of Weil type for $k_{c}\subset\End^{0}(X_{c})$ with the
    two exotic classes as its Weil line; pull-backs of that line along
    homomorphisms with components in $k_{c}$ span, with divisor products,
    $110$ of the $132$ Hodge classes of the square and no exceptional class,
    and an exceptional class is reached exactly by components generating the
    endomorphism algebra over $k_{c}$, \Cref{prop:cmsource}
  & a degree bound for cycles using the real multiplication, at infinitely
    many CM points, \Cref{rem:cmsource} \\
Objects built from line bundles, for the same two classes
  & closed at every point of the curve: the Chern characters stay invariant
    under the Lefschetz group, \Cref{thm:lefschetzclosure}
  & steps that break that symmetry, \Cref{rem:leavel} \\
Twistor deformation of the square, for the same two classes
  & closed: the component of the Hodge locus among all complex tori is the
    Mumford curve, which contains no twistor line, \Cref{prop:notwistor}
  & the Kuga--Satake class, which moves in a seven-dimensional family \\
Five further routes, for the same two classes
  & one statement: three of them are equivalent to the target and two imply
    it, so any one of them suffices, \Cref{cor:mumfordone}
  & the target itself, which is only the first case of (F2),
    \Cref{rem:mumfordneeded} \\
Pull-back from a hyperk\"ahler manifold, for the same two classes
  & the square is holomorphic symplectic and the classes are the twisted dual
    forms of the part of $H^{2}$ its symplectic class generates; a symplectic
    rational map to some $S_{\mu}^{[n]}$, $n\ge5$, suffices, and no K3
    surface or hyperk\"ahler eightfold serves, \Cref{prop:hksquare}
  & such a map, \Cref{rem:hkroute} \\
Reduction modulo $p$, for the same two classes
  & represented by cycles at every closed point of every reduction; the
    target is equivalent to a bound on their Hilbert data on a Zariski dense
    set, \Cref{prop:modp}, \Cref{thm:modpdensity}
  & the bound, \Cref{rem:modpremains} \\
The Weil classes of every CM field
  & the same propagation problem, with a CM base point in every family,
    \Cref{thm:cmbasepoint}, \Cref{thm:cmpropagation}
  & (P2) for the split families \\
Routes from outside algebraic geometry
  & positivity of currents sees nothing; the conjecture is the vanishing of
    the integrality gap of a mass minimisation whose real version is solved,
    \Cref{prop:positivitynotest}, \Cref{thm:massgap}
  & a calibrated integral current in one class, \Cref{ex:mumfordmass} \\
The closure, as a whole
  & five minimal sets of hypotheses: $\{\mathrm{F3}\}$ and
    $\{\mathrm{L},\mathrm{M}\}$, each equivalent to the conjecture, and
    three through the conjecture modulo abelian varieties,
    $\{\mathrm{L},\mathrm{F3}'\}$, $\{\mathrm{V},\mathrm{F3}'\}$ and this
    route, $\{\mathrm{P2}_{\mathrm{split}},\mathrm{F2},\mathrm{F3}'\}$,
    \Cref{thm:closure}, \Cref{prop:f3ishc}, \Cref{cor:fourremain},
    \Cref{rem:bullseye}
  & the Lefschetz standard conjecture, (M), (V), (F2), (F3$'$) \\
Convolutions of the Weil pieces on $E_{0}^{2n}$
  & at $n=3$ every convolution has $\dim\Ext^{2}\ge69>57$,
    \Cref{thm:noconvolution}; at $n=4$, with the pieces at two shifts,
    $\dim\Ext^{2}\ge40960>104$, \Cref{thm:efourtwolevels}, at three
    consecutive shifts $\dim\Ext^{2}\ge22528$, \Cref{thm:efourthreelevels},
    with the shifts within six consecutive values $\dim\Ext^{2}\ge640$,
    \Cref{cor:efoursix}, and with each parity at one shift, at any odd
    distance, $\dim\Ext^{2}\ge1792$, \Cref{prop:efoursingle}
  & the other arrangements of the shifts at $n=4$,
    \Cref{rem:convolutionsopen} \\
Diagonal complete intersections of Vandermonde type
  & quotients of powers of a generalised Fermat curve, so in $\mathcal{A}$;
    the conjecture holds when it holds on the abelian varieties $B_{[a]}$, and
    holds for two diagonal hypersurfaces of degree $3$, $4$, $6$ in
    $\PP^{6}$, \Cref{thm:vandermonde}, \Cref{cor:twodiagonal}; it holds on
    the very general member for $d=2$ in every dimension, for $d=3,4$ up
    to dimension seven and for $d=6$ up to dimension five,
    \Cref{thm:vgvandermonde}
  & $B_{[a]}$ of dimension above five, the Weil classes of
    $\QQ(\zeta_{e})$ with $\varphi(e)\ge4$, and the very general member
    for $d=3,4,6$ from dimension eight on and for $d=6$ in dimension six,
    \Cref{rem:vandermonde} \\
The conjecture modulo abelian varieties, (F3$'$)
  & holds on the class $\mathcal{A}$ of varieties whose cohomology is
    reached from abelian varieties by algebraic correspondences (curves,
    abelian varieties, products, surjective images), in degrees $0$, $2$,
    $2n-2$, $2n$, and in dimension at most three, \Cref{prop:f3prime};
    on the Delsarte hypersurfaces and the diagonal complete intersections
    of Vandermonde type, \Cref{cor:delsarteall}, \Cref{thm:vandermonde};
    on the fourfolds and fivefolds of \Cref{prop:f3primesmall},
    \Cref{prop:f3primefibration} and \Cref{prop:f3primedominant};
    through the Hodge conjecture on the varieties of \Cref{prop:k3powers}
    and \Cref{prop:k3ntype}
  & degree four on fourfolds outside $\mathcal{A}$ not covered by these
    results, and beyond, for
    instance squares of K3 surfaces with real multiplication, such as
    $S_{\lambda}$, \Cref{rem:f3primeopen},
    \Cref{rem:f3primefrontier} \\
\end{longtable}
\endgroup
